% Génétique et hérédité humaines : éradication des préjugés autour des anomalies — SVT 1ère A — Cours signé OnBuch+
% Programme officiel MINESEC. Compiler avec : tectonic NA02-eradication-prejuges-anomalies-caracteres.tex
\newcommand{\DOCMATIERE}{SVT}
\newcommand{\DOCNIVEAU}{1ère A}
\newcommand{\DOCMODULE}{Module 1 : Le Monde Vivant}
\newcommand{\DOCLECON}{NA02}
\newcommand{\DOCTITRE}{Génétique et hérédité humaines : éradication des préjugés autour des anomalies}
% OnBuch+ — préambule « cours signés » ENRICHI (compilé avec tectonic / XeLaTeX)
% Système visuel « Pop » d'OnBuch+ : fond crème (jamais blanc), bordures encre
% épaisses, ombres « dures » décalées, titres Archivo Black en capitales.
% Version MATHÉMATIQUES : graphes (pgfplots/tikz), boîtes pédagogiques
% (définition, propriété, méthode, attention, le savais-tu, exercice résolu,
% exercice, TP, bilan), page de garde et signature OnBuch+ ancrée en bas.
%
% Macros attendues dans main.tex AVANT \input{preamble} :
%   \DOCMATIERE  \DOCNIVEAU  \DOCMODULE  \DOCLECON (numéro)  \DOCTITRE
\documentclass[11pt]{article}

\usepackage{fontspec}
\usepackage[french]{babel}
\usepackage[a4paper,margin=2cm,top=3.4cm,bottom=2.8cm,headheight=1.4cm,footskip=1.3cm]{geometry}
\usepackage{amsmath,amssymb,amsfonts}
\usepackage{mathtools} % \xmapsto, \coloneqq... souvent utilisés par les modèles
\usepackage{cancel} % simplifications barrées (bilans de forces)
% \abs{z} (module, valeur absolue) et \norm{u} (norme), utilisés par les modèles
\providecommand{\abs}{}\let\abs\relax\DeclarePairedDelimiter{\abs}{\lvert}{\rvert}
\providecommand{\norm}{}\let\norm\relax\DeclarePairedDelimiter{\norm}{\lVert}{\rVert}
\usepackage[table]{xcolor} % [table] : fournit aussi \rowcolors (alternance de lignes)
\usepackage{pagecolor}
\usepackage{tikz}
\usetikzlibrary{arrows.meta,positioning,calc,shapes.geometric,shapes.misc,shapes.arrows,decorations.pathmorphing,decorations.pathreplacing,decorations.markings,patterns,shadows,mindmap,trees,fit,backgrounds,matrix,intersections,angles,quotes}
\usepackage{pgfplots}
\usepackage[version=4]{mhchem} % équations nucléaires \ce{^{235}_{92}U}
\usepackage[european,siunitx]{circuitikz} % schémas électriques
\pgfplotsset{compat=1.18}
\usepgfplotslibrary{fillbetween,groupplots} % aires entre courbes (calcul intégral)
\usepackage{enumitem}
\usepackage{fancyhdr}
% [most] charge toutes les bibliothèques tcolorbox (listings, minted, raster,
% documentation...) et gonfle inutilement la mémoire TeX sur un document long
% (~300 boîtes) : on ne charge que ce qui est réellement utilisé.
\usepackage[skins,breakable]{tcolorbox}
\usepackage{graphicx}
\usepackage{colortbl}
\usepackage{booktabs}
\usepackage{tabularx}
\usepackage{diagbox} % case d'en-tête diagonale des tableaux à double entrée
% Colonnes extensibles centrée / gauche / droite (C, L, R), souvent utilisées par les modèles
\newcolumntype{C}{>{\centering\arraybackslash}X}
\newcolumntype{L}{>{\raggedright\arraybackslash}X}
\newcolumntype{R}{>{\raggedleft\arraybackslash}X}
\usepackage{array}
\usepackage{multicol}
\usepackage{siunitx}
% parse-numbers=false : accepte aussi \SI{2,5 \times 10^{-3}}{...}
\sisetup{locale=FR,output-decimal-marker={,},group-minimum-digits=4,parse-numbers=false}
\DeclareSIUnit\ohm{\ensuremath{\symup{\Omega}}} % U+2126 absent de la police
\DeclareSIUnit\division{div} % réglages d'oscilloscope (ms/div, V/div)
\DeclareSIUnit\tr{tr} % tours (vitesse de rotation, ex. \SI{1500}{\tr\per\minute})
\DeclareSIUnit\ch{ch} % cheval-vapeur (puissance moteur, unité usuelle hors SI)
\DeclareSIUnit\franccfa{FCFA} % franc CFA (coûts, contexte camerounais)
\DeclareSIUnit\dioptre{\ensuremath{\delta}} % vergence d'une lentille (dioptrie)
\usepackage{qrcode}
% \degree : commande fréquemment utilisée par les modèles pour un angle ou une
% température en dehors de \SI{}{} (non fournie par les paquets chargés).
\providecommand{\degree}{^{\circ}}
% \qed (fin de démonstration), utilisé par les modèles sans amsthm
\providecommand{\qed}{\hfill$\square$}
% \barre : double barre des valeurs interdites dans un tableau de variations (tkz-tab)
\providecommand{\barre}{\ensuremath{\|}}
% environnement proof (amsthm non chargé) utilisé par les modèles
\ifdefined\proof\else\newenvironment{proof}[1][Démonstration]{\par\noindent\textit{#1.}\ }{\qed\par}\fi
\usepackage{lastpage}
\usepackage{hyperref}
\hypersetup{hidelinks}

% ============================================================
% Palette « Pop » OnBuch+ — mêmes hex que la classe P (plus_ui.dart)
% ============================================================
\definecolor{poporange}{HTML}{F59321}
\definecolor{popdark}{HTML}{E07A0C}
\definecolor{popcream}{HTML}{FFF6E8}
\definecolor{popink}{HTML}{1C1714}
\definecolor{poppurple}{HTML}{7A5AE0}
\definecolor{popgreen}{HTML}{2E9E6A}
\definecolor{poppink}{HTML}{E0457B}
\definecolor{popblue}{HTML}{2F7DE1}
\definecolor{popgold}{HTML}{C98A12}
\definecolor{popmuted}{HTML}{8A7F76}
\definecolor{popline}{HTML}{EADFD2}
\definecolor{popgray}{HTML}{8A7F76}
\colorlet{popgrey}{popgray}
% Alias tolérants (le modèle nomme parfois les couleurs différemment)
\colorlet{popwhite}{white}
\colorlet{popblack}{popink}
\colorlet{popred}{poppink}
\colorlet{popyellow}{popgold}
% Teintes claires pour fonds de tableaux / zones
\colorlet{popblueL}{popblue!10!white}
\colorlet{poporangeL}{poporange!14!white}
\colorlet{popgreenL}{popgreen!12!white}
\colorlet{poppurpleL}{poppurple!12!white}
\colorlet{poppinkL}{poppink!12!white}
\colorlet{popgoldL}{popgold!14!white}

\pagecolor{popcream}

% ============================================================
% Typographie
% ============================================================
\defaultfontfeatures{Ligatures=TeX}
\setmainfont{PlusJakartaSans-Medium.ttf}[Path=../../../fonts/,BoldFont=PlusJakartaSans-Bold.ttf]
% Maths en sans-serif (Fira Math) avec chiffres et lettres droites de Plus
% Jakarta, pour que les formules se fondent dans le texte.
\usepackage[math-style=ISO,bold-style=ISO]{unicode-math}
\setmathfont{FiraMath-Regular.otf}[Path=../../../fonts/]
\setmathfont{PlusJakartaSans-Medium.ttf}[Path=../../../fonts/,range={up/num,up/{Latin,latin},\mathup}]
\usepackage{icomma} % virgule décimale sans espace en mode math
% Caractères mathématiques Unicode absents des polices de texte (Plus Jakarta,
% Archivo Black) : rendus par leur équivalent mathématique, en texte comme en
% formule — sinon ils s'affichent en carré vide (ex. « Arithmétique dans ℤ »).
\usepackage{newunicodechar}
\newunicodechar{ℝ}{\ensuremath{\mathbb{R}}}
\newunicodechar{ℤ}{\ensuremath{\mathbb{Z}}}
\newunicodechar{ℂ}{\ensuremath{\mathbb{C}}}
\newunicodechar{ℕ}{\ensuremath{\mathbb{N}}}
\newunicodechar{ℚ}{\ensuremath{\mathbb{Q}}}
\newunicodechar{𝔻}{\ensuremath{\mathbb{D}}}
\newunicodechar{α}{\ensuremath{\alpha}}
\newunicodechar{β}{\ensuremath{\beta}}
\newunicodechar{γ}{\ensuremath{\gamma}}
\newunicodechar{δ}{\ensuremath{\delta}}
\newunicodechar{ε}{\ensuremath{\varepsilon}}
\newunicodechar{θ}{\ensuremath{\theta}}
\newunicodechar{λ}{\ensuremath{\lambda}}
\newunicodechar{μ}{\ensuremath{\mu}}
\newunicodechar{σ}{\ensuremath{\sigma}}
\newunicodechar{τ}{\ensuremath{\tau}}
\newunicodechar{φ}{\ensuremath{\varphi}}
\newunicodechar{ψ}{\ensuremath{\psi}}
\newunicodechar{ω}{\ensuremath{\omega}}
\newunicodechar{Σ}{\ensuremath{\Sigma}}
% majuscules produites par \MakeUppercase dans les titres (α-aminés -> Α)
\newunicodechar{Α}{\ensuremath{\alpha}}
\newunicodechar{Β}{\ensuremath{\beta}}
\newunicodechar{Γ}{\ensuremath{\Gamma}}
\newunicodechar{Δ}{\ensuremath{\Delta}}
\newunicodechar{Ε}{\ensuremath{\varepsilon}}
\newunicodechar{Θ}{\ensuremath{\Theta}}
\newunicodechar{Λ}{\ensuremath{\Lambda}}
\newunicodechar{Μ}{\ensuremath{\mu}}
\newunicodechar{Τ}{\ensuremath{\tau}}
\newunicodechar{Φ}{\ensuremath{\Phi}}
\newunicodechar{Ψ}{\ensuremath{\Psi}}
\newunicodechar{Ω}{\ensuremath{\Omega}}
\newunicodechar{↦}{\ensuremath{\mapsto}}
\newunicodechar{∈}{\ensuremath{\in}}
\newunicodechar{∉}{\ensuremath{\notin}}
\newunicodechar{∧}{\ensuremath{\wedge}}
\newunicodechar{⇔}{\ensuremath{\Leftrightarrow}}
\newunicodechar{⟺}{\ensuremath{\Longleftrightarrow}}
\newunicodechar{⇒}{\ensuremath{\Rightarrow}}
\newunicodechar{⟹}{\ensuremath{\Longrightarrow}}
\newunicodechar{∘}{\ensuremath{\circ}}
\newunicodechar{∪}{\ensuremath{\cup}}
\newunicodechar{∩}{\ensuremath{\cap}}
\newunicodechar{≡}{\ensuremath{\equiv}}
\newunicodechar{⊂}{\ensuremath{\subset}}
\newunicodechar{⊥}{\ensuremath{\perp}}
\newunicodechar{∃}{\ensuremath{\exists}}
\newunicodechar{∀}{\ensuremath{\forall}}
\newunicodechar{∛}{\ensuremath{\sqrt[3]{\,}}}
\newunicodechar{∜}{\ensuremath{\sqrt[4]{\,}}}
\newunicodechar{⃗}{\ensuremath{{}^{\to}}}
\newunicodechar{ⁿ}{\ifmmode^{n}\else\textsuperscript{\ensuremath{n}}\fi}
\newunicodechar{ˣ}{\ifmmode^{x}\else\textsuperscript{\ensuremath{x}}\fi}
\newunicodechar{ᵇ}{\ifmmode^{b}\else\textsuperscript{\ensuremath{b}}\fi}
\newunicodechar{ʳ}{\ifmmode^{r}\else\textsuperscript{\ensuremath{r}}\fi}
\newunicodechar{ᵘ}{\ifmmode^{u}\else\textsuperscript{\ensuremath{u}}\fi}
\newunicodechar{ʸ}{\ifmmode^{y}\else\textsuperscript{\ensuremath{y}}\fi}
\newunicodechar{ᵃ}{\ifmmode^{a}\else\textsuperscript{\ensuremath{a}}\fi}
\newunicodechar{ᵉ}{\ifmmode^{e}\else\textsuperscript{\ensuremath{e}}\fi}
\newunicodechar{ᵏ}{\ifmmode^{k}\else\textsuperscript{\ensuremath{k}}\fi}
\newunicodechar{ᵖ}{\ifmmode^{p}\else\textsuperscript{\ensuremath{p}}\fi}
\newunicodechar{ᴺ}{\ifmmode^{N}\else\textsuperscript{\ensuremath{N}}\fi}
\newunicodechar{ᶜ}{\ifmmode^{c}\else\textsuperscript{\ensuremath{c}}\fi}
\newunicodechar{ʰ}{\ifmmode^{h}\else\textsuperscript{\ensuremath{h}}\fi}
\newunicodechar{ᵠ}{\ifmmode^{q}\else\textsuperscript{\ensuremath{q}}\fi}
\newunicodechar{ₙ}{\ifmmode_{n}\else\textsubscript{\ensuremath{n}}\fi}
\newunicodechar{ₐ}{\ifmmode_{a}\else\textsubscript{\ensuremath{a}}\fi}
\newunicodechar{ᵢ}{\ifmmode_{i}\else\textsubscript{\ensuremath{i}}\fi}
\newunicodechar{ᵣ}{\ifmmode_{r}\else\textsubscript{\ensuremath{r}}\fi}
\newunicodechar{ₖ}{\ifmmode_{k}\else\textsubscript{\ensuremath{k}}\fi}
\newunicodechar{ᵦ}{\ifmmode_{\beta}\else\textsubscript{\ensuremath{\beta}}\fi}
\newunicodechar{ₓ}{\ifmmode_{x}\else\textsubscript{\ensuremath{x}}\fi}
\newfontfamily\popdisplay{ArchivoBlack-Regular.ttf}[Path=../../../fonts/]
\newfontfamily\poplogo{SpaceGrotesk-Bold.ttf}[Path=../../../fonts/]
\newfontfamily\popsemi{PlusJakartaSans-SemiBold.ttf}[Path=../../../fonts/]
\newfontfamily\popxbold{PlusJakartaSans-ExtraBold.ttf}[Path=../../../fonts/]
\newfontfamily\popxboldls{PlusJakartaSans-ExtraBold.ttf}[Path=../../../fonts/,LetterSpace=3.0]

\setlist[enumerate]{leftmargin=1.7em,itemsep=5pt,topsep=4pt}
\setlist[itemize]{leftmargin=1.7em,itemsep=4pt,topsep=4pt,label={\color{poporange}\textbullet}}
\setlength{\parindent}{0pt}
\setlength{\parskip}{5pt}
\renewcommand{\arraystretch}{1.25}

% pgfplots : style Pop par défaut
\pgfplotsset{
  every axis/.append style={axis line style={popink,line width=1pt},tick style={popink},
    grid style={popline,line width=0.5pt},label style={font=\small},tick label style={font=\footnotesize}},
  popcurve/.style={poporange,line width=1.8pt,no markers,smooth},
}
% Même style disponible dans un \draw TikZ simple (hors pgfplots)
\tikzset{popcurve/.style={poporange,line width=1.8pt}}
\tikzset{popgrid/.style={popline,very thin}}
\colorlet{popcurve}{poporange} % le nom du style est parfois utilisé comme couleur

% ============================================================
% Pill / squiggle
% ============================================================
\newcommand{\poppill}[3]{%
  \tikz[baseline=-0.55ex]\node[draw=popink,line width=1.2pt,rounded corners=2.6pt,%
    fill=#2,inner xsep=6.5pt,inner ysep=3pt]{%
    \popxboldls\fontsize{7}{7}\selectfont\color{#3}\MakeUppercase{#1}};%
}
\newcommand{\popsquiggle}[1]{%
  \tikz[baseline=-0.4ex]{
    \draw[#1,line width=2.6pt,line cap=round]
      plot [smooth] coordinates {(0,0.09) (0.34,0.01) (0.68,0.21) (1.02,0.01) (1.36,0.10)};
  }%
}
% Mot-clé surligné (marqueur orange sous le texte)
\newcommand{\cle}[1]{{\popxbold\color{popdark}#1}}

% ============================================================
% En-tête / pied
% ============================================================
\pagestyle{fancy}
\fancyhf{}
\renewcommand{\headrulewidth}{0pt}
\renewcommand{\footrulewidth}{0pt}
\lhead{\raisebox{-0.15ex}{\poplogo\fontsize{13}{13}\selectfont\color{popink}OnBuch\color{poporange}+}}
\rhead{\poppill{\DOCMATIERE\ \textbullet\ \DOCNIVEAU\ \textbullet\ Leçon \DOCLECON}{white}{popink}}
\lfoot{\popxbold\fontsize{7.6}{7.6}\selectfont\color{popmuted}\MakeUppercase{Cours signé OnBuch+}\ \textbullet\ {\popsemi\DOCTITRE}}
\rfoot{\tikz[baseline=-0.6ex]\node[rounded corners=8.5pt,fill=popink,minimum height=17pt,minimum width=17pt,inner xsep=5pt,inner sep=0]{\popxbold\fontsize{8}{8}\selectfont\color{popcream}\thepage\,/\,\pageref*{LastPage}};}

% ============================================================
% Titres
% ============================================================
\newcommand{\chaptitre}[2]{%
  \begin{tcolorbox}[enhanced,colback=poporange,colframe=popink,boxrule=2.6pt,arc=11pt,
      drop shadow=popink,left=15pt,right=15pt,top=12pt,bottom=11pt,before skip=3pt,after skip=18pt]
    \poppill{#2}{popink}{popcream}\par\vspace{9pt}
    {\raggedright\hyphenpenalty=10000\exhyphenpenalty=10000\popdisplay\fontsize{22}{24}\selectfont\color{popcream}\MakeUppercase{#1}\par}
  \end{tcolorbox}
}
% Section numérotée : pastille orange avec le numéro + titre Archivo
\newcounter{popsec}
\newcommand{\coursec}[1]{%
  \stepcounter{popsec}\setcounter{popsub}{0}%
  \par\vspace{16pt}\noindent
  \begin{minipage}{\linewidth}
  \tikz[baseline=-0.7ex]\node[draw=popink,line width=1.6pt,fill=poporange,rounded corners=4pt,minimum size=22pt,inner sep=0]{\popdisplay\fontsize{12}{12}\selectfont\color{popcream}\thepopsec};%
  \hspace{8pt}{\popdisplay\fontsize{15}{17}\selectfont\color{popink}\MakeUppercase{#1}}\par
  \vspace{4pt}\noindent{\color{poporange}\rule{36pt}{3.6pt}}
  \end{minipage}\par\nopagebreak\vspace{8pt}
}
\newcounter{popsub}
\newcommand{\courssub}[1]{%
  \stepcounter{popsub}%
  \par\vspace{9pt}\noindent{\popxbold\fontsize{11.5}{13}\selectfont\color{popink}{\color{poporange}\thepopsec.\thepopsub}\hspace{6pt}#1}\par\nopagebreak\vspace{4pt}
}

% ============================================================
% Boîtes pédagogiques — même recette : fond blanc, bordure épaisse colorée,
% ombre dure encre, badge plein. Titre optionnel [#1].
% ============================================================
\tcbset{popbase/.style={breakable,enhanced,colback=white,boxrule=2.3pt,arc=9pt,
  shadow={3pt}{-3pt}{0pt}{popink},left=14pt,right=14pt,top=11pt,bottom=11pt,before skip=11pt,after skip=15pt,
  fontupper=\popsemi\fontsize{10.6}{14.5}\selectfont\color{popink},
  fontlower=\popsemi\fontsize{10.6}{14.5}\selectfont\color{popink}}}
% Coche dessinée (le glyphe ✓ n'existe pas dans les polices du document)
\newcommand{\popcheck}[1]{\tikz[baseline=-0.5ex]\draw[#1,line width=1.6pt,line cap=round,line join=round](-0.13,0.0)--(-0.04,-0.09)--(0.13,0.1);}
\providecommand{\coche}{\popcheck{popgreen}}
\providecommand{\resultat}[1]{\par\smallskip\noindent\fcolorbox{popgreen}{popgreenL}{\parbox{\dimexpr\linewidth-2\fboxsep-2\fboxrule\relax}{\textbf{Résultat.} #1}}\par\smallskip}
\newcommand{\popboxhead}[3]{\poppill{#1}{#2}{white}\ifx\relax#3\relax\else\hspace{6pt}{\popxbold\fontsize{10.6}{12}\selectfont\color{popink}#3}\fi\par\vspace{7pt}}

\newtcolorbox{aretenir}[1][]{popbase,colframe=popgreen,before upper={\popboxhead{À retenir}{popgreen}{#1}}}
\newtcolorbox{exemplebox}[1][]{popbase,colframe=poporange,before upper={\popboxhead{Exemple}{poporange}{#1}}}
\newtcolorbox{definition}[1][]{popbase,colframe=popblue,before upper={\popboxhead{Définition}{popblue}{#1}}}
\newtcolorbox{propriete}[1][]{popbase,colframe=poppurple,before upper={\popboxhead{Propriété}{poppurple}{#1}}}
\newtcolorbox{methode}[1][]{popbase,colframe=popgold,before upper={\popboxhead{Méthode}{popgold}{#1}}}
\newtcolorbox{attention}[1][]{popbase,colframe=poppink,before upper={\popboxhead{Attention}{poppink}{#1}}}
\newtcolorbox{experience}[1][]{popbase,colframe=popblue,colback=popblueL,before upper={\popboxhead{Expérience}{popblue}{#1}}}
% « Le savais-tu ? » : carte violette pleine (contexte, culture, Cameroun)
\newtcolorbox{savaistu}[1][]{popbase,colback=poppurple,colframe=popink,
  fontupper=\popsemi\fontsize{10.4}{14.2}\selectfont\color{white},
  before upper={\poppill{Le savais-tu\,?}{white}{poppurple}\ifx\relax#1\relax\else\hspace{6pt}{\popxbold\color{white}#1}\fi\par\vspace{7pt}}}
% Exercice résolu : énoncé en haut, \tcblower puis la solution sur fond crème
\newcounter{popexo}\newcounter{popexr}
\newtcolorbox{exoresolu}[1][]{popbase,colframe=poporange,colbacklower=popcream,
  segmentation style={popink,line width=1.2pt,dashed},
  before upper={\stepcounter{popexr}\popboxhead{Exercice résolu \thepopexr}{poporange}{#1}},
  before lower={\poppill{Solution}{popink}{popcream}\par\vspace{6pt}}}
% Exercice (non résolu en place) — la correction est dans la partie Corrigés
\newtcolorbox{exercice}[1][]{popbase,colframe=popink,
  before upper={\stepcounter{popexo}\popboxhead{Exercice \thepopexo}{popink}{#1}}}
% Corrigé
\newtcolorbox{corrige}[1][]{popbase,colframe=popgreen,colback=popgreenL,
  before upper={\popboxhead{Corrigé}{popgreen}{#1}}}
% Objectifs / prérequis / bilan
\newtcolorbox{objectifs}{popbase,colframe=popink,colback=poporangeL,
  before upper={\popboxhead{Objectifs de la leçon}{popink}{}}}
\newtcolorbox{prerequis}{popbase,colframe=popink,colback=popblueL,
  before upper={\popboxhead{Prérequis}{popblue}{}}}
\newtcolorbox{bilan}{popbase,colframe=popink,colback=white,boxrule=2.8pt,
  before upper={\popboxhead{Fiche bilan}{poporange}{L'essentiel à connaître pour l'examen}}}
% Figure encadrée avec légende
\newtcolorbox{popfigure}[1][]{enhanced,breakable=false,colback=white,colframe=popline,boxrule=1.4pt,arc=8pt,
  left=10pt,right=10pt,top=10pt,bottom=8pt,before skip=10pt,after skip=12pt,halign=center,
  lower separated=false,halign lower=center,
  fontlower=\popsemi\fontsize{9}{11.5}\selectfont\color{popmuted},
  #1}
\newcommand{\legende}[1]{\par\vspace{6pt}{\popsemi\fontsize{9}{11.5}\selectfont\color{popmuted}#1}}

% ============================================================
% Signature OnBuch+
% ============================================================
\newcommand{\poplogoinline}[1]{{\poplogo\fontsize{#1}{#1}\selectfont\color{popink}OnBuch\color{poporange}+}}
\newcommand{\signatureonbuch}{%
  \begin{tcolorbox}[enhanced,colback=popink,colframe=popink,boxrule=2.2pt,arc=10pt,
      drop shadow=poporange,left=14pt,right=14pt,top=10pt,bottom=10pt,before skip=0pt,after skip=0pt]
    \tikz[baseline=-0.7ex]\node[circle,fill=popgreen,draw=popcream,line width=1.3pt,minimum size=22pt,inner sep=0]{\popcheck{white}};%
    \hspace{9pt}\begin{minipage}[c]{0.62\linewidth}
      {\popxbold\fontsize{12}{13}\selectfont\color{popcream}Signé\ \textbullet\ {\poplogo OnBuch\color{poporange}+}}\par\vspace{2pt}
      {\raggedright\popsemi\fontsize{8}{10.5}\selectfont\color{popline}\MakeUppercase{Équipe pédagogique OnBuch+}\\\MakeUppercase{Conforme au programme officiel MINESEC}\par}
    \end{minipage}\hfill
    {\poplogo\fontsize{22}{22}\selectfont\color{popcream}OnBuch\color{poporange}+}
  \end{tcolorbox}%
}

% ============================================================
% Page de garde
% #1 titre · #2 matière · #3 niveau · #4 module · #5 n° leçon · #6 durée ·
% #7 description / objectifs (texte ou itemize)
% ============================================================
\newcommand{\pagegarde}[7]{%
  \thispagestyle{empty}
  \newgeometry{a4paper,margin=1.7cm,top=1.6cm,bottom=1.4cm}%
  \begin{tikzpicture}[remember picture,overlay]
    \fill[poporange!18] ([xshift=-3.2cm,yshift=-3.6cm]current page.north east) circle (4.6cm);
    \fill[poppurple!14] ([xshift=2.4cm,yshift=7.6cm]current page.south west) circle (3.8cm);
  \end{tikzpicture}%
  {\poplogo\fontsize{30}{30}\selectfont\color{popink}OnBuch\color{poporange}+}\hfill
  \raisebox{0.6ex}{\poppill{Cours signé \textbullet\ Premium}{popink}{popcream}}\par
  \vspace{1pt}\popsquiggle{poppurple}\par
  \vspace{8pt}
  \begin{tcolorbox}[enhanced,colback=poporange,colframe=popink,boxrule=2.8pt,arc=13pt,
      drop shadow=popink,left=18pt,right=18pt,top=12pt,bottom=13pt,after skip=10pt]
    \poppill{#2 \textbullet\ #3}{popink}{popcream}\hspace{5pt}\poppill{#4}{white}{popink}\par\vspace{8pt}
    {\popdisplay\fontsize{14}{14}\selectfont\color{popink}LEÇON #5}\par\vspace{4pt}
    {\raggedright\hyphenpenalty=10000\exhyphenpenalty=10000\popdisplay\fontsize{25}{27}\selectfont\color{popcream}\MakeUppercase{#1}\par}
    \vspace{8pt}
    {\popxbold\fontsize{9.5}{9.5}\selectfont\color{popink}\MakeUppercase{Durée conseillée : #6}}
  \end{tcolorbox}
  \begin{tcolorbox}[enhanced,colback=white,colframe=popink,boxrule=2.2pt,arc=10pt,
      drop shadow=popink,left=16pt,right=16pt,top=10pt,bottom=10pt,after skip=8pt]
    \poppill{Dans cette leçon}{poppurple}{white}\par\vspace{5pt}
    \popsemi\fontsize{9.3}{12.6}\selectfont\color{popink}#7
  \end{tcolorbox}
  \noindent\begin{minipage}[t]{0.487\linewidth}
    \begin{tcolorbox}[enhanced,colback=popgreen,colframe=popink,boxrule=2.2pt,arc=10pt,
        drop shadow=popink,left=11pt,right=11pt,top=10pt,bottom=10pt,height=3.1cm,valign=center]
      \tcbox[colback=white,colframe=popink,boxrule=1.3pt,arc=5pt,boxsep=0pt,nobeforeafter,
        left=3pt,right=3pt,top=3pt,bottom=3pt,valign=center]{\qrcode[height=1.9cm]{https://play.google.com/store/apps/details?id=com.luvvix.onbuch&hl=fr}}%
      \hspace{8pt}\begin{minipage}[c]{0.5\linewidth}
        {\popdisplay\fontsize{11}{12.5}\selectfont\color{white}\MakeUppercase{Télécharge OnBuch}}\par\vspace{4pt}
        {\raggedright\popsemi\fontsize{8.4}{11}\selectfont\color{white}Gratuit \textbullet\ Google Play\\Tuteur IA, annales, concours, résultats\par}
      \end{minipage}
    \end{tcolorbox}
  \end{minipage}\hfill
  \begin{minipage}[t]{0.487\linewidth}
    \begin{tcolorbox}[enhanced,colback=poppurple,colframe=popink,boxrule=2.2pt,arc=10pt,
        drop shadow=popink,left=11pt,right=11pt,top=10pt,bottom=10pt,height=3.1cm,valign=center]
      \tikz[baseline=-0.7ex]\node[circle,fill=white,draw=popink,line width=1.3pt,minimum size=1.6cm,inner sep=0]{\popdisplay\fontsize{24}{24}\selectfont\color{poppurple}+};%
      \hspace{8pt}\begin{minipage}[c]{0.6\linewidth}
        {\popdisplay\fontsize{11}{12.5}\selectfont\color{white}\MakeUppercase{Passe à OnBuch+}}\par\vspace{4pt}
        {\raggedright\popsemi\fontsize{8.4}{11}\selectfont\color{white}Tous les cours signés, Tuteur IA illimité, fiches PDF — onglet OnBuch+ de l'app\par}
      \end{minipage}
    \end{tcolorbox}
  \end{minipage}
  \vspace{8pt}
  \popcontenu
  \vfill
  \signatureonbuch
  \restoregeometry
  \newpage
}

% Rangée « Ce cours contient » (page de garde)
\newcommand{\poptile}[3]{% #1 couleur #2 gros texte #3 légende
  \begin{tcolorbox}[enhanced,colback=#1,colframe=popink,boxrule=2pt,arc=9pt,shadow={2.5pt}{-2.5pt}{0pt}{popink},
      left=6pt,right=6pt,top=9pt,bottom=9pt,halign=center,valign=center,height=1.75cm,width=\linewidth,nobeforeafter]
    {\popdisplay\fontsize{12.5}{13}\selectfont\color{white}\MakeUppercase{#2}}\par\vspace{3pt}
    {\centering\popsemi\fontsize{7.8}{9.5}\selectfont\color{white}#3\par}
  \end{tcolorbox}}
\newcommand{\popcontenu}{%
  \par\noindent{\popxboldls\fontsize{7.5}{7.5}\selectfont\color{popmuted}CE COURS SIGNÉ CONTIENT}\par\vspace{7pt}
  \noindent\begin{minipage}[t]{0.235\linewidth}\poptile{popblue}{Cours}{complet, illustré, pas à pas}\end{minipage}\hfill
  \begin{minipage}[t]{0.235\linewidth}\poptile{poporange}{Méthodes}{et exercices résolus}\end{minipage}\hfill
  \begin{minipage}[t]{0.235\linewidth}\poptile{poppink}{Exercices}{type examen + corrigés}\end{minipage}\hfill
  \begin{minipage}[t]{0.235\linewidth}\poptile{popgreen}{Fiche bilan}{l'essentiel à retenir}\end{minipage}\par}

% Sommaire
\newtcolorbox{sommaire}{enhanced,colback=white,colframe=popink,boxrule=2.2pt,arc=10pt,
  drop shadow=popink,left=18pt,right=18pt,top=13pt,bottom=13pt,before skip=6pt,after skip=20pt,
  before upper={\poppill{Sommaire}{poppurple}{white}\par\vspace{9pt}\popsemi\fontsize{10.6}{14.5}\selectfont\color{popink}}}

% Signature de fin de cours — poussée tout en bas de la dernière page
\newcommand{\findecours}{%
  \par\vfill\nopagebreak
  \begin{center}\popsquiggle{poporange}\end{center}\vspace{4pt}
  \signatureonbuch
}
\AtBeginDocument{\def\llbracket{\mathopen{[\mkern-3.5mu[}}\def\rrbracket{\mathclose{]\mkern-3.5mu]}}} % intervalles d'entiers (glyphe ⟦ absent de la police)
% Glyphes absents de Fira Math : redéfinis à partir de glyphes disponibles
\AtBeginDocument{%
  \def\setminus{\mathbin{\backslash}}\let\smallsetminus\setminus
  \def\vdots{\mathord{\vbox{\baselineskip3.2pt\lineskiplimit0pt\kern4pt\hbox{.}\hbox{.}\hbox{.}}}}%
  \def\ddots{\mathinner{\mkern1mu\raise6.5pt\hbox{.}\mkern2mu\raise3.6pt\hbox{.}\mkern2mu\raise0.7pt\hbox{.}\mkern1mu}}%
  \def\triangle{\mathord{\scalebox{0.9}{$\symup{\Delta}$}}}%
  \def\nabla{\mathord{\raisebox{\depth}{\scalebox{1}[-1]{$\symup{\Delta}$}}}}%
  \def\checkmark{\popcheck{popdark}}%
  \def\star{\mathbin{\ast}}\def\bigstar{\mathord{\ast}}%
  \def\diamond{\mathbin{\scalebox{0.6}{$\Diamond$}}}%
  \def\bigcup{\mathop{\mathchoice{\raisebox{-0.3em}{\scalebox{1.7}{$\cup$}}}{\scalebox{1.25}{$\cup$}}{\cup}{\cup}}\displaylimits}%
  \def\bigcap{\mathop{\mathchoice{\raisebox{-0.3em}{\scalebox{1.7}{$\cap$}}}{\scalebox{1.25}{$\cap$}}{\cap}{\cap}}\displaylimits}%
  \def\lesseqqgtr{\mathrel{\lessgtr}}\def\bigoplus{\mathop{\oplus}\displaylimits}%
}
\providecommand{\ssi}{si et seulement si} % abréviation fréquente
\AtBeginDocument{\providecommand{\wideparen}{\overparen}} % arc de cercle

%% FIGFIX-BEGIN v5 — correctifs globaux des figures (docs/onbuch-generation/tools/figfix)
\usepackage{adjustbox}\usepackage{etoolbox}\tcbuselibrary{raster}
% toute figure TikZ/pgfplots plus large que la ligne est réduite à la largeur disponible
\BeforeBeginEnvironment{tikzpicture}{\begin{adjustbox}{max width=\linewidth}}
\AfterEndEnvironment{tikzpicture}{\end{adjustbox}}
% légendes pgfplots lisibles par-dessus les courbes
\pgfplotsset{every axis/.append style={legend style={fill=white,fill opacity=0.92,text opacity=1,draw=black!15,inner sep=3pt}}}
% étiquettes d'arêtes (syntaxe edge["..."]) avec fond blanc
\tikzset{every edge quotes/.append style={fill=white,inner sep=1.2pt}}
%% FIGFIX-END
\begin{document}

\pagegarde{Génétique et hérédité humaines : éradication des préjugés autour des anomalies}{SVT}{1ère A}{Module 1 : Le Monde Vivant}{NA02}{3 h}{Cette leçon étudie le support chromosomique de l'hérédité humaine : le caryotype, la détermination du sexe et les mécanismes de transmission des caractères et anomalies (daltonisme, hémophilie, albinisme, drépanocytose). Elle fournit les outils scientifiques pour interpréter des arbres généalogiques et combattre les préjugés liés aux anomalies héréditaires.
\vspace{4pt}
\begin{multicols}{2}\small
\begin{itemize}
\item À la fin de cette leçon, je sais définir le caryotype et décrire ses caractéristiques chez l'Homme
\item À la fin de cette leçon, je sais identifier le caryotype de cellules humaines à partir d'une photographie de chromosomes
\item À la fin de cette leçon, je sais identifier le sexe d'un individu à partir du caryotype de ses cellules
\item À la fin de cette leçon, je sais écrire et interpréter une formule chromosomique normale ou anormale
\item À la fin de cette leçon, je sais identifier des anomalies chromosomiques (trisomie 21) et géniques sur un document
\item À la fin de cette leçon, je sais expliquer la détermination chromosomique du sexe chez l'Homme
\end{itemize}\end{multicols}}

\begin{sommaire}
\begin{multicols}{2}
\begin{enumerate}
\item Le caryotype humain
\item Anomalies du nombre de chromosomes et formule chromosomique
\item Chromosomes et détermination du sexe
\item Hérédité liée au sexe : daltonisme et hémophilie
\item Hérédité autosomale : albinisme et drépanocytose
\item Synthèse : démarche d'interprétation et lutte contre les préjugés
\item Activité d'intégration
\item Méthodes et exercices résolus
\item Exercices
\item Corrigés des exercices
\item Fiche bilan
\end{enumerate}
\end{multicols}
\end{sommaire}

\begin{objectifs}
\begin{itemize}
\item À la fin de cette leçon, je sais définir le caryotype et décrire ses caractéristiques chez l'Homme
\item À la fin de cette leçon, je sais identifier le caryotype de cellules humaines à partir d'une photographie de chromosomes
\item À la fin de cette leçon, je sais identifier le sexe d'un individu à partir du caryotype de ses cellules
\item À la fin de cette leçon, je sais écrire et interpréter une formule chromosomique normale ou anormale
\item À la fin de cette leçon, je sais identifier des anomalies chromosomiques (trisomie 21) et géniques sur un document
\item À la fin de cette leçon, je sais expliquer la détermination chromosomique du sexe chez l'Homme
\item À la fin de cette leçon, je sais interpréter un arbre généalogique simple pour le daltonisme et l'hémophilie (hérédité liée au sexe)
\item À la fin de cette leçon, je sais interpréter un arbre généalogique simple pour l'albinisme et la drépanocytose (hérédité autosomale)
\item À la fin de cette leçon, je sais argumenter scientifiquement contre les préjugés entourant les anomalies héréditaires
\end{itemize}
\end{objectifs}

\begin{prerequis}
\begin{itemize}
\item Notions de cellule, noyau et chromosomes (classes de collège)
\item Notion de gamètes, fécondation et brassage génétique (début du Module 1 en Première)
\item Notions de gène et d'allèle (introduction à la génétique)
\item Lecture et construction de tableaux de croisement (échiquier de croisement)
\end{itemize}
\end{prerequis}

\newpage

\coursec{Le caryotype humain}

\textbf{Situation d'accroche.} Dans un quartier de Yaoundé, la famille Ngo Bell accueille un bébé albinos : la peau très claire, les cheveux presque blancs, les yeux sensibles à la lumière. Aussitôt, les rumeurs circulent : certains voisins accusent la mère d'infidélité, d'autres parlent de sorcellerie ou de « mauvais sort ». Quelques mois plus tard, on apprend que le fils aîné d'une autre famille ne distingue pas le rouge du vert, exactement comme son grand-père maternel, alors que son père, lui, voit parfaitement les couleurs. Comment expliquer \emph{scientifiquement} la transmission de ces caractères ? Peut-on prévoir leur apparition et démonter les préjugés qui entourent ces anomalies ?

Pour répondre à ces questions, il faut d'abord explorer le « support matériel » de l'hérédité : les \cle{chromosomes}, que l'on observe grâce à une image très particulière appelée \cle{caryotype}. C'est l'objet de cette première section.

\courssub{Qu'est-ce qu'un caryotype ?}

\textbf{Intuition.} Imagine que tu veuilles ranger une bibliothèque de 46 livres écrits en deux exemplaires. Pour vérifier qu'aucun livre ne manque et qu'aucun n'est en trop, on les photographie, on découpe les images et on les classe par paires, du plus gros au plus petit. Le caryotype, c'est exactement cela, appliqué aux chromosomes d'une cellule humaine : une photographie rangée et ordonnée qui permet de tout compter et de tout vérifier.

\begin{definition}[Caryotype]
Le \cle{caryotype} est la \textbf{représentation ordonnée et classée de l'ensemble des chromosomes d'une cellule}, observés en métaphase de mitose. Les chromosomes sont photographiés, découpés, puis appariés (regroupés par paires de chromosomes homologues) et classés par taille décroissante.
\end{definition}

\begin{attention}[Chromosome, chromatide : ne pas confondre]
En métaphase, chaque chromosome est déjà \textbf{dupliqué} : il est formé de \textbf{deux chromatides sœurs} identiques, réunies par un point de constriction appelé \cle{centromère}. Un chromosome dupliqué (deux chromatides) compte pour \textbf{un seul} chromosome. En revanche, deux \textbf{chromosomes homologues} (l'un d'origine paternelle, l'autre d'origine maternelle) sont deux chromosomes distincts qui forment \textbf{une paire}. Piège classique au Probatoire : compter les chromatides au lieu des chromosomes fait croire à tort qu'une cellule humaine possède 92 chromosomes !
\end{attention}

\courssub{Comment obtient-on un caryotype ?}

L'observation directe des chromosomes n'est possible qu'au moment où ils sont le plus condensés et le plus visibles : la \textbf{métaphase} de la division cellulaire (mitose). La technique, réalisée en laboratoire d'analyses médicales, suit un protocole précis.

\begin{methode}[Réalisation d'un caryotype humain]
\begin{enumerate}
\item \textbf{Prélèvement} : on prélève un peu de sang, dont on isole les \cle{lymphocytes} (globules blancs), cellules faciles à cultiver.
\item \textbf{Culture} : les lymphocytes sont placés en milieu de culture pendant environ 72 heures, avec une substance (phytohémagglutinine) qui stimule leur division.
\item \textbf{Blocage en métaphase} : on ajoute de la \cle{colchicine}, qui empêche la polymérisation des microtubules du fuseau mitotique et bloque ainsi les cellules en métaphase, moment où les chromosomes sont courts, épais et bien individualisés.
\item \textbf{Étalement et coloration} : les cellules sont soumises à un choc hypotonique (éclatement), étalées sur lame, puis colorées (coloration de Giemsa), ce qui fait apparaître des bandes caractéristiques sur chaque chromosome (bandes G).
\item \textbf{Photographie et classement} : on photographie les chromosomes au microscope, on les découpe, on les apparie par paires d'homologues et on les classe par taille décroissante, de la paire 1 à la paire 22, puis la paire des chromosomes sexuels.
\end{enumerate}
\end{methode}

\begin{savaistu}[Le caryotype au Cameroun]
La technique du caryotype est pratiquée dans les grands hôpitaux camerounais (Centre Hospitalier et Universitaire de Yaoundé, Hôpital Général de Douala) et dans certains laboratoires privés. Elle sert notamment au \textbf{diagnostic prénatal} (recherche d'anomalies chromosomiques chez le fœtus, comme la trisomie 21) et au \textbf{bilan des stérilités} du couple : certaines difficultés à concevoir s'expliquent par des anomalies du nombre ou de la structure des chromosomes chez l'un des partenaires. Loin de la sorcellerie, c'est une science au service des familles !
\end{savaistu}

\courssub{Caractéristiques du caryotype humain normal}

L'observation de nombreux caryotypes humains, hommes et femmes, conduit à une conclusion fondamentale.

\begin{propriete}[Constitution chromosomique de l'espèce humaine (admise)]
Le noyau de toute \textbf{cellule somatique} humaine (cellule du corps, non reproductrice) possède \textbf{46 chromosomes répartis en 23 paires} : on dit qu'elle est \textbf{diploïde} ($2n = 46$). Les \textbf{gamètes} (spermatozoïdes et ovules) ne possèdent que \textbf{23 chromosomes}, un par paire : ils sont \textbf{haploïdes} ($n = 23$). Parmi les 23 paires :
\begin{itemize}
\item \textbf{22 paires d'autosomes} : chromosomes communs aux deux sexes, numérotés de 1 à 22 ;
\item \textbf{1 paire de gonosomes} (ou \cle{hétérochromosomes}) : les chromosomes sexuels X et Y, qui déterminent le sexe.
\end{itemize}
\end{propriete}

\begin{definition}[Autosome et gonosome]
Un \cle{autosome} est un chromosome qui n'intervient pas dans la détermination du sexe ; il existe en deux exemplaires identiques chez l'homme comme chez la femme (paires 1 à 22). Un \cle{gonosome} (ou hétérochromosome) est un chromosome sexuel : la femme possède deux chromosomes X (paire XX), l'homme un chromosome X et un chromosome Y (paire XY).
\end{definition}

\textbf{Homologie des chromosomes d'une paire.} Deux chromosomes d'une même paire sont dits \cle{homologues} : ils ont la \textbf{même taille}, la \textbf{même position du centromère}, le même aspect après coloration (même profil de bandes), et surtout ils portent les \textbf{mêmes gènes} aux mêmes emplacements (locus). Attention toutefois : pour un gène donné, les deux homologues peuvent porter des versions différentes de ce gène, appelées \cle{allèles}. Par exemple, le gène intervenant dans la production de mélanine (pigment de la peau) existe en version « normale » et en version « albinos » : c'est la clé pour comprendre le bébé de la famille Ngo Bell, comme nous le verrons dans la suite de la leçon.

\textbf{Classification des chromosomes.} Les chromosomes sont classés selon deux critères :
\begin{itemize}
\item la \textbf{taille}, décroissante de la paire 1 (la plus grande) à la paire 22 ;
\item la \textbf{position du centromère}, qui définit trois grands types morphologiques.
\end{itemize}

\begin{definition}[Types de chromosomes selon le centromère]
Un chromosome est dit \cle{métacentrique} si son centromère est au milieu (deux bras égaux), \cle{submétacentrique} si le centromère est décalé (un bras court et un bras long), \cle{acrocentrique} si le centromère est très proche d'une extrémité (un bras très court, un bras très long).
\end{definition}

\begin{popfigure}
\begin{center}
\begin{tikzpicture}[scale=0.62]
% Métacentrique
\draw[fill=popblue!60,draw=popdark] (-0.25,0) rectangle (0.25,2.2);
\draw[fill=popblue!60,draw=popdark] (-0.25,2.2) rectangle (0.25,4.4);
\draw[fill=popblue!60,draw=popdark] (0.45,0) rectangle (0.95,2.2);
\draw[fill=popblue!60,draw=popdark] (0.45,2.2) rectangle (0.95,4.4);
\fill[poporange] (0.35,2.2) circle (0.13);
\node[align=center] at (0.35,-0.9) {\small Métacentrique\\ \small (centromère médian)};
% Submétacentrique
\begin{scope}[xshift=4.2cm]
\draw[fill=popgreen!60,draw=popdark] (-0.25,0) rectangle (0.25,1.5);
\draw[fill=popgreen!60,draw=popdark] (-0.25,1.5) rectangle (0.25,4.4);
\draw[fill=popgreen!60,draw=popdark] (0.45,0) rectangle (0.95,1.5);
\draw[fill=popgreen!60,draw=popdark] (0.45,1.5) rectangle (0.95,4.4);
\fill[poporange] (0.35,1.5) circle (0.13);
\node[align=center] at (0.35,-0.9) {\small Submétacentrique\\ \small (bras inégaux)};
\end{scope}
% Acrocentrique
\begin{scope}[xshift=8.4cm]
\draw[fill=poppurple!50,draw=popdark] (-0.25,0) rectangle (0.25,0.7);
\draw[fill=poppurple!50,draw=popdark] (-0.25,0.7) rectangle (0.25,4.4);
\draw[fill=poppurple!50,draw=popdark] (0.45,0) rectangle (0.95,0.7);
\draw[fill=poppurple!50,draw=popdark] (0.45,0.7) rectangle (0.95,4.4);
\fill[poporange] (0.35,0.7) circle (0.13);
\node[align=center] at (0.35,-0.9) {\small Acrocentrique\\ \small (centromère presque terminal)};
\end{scope}
% Légende centromère
\fill[poporange] (11.4,3.9) circle (0.13);
\draw[popdark] (11.4,3.9) -- (12.6,3.9) node[right,align=left] {\small centromère};
\end{tikzpicture}
\end{center}
\legende{Les trois types de chromosomes selon la position du centromère (en orange). Chaque chromosome en métaphase est formé de deux chromatides sœurs réunies par le centromère.}
\end{popfigure}

\begin{popfigure}
\begin{center}
\begin{tikzpicture}[scale=0.5]
% Autosomes : paires 1 à 11 (première rangée)
\foreach \i in {1,...,11} {
  \pgfmathsetmacro{\h}{3.4 - 0.13*(\i-1)}
  \pgfmathsetmacro{\x}{0.95*(\i-1)}
  \draw[fill=popblue!55,draw=popdark] (\x,0) rectangle (\x+0.28,\h);
  \draw[fill=popblue!55,draw=popdark] (\x+0.34,0) rectangle (\x+0.62,\h);
  \node[below] at (\x+0.31,-0.15) {\tiny \i};
}
% Autosomes : paires 12 à 22 (deuxième rangée)
\begin{scope}[yshift=-5.6cm]
\foreach \i in {12,...,22} {
  \pgfmathsetmacro{\h}{3.4 - 0.13*(\i-1)}
  \pgfmathsetmacro{\x}{0.95*(\i-12)}
  \draw[fill=popblue!55,draw=popdark] (\x,0) rectangle (\x+0.28,\h);
  \draw[fill=popblue!55,draw=popdark] (\x+0.34,0) rectangle (\x+0.62,\h);
  \node[below] at (\x+0.31,-0.15) {\tiny \i};
}
% Paire sexuelle XX encadrée
\draw[fill=poppink!70,draw=popdark] (10.9,0) rectangle (11.18,1.2);
\draw[fill=poppink!70,draw=popdark] (11.24,0) rectangle (11.52,1.2);
\node[below] at (11.2,-0.15) {\tiny X};
\draw[fill=poppink!70,draw=popdark] (11.8,0) rectangle (12.08,1.2);
\draw[fill=poppink!70,draw=popdark] (12.14,0) rectangle (12.42,1.2);
\node[below] at (12.1,-0.15) {\tiny X};
\draw[poporange,thick,rounded corners] (10.6,-0.7) rectangle (12.75,1.6);
\node[poporange,align=center] at (11.67,2.15) {\small paire sexuelle};
\end{scope}
\node[align=left,popdark] at (5.2,4.3) {\small \textbf{Autosomes : paires 1 à 22}};
\end{tikzpicture}
\end{center}
\legende{Caryotype humain schématisé d'une femme : 22 paires d'autosomes classées par taille décroissante (deux rangées), puis la paire de gonosomes XX (encadrée en orange). Chez un homme, la paire sexuelle serait XY, le Y étant beaucoup plus petit que le X.}
\end{popfigure}

\courssub{Cellule somatique ou gamète : deux lectures du caryotype}

Le nombre de chromosomes dépend du type de cellule observée. La fécondation (union d'un spermatozoïde à $n = 23$ et d'un ovule à $n = 23$) restaure le nombre diploïde : $23 + 23 = 46$. Ainsi, chaque enfant reçoit \textbf{un chromosome de chaque paire de son père} et \textbf{un chromosome de chaque paire de sa mère} : voilà pourquoi les caractères « voyagent » d'une génération à l'autre, parfois en sautant une génération, comme le daltonisme du grand-père maternel retrouvé chez son petit-fils.

\begin{center}
\begin{tabularx}{\linewidth}{|l|X|X|}
\hline
\rowcolor{popblueL} \textbf{Critère} & \textbf{Cellule somatique} & \textbf{Gamète} \\
\hline
Nombre de chromosomes & 46 ($2n$, diploïde) & 23 ($n$, haploïde) \\
\hline
Organisation & 23 paires d'homologues & 1 chromosome par paire \\
\hline
Formule chromosomique & 46,XX ou 46,XY & 23,X ou 23,Y \\
\hline
Exemples & lymphocyte, cellule de peau & spermatozoïde, ovule \\
\hline
\end{tabularx}
\end{center}

\begin{definition}[Formule chromosomique]
La \cle{formule chromosomique} résume le caryotype d'un individu : elle indique le \textbf{nombre total de chromosomes}, suivi d'une virgule, puis la \textbf{composition de la paire sexuelle} (pour une cellule somatique) ou du gonosome unique (pour un gamète). Ainsi : \textbf{46,XX} désigne une femme au caryotype normal ; \textbf{46,XY} désigne un homme au caryotype normal ; \textbf{23,X} ou \textbf{23,Y} désignent des gamètes.
\end{definition}

\begin{methode}[Identifier un caryotype]
Face à un caryotype (document fréquent au Probatoire), procède toujours dans l'ordre :
\begin{enumerate}
\item \textbf{Compter} les chromosomes (ou les paires) : 46 chromosomes = 23 paires est la normale ; 45 ou 47 signale une anomalie du nombre.
\item \textbf{Repérer la paire sexuelle} (dernière paire) : XX $\Rightarrow$ individu féminin ; XY $\Rightarrow$ individu masculin.
\item \textbf{Vérifier l'appariement des autosomes} : chaque paire doit contenir deux homologues de même taille et même centromère ; une paire à trois chromosomes (trisomie) ou à un seul (monosomie) est anormale.
\item \textbf{Conclure} par la formule chromosomique, par exemple 46,XX.
\end{enumerate}
\end{methode}

\begin{exoresolu}[Lecture d'un caryotype]
Énoncé. On réalise le caryotype des lymphocytes d'un patient. On observe 23 paires de chromosomes ; les paires 1 à 22 sont normales ; la dernière paire comporte un grand chromosome et un très petit chromosome.
\begin{enumerate}
\item Ce caryotype est-il celui d'un homme ou d'une femme ? Justifier.
\item Donner la formule chromosomique de cet individu.
\item Combien de chromosomes possède chacun de ses gamètes et quelle est leur formule chromosomique ?
\end{enumerate}
\tcblower
Solution détaillée.
\begin{enumerate}
\item La dernière paire est la paire des gonosomes. Elle comporte deux chromosomes \textbf{de tailles différentes} : un grand chromosome X et un petit chromosome Y. La paire est donc \textbf{XY} : il s'agit d'un \textbf{homme}.
\item Le nombre total de chromosomes est $23 \times 2 = 46$ et la paire sexuelle est XY. La formule chromosomique est donc \textbf{46,XY}.
\item Les gamètes (spermatozoïdes) sont haploïdes : chacun possède $n = 23$ chromosomes, soit 22 autosomes plus un gonosome. Par méiose, la paire XY se sépare : la moitié des spermatozoïdes reçoivent le X, l'autre le Y. Leurs formules chromosomiques sont donc \textbf{23,X} ou \textbf{23,Y}.
\end{enumerate}
\end{exoresolu}

\begin{attention}[Erreurs fréquentes au Probatoire]
\begin{itemize}
\item Dire « 23 chromosomes dans une cellule somatique » : faux, c'est \textbf{23 paires}, soit 46 chromosomes ($2n = 46$).
\item Confondre \textbf{paire d'homologues} et \textbf{chromosome à deux chromatides} : une paire = deux chromosomes différents (paternel + maternel) ; deux chromatides sœurs = un seul chromosome dupliqué.
\item Écrire la formule chromosomique « XX,46 » : le nombre total s'écrit \textbf{toujours en premier} : 46,XX.
\item Affirmer que le chromosome Y est un autosome : c'est un \textbf{gonosome}.
\item Oublier que les gamètes ont une formule à deux termes (ex: 23,X) et non trois (ex: 23,XX).
\end{itemize}
\end{attention}

\begin{aretenir}[L'essentiel de la section]
\begin{itemize}
\item Le \cle{caryotype} est la représentation ordonnée des chromosomes d'une cellule, obtenue après culture, blocage en métaphase (colchicine), coloration (Giemsa), photographie et classement.
\item Le caryotype humain normal comporte \textbf{46 chromosomes en 23 paires} : 22 paires d'\textbf{autosomes} et 1 paire de \textbf{gonosomes} (XX chez la femme, XY chez l'homme).
\item Les chromosomes homologues portent les mêmes gènes aux mêmes locus, mais parfois des allèles différents : c'est la base de la transmission des caractères.
\item Les cellules somatiques sont diploïdes ($2n = 46$), les gamètes haploïdes ($n = 23$).
\item La \textbf{formule chromosomique} s'écrit : 46,XX (femme) ou 46,XY (homme) pour une cellule somatique ; 23,X ou 23,Y pour un gamète.
\end{itemize}
\end{aretenir}

Maintenant que nous savons lire cette « carte d'identité chromosomique », nous pouvons, dans la section suivante, repérer ce qui se passe lorsque le nombre de chromosomes n'est pas exactement 46 : c'est le cas de certaines anomalies comme la trisomie 21.

\coursec{Anomalies du nombre de chromosomes et formule chromosomique}

Dans la section précédente, tu as découvert que le \cle{caryotype} humain normal comporte \cle{46 chromosomes} répartis en 23 paires : 22 paires d'\cle{autosomes} et une paire de \cle{gonosomes} (XX chez la fille, XY chez le garçon). Mais dans certaines familles, on observe la naissance d'enfants porteurs de caractères inhabituels, parfois source de rejet social ou de superstitions. Au Cameroun comme ailleurs, un enfant né avec une trisomie 21 est encore trop souvent présenté comme victime d'une « malédiction » ou comme la conséquence d'une faute commise par ses parents. La génétique permet de comprendre ce qui s'est réellement passé : un simple accident mécanique survenu lors de la formation des gamètes. C'est ce mécanisme que nous allons maintenant étudier, et tu verras qu'il obéit à des lois biologiques précises, sans aucune part de mystère.

\courssub{Situation-problème : un caryotype qui ne compte pas 46 chromosomes}

\textbf{Observation.} Dans un service de pédiatrie d'un hôpital de référence de Yaoundé, un nourrisson présente un ensemble de signes caractéristiques : visage arrondi, yeux bridés, petite taille, hypotonie (manque de tonus musculaire). Le médecin prescrit une analyse du caryotype. Le résultat montre \num{47} chromosomes au lieu de \num{46}, avec \textbf{trois} chromosomes 21 au lieu de deux.

\textbf{Problème.} D'où vient ce chromosome surnuméraire ? Est-ce une maladie transmise par les parents ? Est-ce la faute de quelqu'un ?

\textbf{Hypothèse.} Le chromosome supplémentaire provient d'un gamète anormal, c'est-à-dire d'un spermatozoïde ou d'un ovule qui aurait reçu deux chromosomes 21 au lieu d'un seul.

Pour vérifier cette hypothèse, il faut se rappeler ce qui se passe lors de la \cle{méiose}, la division cellulaire qui produit les gamètes.

\courssub{Origine des anomalies de nombre : la non-disjonction méiotique}

\begin{aretenir}[Rappel : la méiose]
La \cle{méiose} est la division cellulaire qui produit les gamètes (spermatozoïdes chez l'homme, ovules chez la femme). À partir d'une cellule mère à $2n = 46$ chromosomes, elle aboutit à la formation de gamètes haploïdes à $n = 23$ chromosomes : chaque gamète reçoit \textbf{un seul exemplaire} de chaque paire de chromosomes homologues. On dit que les chromosomes d'une même paire \emph{se disjoignent} (se séparent) au cours de la première division de la méiose (méiose I).
\end{aretenir}

Normalement, les deux chromosomes 21 de la cellule mère se séparent lors de la méiose I : chaque gamète reçoit un seul chromosome 21. Mais il arrive que cette séparation ne se produise pas correctement.

\begin{definition}[Non-disjonction]
La \cle{non-disjonction} est un accident de la méiose (le plus souvent lors de la \textbf{métaphase I / anaphase I}) au cours duquel les deux chromosomes homologues d'une même paire \textbf{ne se séparent pas} et migrent ensemble dans le même gamète. Il en résulte des gamètes anormaux :
\begin{itemize}
\item un gamète à $n + 1 = 24$ chromosomes (il possède le chromosome en double) ;
\item un gamète à $n - 1 = 22$ chromosomes (il ne possède pas ce chromosome).
\end{itemize}
\end{definition}

Que deviennent ces gamètes anormaux après la fécondation ? Suivons le raisonnement :

\begin{itemize}
\item Si un gamète à $24$ chromosomes (avec deux chromosomes 21) fusionne avec un gamète normal à $23$ chromosomes (un chromosome 21), le \cle{zygote} (œuf) obtenu possède $24 + 23 = 47$ chromosomes, dont \textbf{trois chromosomes 21} : c'est une \cle{trisomie} 21.
\item Si un gamète à $22$ chromosomes (sans chromosome 21) fusionne avec un gamète normal à $23$ chromosomes, le zygote possède $22 + 23 = 45$ chromosomes, dont \textbf{un seul chromosome 21} : c'est une \cle{monosomie} 21, en général non viable (la grossesse s'arrête précocement, souvent avant que la femme ne sache qu'elle est enceinte).
\end{itemize}

\begin{experience}[Démonstration guidée : compter les chromosomes après non-disjonction]
\textbf{Données.} Une cellule mère humaine contient $2n = 46$ chromosomes, soit 23 paires. Considérons uniquement la paire 21.

\textbf{Étape 1 — Méiose normale.} Les deux chromosomes 21 (homologues) se disjoignent lors de la méiose I : chaque gamète reçoit $1$ chromosome 21, soit $n = 23$ chromosomes au total.

\textbf{Étape 2 — Méiose avec non-disjonction (méiose I).} Les deux chromosomes 21 homologues migrent ensemble vers le même pôle : un gamète reçoit $2$ chromosomes 21 ($n+1 = 24$ chromosomes), l'autre n'en reçoit aucun ($n-1 = 22$ chromosomes).

\textbf{Étape 3 — Fécondation par un gamète normal ($n = 23$).}
\begin{align*}
24 + 23 &= 47 \quad \text{(zygote trisomique : trois chromosomes 21)}\\
22 + 23 &= 45 \quad \text{(zygote monosomique : un seul chromosome 21)}
\end{align*}

\textbf{Conclusion.} La trisomie 21 résulte de la fécondation d'un gamète anormal à 24 chromosomes par un gamète normal à 23 chromosomes. C'est un \textbf{accident méiotique}, survenu au hasard, sans que personne n'en soit responsable.
\end{experience}

\begin{popfigure}
\begin{tikzpicture}[scale=0.95, every node/.style={font=\small}]
% Cellule mère
\node[draw, ellipse, fill=popblueL, minimum width=2.6cm, minimum height=1.5cm] (mere) at (0,0) {};
\node[font=\small\bfseries] at (0,-1.15) {Cellule mère $2n = 46$};
% Chromosomes 21 homologues (couleurs différentes)
\draw[thick, popdark] (-0.5,0.15) -- (-0.2,0.15);
\draw[thick, popdark] (-0.5,-0.15) -- (-0.2,-0.15);
\draw[thick, poporange] (0.2,0.15) -- (0.5,0.15);
\draw[thick, poporange] (0.2,-0.15) -- (0.5,-0.15);
\node[font=\scriptsize, popmuted] at (0,0.55) {Paire 21 (homologues)};

% Branche normale (gauche)
\draw[-{Stealth}, thick, popgreen] (-0.8,-0.9) -- (-3.6,-1.9);
\node[font=\small\bfseries, popgreen] at (-4.4,-1.5) {Méiose normale};
\node[draw, ellipse, fill=popgreenL, minimum width=1.9cm, minimum height=1.1cm] (g1) at (-4.6,-2.7) {};
\node[font=\scriptsize] at (-4.6,-3.5) {Gamète $n = 23$};
\draw[thick, popdark] (-4.75,-2.7) -- (-4.45,-2.7); % un chr 21
\node[draw, ellipse, fill=popgreenL, minimum width=1.9cm, minimum height=1.1cm] (g2) at (-2.3,-2.7) {};
\node[font=\scriptsize] at (-2.3,-3.5) {Gamète $n = 23$};
\draw[thick, poporange] (-2.45,-2.7) -- (-2.15,-2.7); % un chr 21

% Branche non-disjonction (droite)
\draw[-{Stealth}, thick, poppink] (0.8,-0.9) -- (3.6,-1.9);
\node[font=\small\bfseries, poppink] at (4.5,-1.5) {Non-disjonction (Méiose I)};
\node[draw, ellipse, fill=poppinkL, minimum width=1.9cm, minimum height=1.1cm] (g3) at (2.3,-2.7) {};
\node[font=\scriptsize] at (2.3,-3.5) {Gamète $n+1 = 24$};
\draw[thick, popdark] (2.05,-2.55) -- (2.35,-2.55); % deux chr 21
\draw[thick, poporange] (2.05,-2.85) -- (2.35,-2.85);
\node[draw, ellipse, fill=poppinkL, minimum width=1.9cm, minimum height=1.1cm] (g4) at (4.6,-2.7) {};
\node[font=\scriptsize] at (4.6,-3.5) {Gamète $n-1 = 22$};
% pas de chr 21

% Fécondations
\draw[-{Stealth}, thick, popblue] (2.3,-3.6) -- (0.6,-4.6);
\draw[-{Stealth}, thick, popblue] (-2.3,-3.6) -- (-0.6,-4.6);
\node[font=\scriptsize, popblue] at (2.6,-4.35) {Fécondation};
\node[draw, ellipse, fill=poporangeL, minimum width=2.2cm, minimum height=1.2cm] (z1) at (-1.6,-5.3) {};
\node[font=\scriptsize\bfseries] at (-1.6,-6.15) {Zygote $2n+1 = 47$};
\node[font=\scriptsize, poppink] at (-1.6,-6.55) {Trisomie 21};
\draw[thick, popdark] (-1.85,-5.15) -- (-1.55,-5.15);
\draw[thick, poporange] (-1.85,-5.45) -- (-1.55,-5.45);
\draw[thick, popgreen] (-1.35,-5.3) -- (-1.05,-5.3); % 3 traits

\draw[-{Stealth}, thick, popblue] (4.6,-3.6) -- (3.2,-4.6);
\node[draw, ellipse, fill=popgoldL, minimum width=2.2cm, minimum height=1.2cm] (z2) at (2.9,-5.3) {};
\node[font=\scriptsize\bfseries] at (2.9,-6.15) {Zygote $2n-1 = 45$};
\node[font=\scriptsize, poppink] at (2.9,-6.55) {Monosomie (non viable)};
\draw[thick, popgreen] (2.75,-5.3) -- (3.05,-5.3); % 1 trait
\end{tikzpicture}
\legende{Conséquences de la non-disjonction de la paire de chromosomes 21 au cours de la méiose I. À gauche, une méiose normale produit des gamètes à $n = 23$ chromosomes (un chromosome 21 chacun). À droite, la non-disjonction produit un gamète à $n+1 = 24$ chromosomes (deux chromosomes 21) et un gamète à $n-1 = 22$ chromosomes (aucun chromosome 21). Après fécondation par un gamète normal ($n=23$), on obtient un zygote trisomique (47 chromosomes, trois chromosomes 21) ou monosomique (45 chromosomes, un chromosome 21).}
\end{popfigure}

\courssub{Exemple de la trisomie 21 : un caryotype à 47 chromosomes}

\begin{definition}[Trisomie 21]
La \cle{trisomie 21} est une anomalie chromosomique caractérisée par la présence de \textbf{trois chromosomes 21} au lieu de deux dans toutes les cellules de l'organisme. Le caryotype comporte donc $46 + 1 = 47$ chromosomes.
\end{definition}

L'examen du caryotype d'une personne atteinte de trisomie 21 montre 22 paires d'autosomes normales, sauf la paire 21 qui comporte trois éléments (on parle de \cle{trisomie libre}), et une paire de gonosomes normale (XX ou XY selon le sexe).

\begin{popfigure}
\begin{tikzpicture}[scale=0.92, every node/.style={font=\scriptsize}]
% ===== Caryotype normal 46,XY =====
\node[font=\small\bfseries] at (-4.2,3.6) {Caryotype normal : 46,XY};
% Groupes de paires (représentation schématique par rangées)
% Rangée 1 : paires 1-6 (grands)
\foreach \i/\x in {1/-6.6, 2/-5.9, 3/-5.2, 4/-4.5, 5/-3.8, 6/-3.1}{
  \draw[thick, popblue] (\x,2.9) -- (\x,2.5);
  \draw[thick, popblue] (\x+0.18,2.9) -- (\x+0.18,2.5);
}
\node at (-4.8,3.15) {Paires 1 à 6};
% Rangée 2 : paires 7-12 (moyens)
\foreach \i/\x in {1/-6.6, 2/-5.9, 3/-5.2, 4/-4.5, 5/-3.8, 6/-3.1}{
  \draw[thick, popblue] (\x,2.1) -- (\x,1.75);
  \draw[thick, popblue] (\x+0.18,2.1) -- (\x+0.18,1.75);
}
\node at (-4.85,2.35) {Paires 7 à 12};
% Rangée 3 : paires 13-18 (petits)
\foreach \i/\x in {1/-6.6, 2/-5.9, 3/-5.2, 4/-4.5, 5/-3.8}{
  \draw[thick, popblue] (\x,1.35) -- (\x,1.05);
  \draw[thick, popblue] (\x+0.18,1.35) -- (\x+0.18,1.05);
}
\node at (-4.9,1.6) {Paires 13 à 18};
% Rangée 4 : paires 19-21 (très petits acrocentriques)
\foreach \i/\x in {1/-6.6, 2/-5.9, 3/-5.2}{
  \draw[thick, popblue] (\x,0.6) -- (\x,0.35);
  \draw[thick, popblue] (\x+0.18,0.6) -- (\x+0.18,0.35);
}
\node at (-5.35,0.85) {Paires 19 à 21};
% Gonosomes XY
\draw[thick, popgreen] (-4.3,0.6) -- (-4.3,0.3); % X
\draw[thick, poporange] (-4.0,0.55) -- (-4.0,0.35); % Y
\node at (-4.15,0.85) {X \quad Y};
\node[font=\scriptsize\bfseries, popblue] at (-4.2,-0.2) {Total : 46 chromosomes};

% ===== Caryotype trisomique 47,XY,+21 =====
\node[font=\small\bfseries] at (3.4,3.6) {Caryotype trisomique : 47,XY,+21};
% Rangée 1
\foreach \i/\x in {1/0.9, 2/1.6, 3/2.3, 4/3.0, 5/3.7, 6/4.4}{
  \draw[thick, popblue] (\x,2.9) -- (\x,2.5);
  \draw[thick, popblue] (\x+0.18,2.9) -- (\x+0.18,2.5);
}
\node at (2.75,3.15) {Paires 1 à 6};
% Rangée 2
\foreach \i/\x in {1/0.9, 2/1.6, 3/2.3, 4/3.0, 5/3.7, 6/4.4}{
  \draw[thick, popblue] (\x,2.1) -- (\x,1.75);
  \draw[thick, popblue] (\x+0.18,2.1) -- (\x+0.18,1.75);
}
\node at (2.7,2.35) {Paires 7 à 12};
% Rangée 3
\foreach \i/\x in {1/0.9, 2/1.6, 3/2.3, 4/3.0, 5/3.7}{
  \draw[thick, popblue] (\x,1.35) -- (\x,1.05);
  \draw[thick, popblue] (\x+0.18,1.35) -- (\x+0.18,1.05);
}
\node at (2.65,1.6) {Paires 13 à 18};
% Rangée 4 : 19, 20, et 21 (trois)
\foreach \i/\x in {1/0.9, 2/1.6}{
  \draw[thick, popblue] (\x,0.6) -- (\x,0.35);
  \draw[thick, popblue] (\x+0.18,0.6) -- (\x+0.18,0.35);
}
\node at (1.15,0.85) {19 et 20};
% Trois chromosomes 21 entourés
\draw[thick, poppink] (2.45,0.6) -- (2.45,0.35);
\draw[thick, poppink] (2.63,0.6) -- (2.63,0.35);
\draw[thick, poppink] (2.81,0.6) -- (2.81,0.35);
\draw[thick, poppink, rounded corners] (2.32,0.22) rectangle (2.94,0.72);
\node[font=\scriptsize\bfseries, poppink] at (2.63,0.0) {3 chr 21};
% Gonosomes XY
\draw[thick, popgreen] (3.7,0.6) -- (3.7,0.3);
\draw[thick, poporange] (4.0,0.55) -- (4.0,0.35);
\node at (3.85,0.85) {X \quad Y};
\node[font=\scriptsize\bfseries, poppink] at (3.4,-0.2) {Total : 47 chromosomes};
\end{tikzpicture}
\legende{Comparaison schématique de deux caryotypes masculins. À gauche, caryotype normal : 46 chromosomes, dont deux chromosomes 21 (paire 21) et les gonosomes XY. À droite, caryotype d'un garçon atteint de trisomie 21 : 47 chromosomes, avec \textbf{trois chromosomes 21} (entourés en rose) ; la formule chromosomique s'écrit 47,XY,+21.}
\end{popfigure}

\textbf{Conséquences : modification des caractères.} La présence d'un chromosome 21 surnuméraire perturbe le développement de l'organisme dès la vie embryonnaire, car les gènes portés par ce chromosome sont présents en trois exemplaires au lieu de deux (effet de \cle{dosage génique}). Il en résulte un \cle{phénotype} caractéristique, c'est-à-dire un ensemble de caractères observables modifiés :

\begin{itemize}
\item un visage arrondi, des yeux bridés (plis épicanthiques), une petite taille ;
\item une hypotonie musculaire (flaccidité) ;
\item une \cle{déficience intellectuelle variable} d'un individu à l'autre, allant de légère à modérée ;
\item parfois des malformations associées (cardiaques notamment, comme la communication interventriculaire), qui justifient un suivi médical régulier dès la naissance.
\end{itemize}

Avec un accompagnement adapté (éducation spécialisée, suivi médical, soutien familial, insertion professionnelle), les personnes trisomiques peuvent apprendre, travailler et s'intégrer pleinement dans la vie sociale.

\courssub{La formule chromosomique : écrire et lire un caryotype}

Pour décrire rapidement et sans ambiguïté le caryotype d'un individu, les biologistes utilisent une écriture codifiée internationale (norme ISCN) : la \cle{formule chromosomique}.

\begin{definition}[Formule chromosomique]
La \cle{formule chromosomique} est l'écriture conventionnelle d'un caryotype. Elle comporte trois éléments, séparés par des virgules :
\begin{enumerate}
\item le \textbf{nombre total de chromosomes} de la cellule somatique ;
\item le \textbf{type de gonosomes} (XX ou XY) ;
\item l'\textbf{anomalie éventuelle}, indiquée par le signe $+$ (chromosome surnuméraire) ou $-$ (chromosome manquant) suivi du numéro du chromosome concerné (ou de la description de l'anomalie structurelle).
\end{enumerate}
\end{definition}

\begin{methode}[Écrire et lire une formule chromosomique]
\begin{enumerate}
\item \textbf{Compter} tous les chromosomes du caryotype et écrire ce nombre en premier.
\item \textbf{Identifier les gonosomes} : écrire XX (fille) ou XY (garçon) après une virgule.
\item \textbf{Repérer l'anomalie} : chercher une paire comportant trois chromosomes (noter $+$ suivi du numéro, ex: $+21$) ou un chromosome seul là où il devrait y en avoir deux (noter $-$ suivi du numéro, ex: $-21$ ou $-X$). Si tout est normal, on n'écrit rien de plus.
\item \textbf{Vérifier la cohérence} : le nombre total doit correspondre à $46$ plus le nombre de chromosomes surnuméraires, moins le nombre de chromosomes manquants.
\end{enumerate}
\textbf{Lecture de la formule 47,XX,+21} : 47 chromosomes au total ; individu de sexe féminin (XX) ; un chromosome 21 surnuméraire, donc trois chromosomes 21 : il s'agit d'une fille atteinte de trisomie 21.
\end{methode}

\begin{exemplebox}[Formules chromosomiques de la trisomie 21]
\begin{itemize}
\item \textbf{47,XX,+21} : fille atteinte de trisomie 21 (47 chromosomes, gonosomes XX, un chromosome 21 en plus).
\item \textbf{47,XY,+21} : garçon atteint de trisomie 21 (47 chromosomes, gonosomes XY, un chromosome 21 en plus).
\item \textbf{46,XX} et \textbf{46,XY} : caryotypes normaux d'une fille et d'un garçon (aucune anomalie à signaler, la formule s'arrête aux gonosomes).
\end{itemize}
\end{exemplebox}

\begin{exoresolu}[Identifier une anomalie à partir d'un caryotype]
\textbf{Énoncé.} Le caryotype des cellules d'un nouveau-né montre 47 chromosomes. L'examen attentif révèle trois chromosomes 21 et les gonosomes sont de type XY.
\begin{enumerate}
\item Identifier l'anomalie chromosomique présentée par cet enfant.
\item Rédiger la formule chromosomique de ce caryotype.
\item Expliquer en quelques lignes l'origine de cette anomalie.
\end{enumerate}
\tcblower
\textbf{Solution détaillée.}
\begin{enumerate}
\item Le caryotype comporte 47 chromosomes au lieu de 46, avec \textbf{trois chromosomes 21} au lieu de deux : il s'agit d'une \cle{trisomie 21} (trisomie libre).
\item Application de la méthode : nombre total $= 47$ ; gonosomes $=$ XY (garçon) ; anomalie $=$ un chromosome 21 surnuméraire, notée $+21$. La formule chromosomique est donc : \textbf{47,XY,+21}. Vérification : $46 + 1 = 47$, la formule est cohérente.
\item Cette anomalie provient d'une \cle{non-disjonction} de la paire de chromosomes 21 au cours de la méiose I chez l'un des parents : un gamète anormal à 24 chromosomes (portant deux chromosomes 21) a fusionné avec un gamète normal à 23 chromosomes, donnant un zygote à $24 + 23 = 47$ chromosomes, dont trois chromosomes 21.
\end{enumerate}
\end{exoresolu}

\courssub{Autres exemples d'anomalies du nombre de chromosomes (gonosomes)}

La non-disjonction peut toucher d'autres chromosomes que le chromosome 21, notamment les gonosomes. Deux exemples sont au programme et doivent être connus.

\begin{exemplebox}[Monosomie X et trisomie XXY]
\begin{itemize}
\item \textbf{Monosomie X ou syndrome de Turner : formule 45,X.} L'individu ne possède que 45 chromosomes, avec un \textbf{seul chromosome X} et aucun second gonosome (ni X, ni Y). Il s'agit d'une fille (phénotype féminin) présentant notamment une petite taille, un cou ptérygoïde, et une absence de développement des ovaires (dysgénésie gonadique entraînant stérilité et absence de puberté spontanée). C'est l'unique monosomie d'un chromosome entier compatible avec la vie.
\item \textbf{Trisomie XXY ou syndrome de Klinefelter : formule 47,XXY.} L'individu possède 47 chromosomes, avec \textbf{deux chromosomes X et un chromosome Y}. Il s'agit d'un garçon (la présence de Y détermine le sexe masculin) présentant souvent une taille élevée, un développement insuffisant des testicules (hypogonadisme), une gynécomastie et une stérilité (azoospermie).
\end{itemize}
\end{exemplebox}

\begin{attention}[Bien distinguer les formules]
Ne confonds pas le \textbf{nombre total} de chromosomes et l'\textbf{anomalie} : dans la formule 45,X, le nombre 45 signifie qu'il manque un chromosome (un gonosome), tandis que dans 47,XXY, le nombre 47 signifie qu'il y a un chromosome surnuméraire (ici un gonosome X). Vérifie toujours que le total correspond à $46 \pm$ le nombre de chromosomes anormaux.
\end{attention}

\courssub{Anomalie chromosomique ou anomalie génique : une distinction essentielle}

Toutes les anomalies héréditaires ne se ressemblent pas. Il faut distinguer deux grandes catégories, car les méthodes de détection et les modes de transmission diffèrent radicalement.

\begin{definition}[Anomalie chromosomique et anomalie génique]
\begin{itemize}
\item Une \cle{anomalie chromosomique} est une modification du \textbf{nombre} de chromosomes (ex: trisomie 21, monosomie X) ou de la \textbf{structure} d'un chromosome (délétion, translocation, inversion, duplication d'un fragment). Elle est \textbf{visible sur le caryotype} (au microscope optique après bandage).
\item Une \cle{anomalie génique} est la modification (\cle{mutation}) d'un \textbf{gène}, c'est-à-dire d'une petite portion de la molécule d'ADN (changement d'une base, insertion/délétion de quelques bases). Le nombre et l'aspect macroscopique des chromosomes restent normaux : elle est \textbf{invisible sur le caryotype standard}.
\end{itemize}
\end{definition}

\begin{attention}[Piège classique au Probatoire]
L'\cle{albinisme} et la \cle{drépanocytose} sont des \textbf{anomalies géniques} (maladies mongéniques autosomales récessives), dues à la mutation d'un gène situé sur un autosome (respectivement gène \emph{TYR} sur chr 11 et gène \emph{HBB} sur chr 11). Le caryotype d'une personne albinos ou drépanocytaire est \textbf{parfaitement normal} : 46,XX ou 46,XY, sans aucune anomalie visible. Ne dis jamais qu'une personne albinos « a un chromosome en trop » ou « un chromosome modifié » au sens chromosomique : c'est faux. Seules les anomalies de nombre ou de structure des chromosomes se détectent sur un caryotype ; les anomalies géniques demandent d'autres analyses (dosage enzymatique, électrophorèse d'hémoglobine, séquençage de l'ADN).
\end{attention}

\begin{center}
\begin{tabularx}{\linewidth}{|l|X|X|}
\hline
\rowcolor{popblueL} \textbf{Critère} & \textbf{Anomalie chromosomique} & \textbf{Anomalie génique} \\
\hline
Nature & Nombre ou structure d'un chromosome modifié & Mutation d'un gène (ADN) \\
\hline
Visible sur le caryotype standard ? & \textbf{Oui} & \textbf{Non} \\
\hline
Exemples au programme & Trisomie 21 ; Turner (45,X) ; Klinefelter (47,XXY) & Albinisme ; drépanocytose \\
\hline
Transmission & Généralement \textbf{non héréditaire} (accident méiotique \emph{de novo}), sauf translocations équilibrées parentales & \textbf{Héréditaire} (lois de Mendel : autosomale récessive pour albinisme et drépanocytose) \\
\hline
\end{tabularx}
\end{center}

\courssub{Éradication des préjugés : ce que la science nous apprend}

\begin{aretenir}[La trisomie 21 est un accident méiotique, pas une faute]
La trisomie 21 résulte d'une \textbf{non-disjonction accidentelle} des chromosomes 21 lors de la méiose, événement aléatoire qui peut survenir chez tout couple, quelle que soit son origine, sa religion, son niveau de vie ou son mode de vie. Ce n'est ni une \textbf{malédiction}, ni une \textbf{punition divine}, ni la \textbf{faute des parents} (ni de la mère, ni du père). Accuser une mère ou une famille d'être responsable de la naissance d'un enfant trisomique est scientifiquement faux et socialement injuste. Les personnes atteintes de trisomie 21, comme celles atteintes d'albinisme ou de drépanocytose, ont droit au respect, à l'éducation, aux soins et à l'affection : c'est un devoir de solidarité pour toute la communauté.
\end{aretenir}

\begin{savaistu}[Trisomie 21 et âge de la mère]
Les statistiques médicales montrent que la fréquence de la trisomie 21 augmente avec l'\textbf{âge de la mère} : elle est d'environ 1 naissance sur \num{1500} chez les mères de 20 ans, mais peut atteindre environ 1 sur 100 après 40 ans. Cela s'explique par le vieillissement des ovocytes, bloqués en prophase I de la méiose depuis la vie fœtale, ce qui augmente le risque de non-disjonction lors de la reprise méiotique à l'ovulation. C'est pourquoi le \textbf{suivi prénatal} (consultations régulières, échographies, dosage des marqueurs sériques, parfois analyse du caryotype fœtal par amniocentèse ou trophobiopsie) est particulièrement recommandé aux femmes enceintes après 35 ans. Rappelons toutefois qu'une non-disjonction peut survenir à tout âge (y compris chez le père, plus rarement) : aucun parent n'est à l'abri, et aucun n'est responsable.
\end{savaistu}

\begin{exercice}[À toi de jouer]
\begin{enumerate}
\item Le caryotype d'une fillette montre 45 chromosomes, dont un seul chromosome X. Rédige sa formule chromosomique et nomme l'anomalie.
\item Un garçon présente la formule chromosomique 47,XXY. Combien de chromosomes possèdent ses cellules somatiques ? Nomme l'anomalie et donne le phénotype sexuel.
\item Vrai ou faux (justifie) : « Le caryotype d'une personne atteinte de drépanocytose montre 47 chromosomes. »
\item Un couple a un enfant trisomique 21. Le père dit : « C'est de ta faute, tes ovules étaient vieux. » En t'appuyant sur les connaissances scientifiques, explique pourquoi cette affirmation est injuste et fausse.
\end{enumerate}
\end{exercice}

\begin{corrige}[Exercice : À toi de jouer]
\begin{enumerate}
\item Nombre total $= 45$ ; un seul gonosome X ; il manque un chromosome sexuel. Formule : \textbf{45,X}. Il s'agit de la \textbf{monosomie X} (syndrome de Turner).
\item Ses cellules somatiques possèdent \textbf{47 chromosomes} (deux X et un Y). Il s'agit de la \textbf{trisomie XXY} (syndrome de Klinefelter). Le phénotype sexuel est \textbf{masculin} (présence du chromosome Y).
\item \textbf{Faux.} La drépanocytose est une \textbf{anomalie génique} (mutation ponctuelle du gène $\beta$-globine sur le chromosome 11), invisible sur le caryotype. Le caryotype d'une personne drépanocytaire est normal : 46,XX ou 46,XY.
\item Cette affirmation est fausse car : (1) La non-disjonction est un \textbf{accident aléatoire} de la méiose, pas une « usure » imputable à la mère ; elle peut survenir lors de la spermatogenèse (père) ou de l'ovogenèse (mère). (2) Même si le risque statistique augmente avec l'âge maternel, la très grande majorité des femmes de plus de 35 ans ont des enfants non trisomiques, et des trisomies naissent de mères jeunes. (3) Attribuer une « faute » à un accident biologique aléatoire est une confusion entre causalité biologique et responsabilité morale, source de stigmatisation injustifiée.
\end{enumerate}
\end{corrige}

Tu sais désormais reconnaître une anomalie du nombre de chromosomes, rédiger et lire une formule chromosomique, et distinguer anomalie chromosomique et anomalie génique. Dans la section suivante, nous verrons comment les chromosomes, et en particulier les gonosomes X et Y, déterminent le sexe de l'enfant dès la fécondation.

\coursec{Chromosomes et détermination du sexe}

Dans la section précédente, nous avons vu que le caryotype humain comporte 46 chromosomes répartis en 23 paires, et que toute anomalie du nombre de chromosomes (comme la trisomie 21) entraîne une modification des caractères. Nous avons aussi remarqué que la 23\ieme{} paire, celle des \cle{gonosomes} ou chromosomes sexuels, présente une particularité remarquable : elle n'est pas toujours constituée de deux chromosomes identiques. C'est précisément cette particularité qui va nous permettre de répondre à une question que se posent toutes les familles : \emph{pourquoi naît-on fille ou garçon ?}

\begin{attention}[Un préjugé à démanteler]
Dans certaines familles camerounaises, on accuse encore la femme de «~ne donner que des filles~» ou de «~ne pas donner de garçon~». Ce préjugé, source de divorces et de souffrances, repose sur une ignorance du mécanisme biologique de la détermination du sexe. Cette section va te montrer, preuves à l'appui, que c'est le \textbf{spermatozoïde du père} qui détermine le sexe de l'enfant.
\end{attention}

\courssub{Observation : la paire des gonosomes diffère selon le sexe}

\begin{experience}[Comparaison des caryotypes d'un homme et d'une femme]
\textbf{Matériel :} deux planches de caryotypes humains normaux, établis à partir de lymphocytes en mitose (technique vue en section 1).

\textbf{Protocole :} on observe et on compare les 23 paires de chromosomes des deux caryotypes, en particulier la dernière paire.

\textbf{Observations :}
\begin{itemize}
\item Les 22 premières paires (\cle{autosomes}) sont identiques chez l'homme et chez la femme ;
\item chez la femme, la 23\ieme{} paire est formée de \textbf{deux chromosomes identiques}, de grande taille, en forme de X : on la note \textbf{XX} ;
\item chez l'homme, la 23\ieme{} paire est formée de \textbf{deux chromosomes différents} : un grand chromosome X et un chromosome beaucoup plus petit, en forme de V ou de petit bâtonnet, appelé chromosome \textbf{Y} : on la note \textbf{XY}.
\end{itemize}

\textbf{Interprétation :} la seule différence chromosomique entre un homme et une femme porte sur les gonosomes. On émet alors l'hypothèse que ce sont eux qui déterminent le sexe.
\end{experience}

\begin{definition}[Gonosomes et formules chromosomiques des sexes]
Les \cle{gonosomes} (ou chromosomes sexuels) sont les chromosomes de la 23\ieme{} paire, responsables de la détermination du sexe.
\begin{itemize}
\item La \textbf{femme} possède deux chromosomes X : sa formule chromosomique est \textbf{44 autosomes + XX} (soit 2n = 46, XX) ;
\item l'\textbf{homme} possède un chromosome X et un chromosome Y : sa formule chromosomique est \textbf{44 autosomes + XY} (soit 2n = 46, XY).
\end{itemize}
Le chromosome \textbf{Y} est nettement \textbf{plus petit} que le chromosome X et porte beaucoup moins de gènes. Il contient notamment le gène \textbf{SRY} (\emph{Sex-determining Region Y}) qui déclenche le développement des testicules chez l'embryon.
\end{definition}

\begin{popfigure}
\begin{center}
\begin{tikzpicture}[scale=1]
% Chromosome X (grand, en X)
\begin{scope}[shift={(0,0)}]
\draw[line width=4.5pt, popblue, line cap=round] (-0.55,-1.5) -- (0.55,1.5);
\draw[line width=4.5pt, popblue, line cap=round] (0.55,-1.5) -- (-0.55,1.5);
\draw[line width=2.2pt, popblueL, line cap=round] (-0.45,-1.35) -- (0.45,1.35);
\draw[line width=2.2pt, popblueL, line cap=round] (0.45,-1.35) -- (-0.45,1.35);
\node[popdark, font=\large\bfseries] at (0,-2.1) {Chromosome X};
\node[popmuted, font=\small] at (0,-2.6) {grand, en forme de X} ;
\end{scope}
% Chromosome Y (petit)
\begin{scope}[shift={(4.5,0)}]
\draw[line width=4.5pt, poporange, line cap=round] (0,0.55) -- (-0.3,-0.35);
\draw[line width=4.5pt, poporange, line cap=round] (0,0.55) -- (0.3,-0.35);
\draw[line width=4.5pt, poporange, line cap=round] (0,0.55) -- (0,1.15);
\draw[line width=2.2pt, poporangeL, line cap=round] (0,0.45) -- (-0.22,-0.28);
\draw[line width=2.2pt, poporangeL, line cap=round] (0,0.45) -- (0.22,-0.28);
\draw[line width=2.2pt, poporangeL, line cap=round] (0,0.5) -- (0,1.05);
\node[popdark, font=\large\bfseries] at (0,-2.1) {Chromosome Y};
\node[popmuted, font=\small] at (0,-2.6) {petit, en forme de V};
\end{scope}
% Echelle comparative
\draw[<->, popmuted] (2.2,1.5) -- (2.2,-1.5) node[midway, right, font=\scriptsize, popmuted] {taille du X};
\draw[<->, popmuted] (5.6,1.15) -- (5.6,-0.35) node[midway, right, font=\scriptsize, popmuted] {taille du Y};
\end{tikzpicture}
\end{center}
\legende{Comparaison schématique des gonosomes humains : le chromosome X (bleu) est grand et en forme de X ; le chromosome Y (orange) est beaucoup plus petit et en forme de V. La femme possède deux X ; l'homme possède un X et un Y.}
\end{popfigure}

\courssub{Les gamètes produits par chaque sexe}

Rappelons que les gamètes (ovules et spermatozoïdes) sont produits par \cle{méiose} : chaque gamète ne reçoit qu'\textbf{un chromosome de chaque paire} (n = 23 chromosomes).

Raisonnons sur la paire de gonosomes :

\begin{itemize}
\item Chez la \textbf{femme} (XX) : les deux gonosomes sont identiques. Tous les ovules reçoivent donc un chromosome X. La femme ne produit qu'\textbf{un seul type d'ovule} : \textbf{22 autosomes + X}. On dit que la femme est \cle{homogamétique}.
\item Chez l'\textbf{homme} (XY) : les deux gonosomes sont différents. La moitié des spermatozoïdes reçoit le chromosome X, l'autre moitié reçoit le chromosome Y. L'homme produit donc \textbf{deux types de spermatozoïdes} en proportions égales : \textbf{22 autosomes + X} (50~\%) et \textbf{22 autosomes + Y} (50~\%). On dit que l'homme est \cle{hétérogamétique}.
\end{itemize}

\begin{aretenir}[Gamètes et sexes]
\begin{itemize}
\item Ovule (unique type) : 22 autosomes + X ;
\item Spermatozoïdes (deux types équiprobables) : 22 autosomes + X \quad ou \quad 22 autosomes + Y.
\end{itemize}
\end{aretenir}

\courssub{Démonstration guidée : l'échiquier de croisement de la fécondation}

Construisons pas à pas l'échiquier qui croise les gamètes possibles des deux parents.

\textbf{Étape 1.} On place en ligne les gonosomes apportés par les gamètes du père : X et Y (spermatozoïdes).

\textbf{Étape 2.} On place en colonne le gonosome apporté par le gamète unique de la mère : X (ovule).

\textbf{Étape 3.} On remplit chaque case en réunissant les gonosomes apportés par les deux gamètes : c'est la fécondation, qui rétablit la diploïdie (2n = 46).

\textbf{Étape 4.} On lit le sexe correspondant à chaque case et on compte les proportions.

\begin{popfigure}
\begin{center}
\begin{tikzpicture}[scale=1.15]
% Grille
\draw[popline, thick] (0,0) grid (3,2);
% En-têtes
\node[font=\small\bfseries, popmuted] at (-1.15,1.5) {Ovules de la mère};
\node[font=\small\bfseries, popmuted, rotate=0] at (1.5,2.55) {Spermatozoïdes du père};
% Gametes mère
\node[font=\Large\bfseries, popblue] at (0.5,1.5) {X};
% Gametes père
\node[font=\Large\bfseries, poporange] at (1.5,0.5) {X};
\node[font=\Large\bfseries, poporange] at (2.5,0.5) {Y};
% Cases
\node[font=\Large\bfseries, popblue] at (1.5,1.5) {XX};
\node[font=\Large\bfseries, poporange] at (2.5,1.5) {XY};
% Labels sexe
\node[font=\small, popblue] at (1.5,1.12) {fille};
\node[font=\small, poporange] at (2.5,1.12) {garçon};
% Proportions
\node[font=\small\bfseries, popblue] at (1.5,-0.45) {50~\%};
\node[font=\small\bfseries, poporange] at (2.5,-0.45) {50~\%};
\end{tikzpicture}
\end{center}
\legende{Échiquier de croisement de la fécondation : l'ovule apporte toujours un X ; le spermatozoïde apporte X ou Y avec la même probabilité. Il en résulte 1/2 de zygotes XX (filles) et 1/2 de zygotes XY (garçons).}
\end{popfigure}

\textbf{Lecture des résultats :} sur les deux cases possibles (en tenant compte des proportions égales des spermatozoïdes), on obtient :
\begin{itemize}
\item X (mère) + X (père) = \textbf{XX} : une fille ;
\item X (mère) + Y (père) = \textbf{XY} : un garçon.
\end{itemize}

\begin{propriete}[Probabilité des sexes à la naissance]
À chaque fécondation, la probabilité d'avoir une fille est égale à la probabilité d'avoir un garçon :
\[ P(\text{fille}) = P(\text{garçon}) = \frac{1}{2} \]
Cette propriété se démontre par l'échiquier de croisement ci-dessus (démonstration exigible au Probatoire). Le sexe de l'enfant dépend \textbf{uniquement du gonosome apporté par le spermatozoïde}, donc du \textbf{père}.
\end{propriete}

\begin{exemplebox}[Application chiffrée]
Dans une maternité enregistrant 1000 naissances, on attend statistiquement environ :
\[ 1000 \times \frac{1}{2} = 500 \text{ filles} \quad \text{et} \quad 500 \text{ garçons.} \]
Les écarts observés dans une famille donnée (par exemple 4 filles de suite) relèvent du hasard des tirages indépendants : chaque naissance «~rejoue~» la probabilité 1/2, sans mémoire des naissances précédentes.
\end{exemplebox}

\begin{attention}[Erreurs fréquentes et préjugés]
\begin{itemize}
\item Le sexe n'est déterminé \textbf{ni par la mère} (elle ne fournit que des X), \textbf{ni par le hasard mystique, ni par la volonté des ancêtres} : il dépend du gonosome (X ou Y) apporté par le \textbf{spermatozoïde} au moment de la fécondation.
\item Ne dis pas «~la probabilité devient plus forte d'avoir un garçon après trois filles~» : les naissances sont des événements indépendants, la probabilité reste 1/2 à chaque fois.
\item Ne confonds pas gonosomes (23\ieme{} paire) et autosomes (22 premières paires).
\end{itemize}
\end{attention}

\begin{methode}[Déterminer le sexe d'un individu à partir de son caryotype]
\begin{enumerate}
\item Repérer la \textbf{dernière paire} du caryotype (paire 23, les gonosomes).
\item Observer les deux chromosomes de cette paire :
\begin{itemize}
\item deux chromosomes \textbf{identiques et grands} (XX) $\Rightarrow$ l'individu est une \textbf{femme} ;
\item deux chromosomes \textbf{différents}, l'un grand (X) et l'autre petit (Y) $\Rightarrow$ l'individu est un \textbf{homme}.
\end{itemize}
\item Conclure en écrivant la formule chromosomique complète : 44 autosomes + XX (femme) ou 44 autosomes + XY (homme).
\end{enumerate}
\end{methode}

\begin{exoresolu}[Sexe d'un individu à partir du caryotype de ses cellules]
\textbf{Énoncé.} On réalise le caryotype des cellules reproductrices (cellules de la lignée germinale, avant méiose) d'un individu A. On y observe 46 chromosomes ; la dernière paire est formée d'un grand chromosome et d'un chromosome nettement plus petit.

1. Quel est le sexe de l'individu A ? Justifier.
2. Écrire sa formule chromosomique.
3. Quels types de gamètes produira-t-il et dans quelles proportions ?
\tcblower
\textbf{Solution détaillée.}

1. La dernière paire (gonosomes) est constituée de deux chromosomes \textbf{différents} : un grand chromosome X et un petit chromosome Y. La paire est donc \textbf{XY} : l'individu A est un \textbf{homme}.

2. Formule chromosomique : \textbf{44 autosomes + XY}, soit 2n = 46, XY.

3. Par méiose, chaque spermatozoïde ne reçoit qu'un seul gonosome. L'individu A produira donc \textbf{deux types de spermatozoïdes en proportions égales} : 50~\% de spermatozoïdes 22 autosomes + X et 50~\% de spermatozoïdes 22 autosomes + Y. C'est le type de spermatozoïde qui féconde l'ovule (toujours 22 autosomes + X) qui déterminera le sexe de l'enfant.
\end{exoresolu}

\begin{savaistu}[Et chez les autres espèces ?]
Le système XX/XY n'est pas universel ! Chez les \textbf{oiseaux} et les \textbf{papillons}, c'est la \textbf{femelle} qui est hétérogamétique : elle possède deux gonosomes différents notés \textbf{ZW}, tandis que le mâle est ZZ. C'est donc l'ovule (Z ou W) qui détermine le sexe de l'oisillon. Chez les \textbf{abeilles}, c'est encore autre chose : les femelles (reines et ouvrières) sont diploïdes et les mâles (faux bourdons) sont haploïdes, issus d'œufs non fécondés (parthénogenèse arrhénotoque). Ce complément est utile au Probatoire pour montrer que la détermination chromosomique du sexe admet plusieurs modalités dans le monde vivant.
\end{savaistu}

\begin{aretenir}[L'essentiel de la section]
\begin{itemize}
\item La femme est \textbf{XX} (homogamétique) : un seul type d'ovule, 22 autosomes + X.
\item L'homme est \textbf{XY} (hétérogamétique) : deux types de spermatozoïdes, 22 autosomes + X et 22 autosomes + Y, en proportions égales.
\item L'échiquier de croisement montre que $P(\text{fille}) = P(\text{garçon}) = \dfrac{1}{2}$ à chaque naissance.
\item Le sexe de l'enfant dépend du \textbf{spermatozoïde}, donc du père : rendre la femme responsable du sexe de l'enfant est scientifiquement injuste.
\item Sur un caryotype, la paire 23 révèle le sexe : XX = femme, XY = homme.
\end{itemize}
\end{aretenir}

\coursec{Hérédité liée au sexe : daltonisme et hémophilie}

\courssub{Transition depuis la détermination du sexe}
Dans la section précédente, nous avons vu que le sexe d'un individu est déterminé par la paire de chromosomes sexuels : XX chez la femme, XY chez l'homme. Cette différence fondamentale a une conséquence directe sur la transmission des gènes portés par le chromosome X. En effet, le chromosome Y ne possède pas d'allèle correspondant pour la plupart des gènes situés sur le X (région non homologue). C'est ce qui crée les particularités de l'hérédité liée au sexe, que nous allons explorer à travers deux maladies humaines classiques : le daltonisme et l'hémophilie.

\courssub{Définition et bases génétiques de l'hérédité liée au sexe}
\begin{definition}[Hérédité liée au sexe (hétérochromosomique)]
On appelle \textbf{hérédité liée au sexe} (ou \textbf{hétérochromosomique}) la transmission d'un caractère dont le gène responsable est situé sur un chromosome sexuel, le plus souvent le chromosome X. Comme le chromosome Y ne porte pas d'allèle homologue pour ce gène (absence de région homologue), les mâles (XY) n'ont qu'un seul allèle pour ce locus (hémizygotes), alors que les femelles (XX) en ont deux.
\end{definition}

\begin{definition}[Femme conductrice (hétérozygote porteuse saine)]
Une \textbf{femme conductrice} est une femme hétérozygote pour un allèle récessif porté par le chromosome X (génotype $X^D X^d$ ou $X^H X^h$). Elle ne manifeste pas la maladie (phénotype sain) mais peut transmettre l'allèle muté à sa descendance.
\end{definition}

\emph{Observation concrète :} Dans de nombreuses familles camerounaises, on observe que le daltonisme « saute » une génération : un grand-père daltonien a des petits-fils daltoniens, alors que ses fils sont tous sains. Cette régularité trouve son explication dans la génétique du chromosome X.

\courssub{Le daltonisme : anomalie de la vision des couleurs}
Le daltonisme le plus fréquent (protanopie/déutéranopie) est dû à un allèle récessif $d$ porté par le chromosome X. L'allèle normal est noté $D$. Les génotypes et phénotypes possibles sont :

\begin{center}
\begin{tabularx}{\linewidth}{|l|X|X|}\hline
\textbf{Génotype} & \textbf{Sexe} & \textbf{Phénotype} \\ \hline
$X^D X^D$ & Femme & Vision normale (non conductrice) \\
$X^D X^d$ & Femme & Vision normale (conductrice) \\
$X^d X^d$ & Femme & Daltonienne (rare) \\
$X^D Y$ & Homme & Vision normale \\
$X^d Y$ & Homme & Daltonien \\ \hline
\end{tabularx}
\end{center}

\begin{propriete}[Transmission du daltonisme — démontrée par échiquier]
Un garçon daltonien ($X^d Y$) a obligatoirement hérité de l'allèle $d$ de sa mère. Un père daltonien transmet son chromosome $X^d$ à toutes ses filles (qui deviennent conductrices) et jamais à ses fils (qui reçoivent son chromosome Y).
\end{propriete}

\begin{experience}[Échiquier de croisement : femme conductrice $\times$ homme sain]
\textbf{Matériel :} Papier, crayon, connaissance des allèles $D$ et $d$. \\
\textbf{Protocole :}
\begin{enumerate}
\item Écrire les gamètes possibles de la mère ($X^D$, $X^d$) en ligne (ou en colonne).
\item Écrire les gamètes possibles du père ($X^D$, $Y$) en colonne (ou en ligne).
\item Remplir chaque case par le génotype du zygote obtenu par fusion des gamètes.
\item Déduire le phénotype et le sexe de chaque descendance possible.
\end{enumerate}
\textbf{Observations (échiquier de croisement $X^D X^d \times X^D Y$) :}
\begin{popfigure}
\begin{tikzpicture}[scale=1.1, every node/.style={font=\small}]
% Grid 2x2
\draw (0,0) grid (2,2);
% Row labels (mother gametes) on left
\node[left=0.2cm] at (0,1.5) {$X^D$};
\node[left=0.2cm] at (0,0.5) {$X^d$};
% Column labels (father gametes) on top
\node[above=0.2cm] at (0.5,2) {$X^D$};
\node[above=0.2cm] at (1.5,2) {$Y$};
% Offspring genotypes (center of each cell)
\node at (0.5,1.5) {$X^D X^D$};
\node at (1.5,1.5) {$X^D Y$};
\node at (0.5,0.5) {$X^D X^d$};
\node at (1.5,0.5) {$X^d Y$};
% Phenotypes below each genotype
\node[below=0.15cm, font=\tiny, text width=1.8cm, align=center] at (0.5,1.35) {Fille saine\\non conductrice};
\node[below=0.15cm, font=\tiny, text width=1.8cm, align=center] at (1.5,1.35) {Garçon sain};
\node[below=0.15cm, font=\tiny, text width=1.8cm, align=center] at (0.5,0.35) {Fille conductrice};
\node[below=0.15cm, font=\tiny, text width=1.8cm, align=center] at (1.5,0.35) {Garçon daltonien};
% Labels for axes
\node[below=0.6cm] at (1,0) {Gamètes paternels};
\node[left=0.6cm, rotate=90] at (0,1) {Gamètes maternels};
\end{tikzpicture}
\legende{Échiquier de croisement $X^D X^d \times X^D Y$ : 50\% des fils sont daltoniens ($X^d Y$), 50\% des filles sont conductrices ($X^D X^d$). Aucune fille n'est malade car l'allèle $d$ est récessif.}
\end{popfigure}
\textbf{Interprétation :} Parmi les descendants, 50\% des garçons sont daltoniens ($X^d Y$) et 50\% des filles sont conductrices ($X^D X^d$). Aucune fille n'est malade car l'allèle $d$ est récessif.
\textbf{Conclusion :} La transmission du daltonisme suit strictement les règles de l'hérédité liée au X.
\end{experience}

\courssub{L'hémophilie : trouble de la coagulation}
L'hémophilie A (déficit en facteur VIII) et l'hémophilie B (déficit en facteur IX) sont également dues à des allèles récessifs portés par le chromosome X, notés $h$ (allèle muté) et $H$ (allèle normal). Le tableau des génotypes/phénotypes est analogue à celui du daltonisme :

\begin{center}
\begin{tabularx}{\linewidth}{|l|X|X|}\hline
\textbf{Génotype} & \textbf{Sexe} & \textbf{Phénotype} \\ \hline
$X^H X^H$ & Femme & Normale (non conductrice) \\
$X^H X^h$ & Femme & Normale (conductrice) \\
$X^h X^h$ & Femme & Hémophile (très rare) \\
$X^H Y$ & Homme & Normal \\
$X^h Y$ & Homme & Hémophile \\ \hline
\end{tabularx}
\end{center}

La particularité clinique majeure est que les hommes hémophiles saignent abondamment après le moindre traumatisme, ce qui posait jadis de graves problèmes de santé publique avant l'avènement des concentrés de facteurs de coagulation (disponibles dans les grands hôpitaux camerounais comme l'Hôpital Général de Yaoundé ou de Douala).

\courssub{Transmission et échiquiers de croisement types}
Deux situations classiques sont au programme et doivent être maîtrisées pour le Probatoire :

\begin{enumerate}
\item \textbf{Femme conductrice $\times$ Homme sain} ($X^D X^d \times X^D Y$ ou $X^H X^h \times X^H Y$) : 50\% des fils atteints, 50\% des filles conductrices, 50\% des fils sains, 50\% des filles non conductrices.
\item \textbf{Femme conductrice $\times$ Homme atteint} ($X^D X^d \times X^d Y$ ou $X^H X^h \times X^h Y$) : 50\% des fils atteints, 50\% des fils sains ; 50\% des filles atteintes (homozygotes $X^d X^d$ ou $X^h X^h$), 50\% des filles conductrices.
\end{enumerate}

\begin{exoresolu}[Croisement femme conductrice $\times$ homme daltonien]
\textbf{Énoncé :} Une femme conductrice de daltonisme ($X^D X^d$) épouse un homme daltonien ($X^d Y$). Déterminer les proportions de phénotypes parmi leurs enfants.
\tcblower
\textbf{Solution détaillée :}
\begin{enumerate}
\item \textbf{Gamètes maternels :} $X^D$, $X^d$ (méiose normale).
\item \textbf{Gamètes paternels :} $X^d$, $Y$ (méiose normale).
\item \textbf{Échiquier de croisement (4 combinaisons équiprobables) :}
\begin{center}
\begin{tabular}{|c|c|c|}\hline
 & $X^d$ & $Y$ \\ \hline
$X^D$ & $X^D X^d$ (fille conductrice) & $X^D Y$ (garçon sain) \\ \hline
$X^d$ & $X^d X^d$ (fille daltonienne) & $X^d Y$ (garçon daltonien) \\ \hline
\end{tabular}
\end{center}
\item \textbf{Phénotypes attendus (proportions théoriques) :}
\begin{itemize}
\item Filles : 50\% conductrices ($X^D X^d$), 50\% daltoniennes ($X^d X^d$).
\item Garçons : 50\% sains ($X^D Y$), 50\% daltoniens ($X^d Y$).
\end{itemize}
\end{enumerate}
\end{exoresolu}

\courssub{Interprétation d'arbres généalogiques : méthode et exemple}
\begin{methode}[Interpréter un arbre généalogique d'hérédité liée au sexe]
\begin{enumerate}
\item \textbf{Repérer les symboles :} carré = homme, cercle = femme ; symbole plein (noir) = individu atteint ; symbole vide = sain ; symbole hachuré ou pointillé = conductrice (souvent indiqué explicitement).
\item \textbf{Identifier les hommes atteints :} tout homme noir porte l'allèle muté sur son unique chromosome X (hémizygote).
\item \textbf{Remonter à la mère :} chaque homme atteint a obligatoirement une mère conductrice (ou atteinte). Vérifier que la mère n'est pas atteinte (sinon elle serait homozygote, cas rare).
\item \textbf{Vérifier l'absence de transmission père $\rightarrow$ fils :} un père transmet son chromosome Y à ses fils, jamais son X muté. Si un père atteint a un fils atteint, l'hypothèse d'hérédité liée au X est exclue (ce serait une autosomale dominante ou récessive).
\item \textbf{Conclure sur le mode de transmission :} si toutes les règles ci-dessus sont respectées, le caractère est récessif lié au X.
\end{enumerate}
\end{methode}

\begin{popfigure}
\begin{tikzpicture}[scale=0.8, every node/.style={font=\small}, >=Stealth]
% Generation I
\node[draw, circle, minimum size=0.7cm] (I1) at (0,4) {};
\node[draw, rectangle, minimum size=0.7cm, fill=black] (I2) at (3,4) {};
\node[above=0.2cm] at (1.5,4.5) {Génération I};
\draw (I1) -- (I2);
% Generation II
\node[draw, circle, minimum size=0.7cm, pattern=north east lines] (II1) at (-2,2) {};
\node[draw, rectangle, minimum size=0.7cm] (II2) at (0,2) {};
\node[draw, circle, minimum size=0.7cm] (II3) at (2,2) {};
\node[draw, rectangle, minimum size=0.7cm, fill=black] (II4) at (4,2) {};
\node[above=0.2cm] at (1,2.5) {Génération II};
% Links I -> II
\draw (I1.south) -- ++(0,-0.3) -| (II1.north);
\draw (I1.south) -- ++(0,-0.3) -| (II2.north);
\draw (I2.south) -- ++(0,-0.3) -| (II3.north);
\draw (I2.south) -- ++(0,-0.3) -| (II4.north);
% Generation III
\node[draw, rectangle, minimum size=0.7cm, fill=black] (III1) at (-3,0) {};
\node[draw, circle, minimum size=0.7cm] (III2) at (-1,0) {};
\node[draw, rectangle, minimum size=0.7cm] (III3) at (1,0) {};
\node[draw, circle, minimum size=0.7cm, pattern=north east lines] (III4) at (3,0) {};
\node[draw, rectangle, minimum size=0.7cm, fill=black] (III5) at (5,0) {};
\node[above=0.2cm] at (1,0.5) {Génération III};
% Links II -> III (only for couples that have children shown)
\draw (II1.south) -- ++(0,-0.3) -| (III1.north);
\draw (II1.south) -- ++(0,-0.3) -| (III2.north);
\draw (II4.south) -- ++(0,-0.3) -| (III4.north);
\draw (II4.south) -- ++(0,-0.3) -| (III5.north);
% Legend
\begin{scope}[xshift=-6cm, yshift=-1.5cm]
\node[draw, circle, minimum size=0.5cm] (L1) at (0,0) {};
\node[right=0.2cm] at (L1.east) {Femme saine};
\node[draw, circle, minimum size=0.5cm, pattern=north east lines] (L2) at (0,-0.8) {};
\node[right=0.2cm] at (L2.east) {Femme conductrice};
\node[draw, rectangle, minimum size=0.5cm] (L3) at (0,-1.6) {};
\node[right=0.2cm] at (L3.east) {Homme sain};
\node[draw, rectangle, minimum size=0.5cm, fill=black] (L4) at (0,-2.4) {};
\node[right=0.2cm] at (L4.east) {Homme atteint};
\end{scope}
\end{tikzpicture}
\legende{Arbre généalogique sur trois générations montrant la transmission du daltonisme (ou hémophilie). Génération I : père atteint, mère saine. Génération II : filles conductrices (dont une hachurée), fils sains. Génération III : réapparition d'hommes atteints nés de mères conductrices. Aucune transmission père $\rightarrow$ fils.}
\end{popfigure}

\begin{exemplebox}[Analyse de l'arbre ci-dessus]
\begin{itemize}
\item \textbf{Génération I :} Le père (carré noir) est atteint ($X^d Y$) ; la mère (cercle vide) est saine ($X^D X^D$). Tous leurs fils (II2, II4) reçoivent le Y du père $\rightarrow$ sains ($X^D Y$). Toutes leurs filles (II1, II3) reçoivent le $X^d$ du père $\rightarrow$ conductrices ($X^D X^d$). II1 est hachurée (conductrice identifiée), II3 ne l'est pas mais l'est obligatoirement.
\item \textbf{Génération II :} La fille conductrice II1 ($X^D X^d$) épouse un homme sain ($X^D Y$) $\rightarrow$ parmi leurs enfants, un fils atteint III1 ($X^d Y$, a reçu $X^d$ de sa mère) et une fille saine III2 ($X^D X^D$ ou $X^D X^d$).
\item \textbf{Génération II :} Le fils sain II2 ($X^D Y$) n'a pas de descendance montrée. La fille conductrice II3 ($X^D X^d$, non hachurée) n'a pas de descendance montrée. Le fils atteint II4 ($X^d Y$, a reçu $X^d$ de sa mère II3) épouse une femme saine ($X^D X^D$) $\rightarrow$ leurs enfants : une fille conductrice III4 ($X^D X^d$) et un fils atteint III5 ($X^d Y$).
\item \textbf{Vérification critique :} Aucune transmission père $\rightarrow$ fils n'est observée (I2 $\rightarrow$ II2/II4 sont sains ; II4 $\rightarrow$ III5 est atteint mais III5 a reçu l'allèle de sa mère, pas de son père II4 qui a donné son Y). Cela confirme l'hérédité récessive liée au X.
\end{itemize}
\end{exemplebox}

\courssub{Pièges à éviter et erreurs fréquentes}
\begin{attention}[Erreurs fréquentes au Probatoire]
\begin{itemize}
\item \textbf{Confondre conductrice et malade :} Une femme $X^D X^d$ (ou $X^H X^h$) est \emph{conductrice}, pas malade. Écrire « femme malade » ou « femme atteinte » pour un génotype hétérozygote est une erreur grave qui coûte des points.
\item \textbf{Oublier que le père ne transmet jamais son X à ses fils :} Tout croisement ou arbre généalogique montrant un père atteint et un fils atteint est incompatible avec l'hérédité liée au X. C'est le critère d'exclusion n°1.
\item \textbf{Ne pas distinguer les allèles selon la pathologie :} Utiliser la même lettre pour le daltonisme et l'hémophilie (ex. $d$ pour les deux) entraîne des confusions. Respecter la notation officielle : $D/d$ pour le daltonisme, $H/h$ pour l'hémophilie.
\item \textbf{Lire l'arbre généalogique sans vérifier la génération :} Un homme atteint en génération I ne prouve pas seul le mode de transmission ; il faut observer la descendance (générations II et III) pour valider l'absence de transmission père $\rightarrow$ fils et la présence de femmes conductrices.
\item \textbf{Oublier le sexe dans l'échiquier :} L'échiquier donne des génotypes ; il faut \emph{toujours} déduire le sexe (XX = fille, XY = garçon) pour répondre à la question « proportions de phénotypes parmi les enfants » ou « parmi les fils/filles ».
\end{itemize}
\end{attention}

\courssub{Le savais-tu ?}
\begin{savaistu}[La reine Victoria et l'hémophilie royale]
La reine Victoria d'Angleterre (1819–1901) était conductrice de l'hémophilie B (mutation sur le gène \emph{F9} codant le facteur IX). Elle a transmis l'allèle muté à plusieurs de ses enfants, qui l'ont à leur tour transmis aux cours royales d'Europe (Russie, Espagne, Allemagne). Le tsarévitch Alexis de Russie, arrière-petit-fils de Victoria, était hémophile ; sa maladie a influencé l'histoire politique de la Russie (influence de Raspoutine sur la famille impériale). Cet exemple historique illustre comment une mutation sur le chromosome X peut traverser les générations et les frontières, sans jamais passer de père en fils, et pourquoi le conseil génétique est crucial pour les familles concernées.
\end{savaistu}

\courssub{Synthèse : éradication des préjugés}
La compréhension rigoureuse de l'hérédité liée au sexe permet de déconstruire plusieurs préjugés encore vivaces dans nos sociétés :
\begin{itemize}
\item « La maladie vient toujours du père » $\rightarrow$ \textbf{Faux} : Pour le daltonisme et l'hémophilie récessifs liés au X, l'allèle muté vient de la \textbf{mère} pour les fils atteints. Le père atteint ne transmet son allèle muté qu'à ses filles.
\item « Une femme conductrice est malade » $\rightarrow$ \textbf{Faux} : Elle est phénotypiquement saine (hétérozygote). Seule la femme homozygote mutée ($X^d X^d$ ou $X^h X^h$) est malade, ce qui est très rare.
\item « Le daltonisme ou l'hémophilie sont des « malédictions familiales » ou une « sorcellerie » » $\rightarrow$ \textbf{Faux} : Ce sont des traits génétiques à transmission prévisible, suivant les lois de Mendel appliquées aux chromosomes sexuels. Le conseil génétique (disponible au Centre de Génétique Médicale de la Faculté de Médecine de Yaoundé) permet d'informer les couples sur les risques de récidive.
\end{itemize}
En maîtrisant les échiquiers de croisement et la lecture d'arbres généalogiques, tu es désormais capable d'expliquer scientifiquement la récurrence de ces anomalies dans les familles et de lutter contre les stigmatisations infondées.

\coursec{Hérédité autosomale : albinisme et drépanocytose}

Dans la section précédente, nous avons étudié le daltonisme et l'hémophilie : leurs gènes responsables sont portés par le chromosome sexuel X, ce qui explique que ces anomalies touchent presque exclusivement les garçons et que leur transmission suive des règles particulières. Mais toutes les anomalies héréditaires ne se transmettent pas ainsi. Dans un quartier de Douala, un couple à la peau normalement pigmentée donne naissance à un enfant albinos ; dans un village de l'Adamaoua, deux parents en bonne santé apprennent que leur fille est drépanocytaire. Comment expliquer qu'un caractère apparaisse chez un enfant alors qu'aucun des deux parents ne le présente ? Et pourquoi, cette fois, filles et garçons sont-ils touchés de la même façon ? C'est toute la logique de l'\cle{hérédité autosomale} que nous allons maintenant découvrir.

\courssub{Notion d'hérédité autosomale}

\textbf{Observation.}\par
Recensons dans une même famille les cas d'albinisme sur trois générations : on constate que des filles et des garçons sont atteints, en proportions comparables, et que l'anomalie peut « sauter » une génération. Ce schéma de transmission est très différent de celui du daltonisme.

\textbf{Hypothèse.}\par
Le gène responsable n'est pas situé sur un chromosome sexuel (gonosome X ou Y), mais sur l'une des 22 paires de chromosomes non sexuels, les \cle{autosomes}.

\textbf{Vérification.}\par
L'étude des caryotypes et des arbres généalogiques confirme cette hypothèse : le gène de l'albinisme est porté par un autosome, et il en va de même pour le gène de l'hémoglobine impliqué dans la drépanocytose.

\begin{definition}[Hérédité autosomale et vocabulaire des allèles]
On parle d'\cle{hérédité autosomale} lorsque le gène étudié est porté par un \cle{autosome} (l'un des 22 chromosomes non sexuels). La transmission du caractère est alors \textbf{indépendante du sexe} : filles et garçons héritent du gène et expriment le caractère avec les mêmes probabilités.
\begin{itemize}
\item Un \cle{allèle dominant} est un allèle qui s'exprime dès qu'il est présent en un seul exemplaire. On le note par une lettre majuscule : $A$.
\item Un \cle{allèle récessif} ne s'exprime que s'il est présent en deux exemplaires. On le note par une lettre minuscule : $a$.
\item Un individu est \cle{homozygote} pour un gène s'il possède deux allèles identiques ($AA$ ou $aa$).
\item Un individu est \cle{hétérozygote} s'il possède deux allèles différents ($Aa$).
\end{itemize}
\end{definition}

\begin{aretenir}[L'essentiel]
Dans une hérédité autosomale, chaque parent transmet un exemplaire du gène à son enfant, quel que soit le sexe de celui-ci. Le phénotype dépend de la dominance : un allèle récessif peut être « caché » chez un hétérozygote de phénotype sain, puis réapparaître chez un descendant homozygote.
\end{aretenir}

\courssub{L'albinisme : un caractère autosomique récessif}

\textbf{Situation concrète.}\par
Dans les écoles camerounaises, on rencontre parfois des élèves albinos : peau très claire, cheveux blancs ou jaunâtres, yeux clairs et sensibles à la lumière. Leurs parents ont le plus souvent une pigmentation tout à fait normale. D'où vient ce caractère ?

\textbf{Explication scientifique.}\par
La couleur de la peau, des cheveux et des yeux est due à un pigment, la \cle{mélanine}, fabriquée par des cellules appelées mélanocytes. La synthèse de la mélanine exige une enzyme (la tyrosinase) dont le gène existe sous deux formes (allèles) :
\begin{itemize}
\item l'allèle $A$ (dominant) : permet la production normale de mélanine ;
\item l'allèle $a$ (récessif) : ne permet pas la production de mélanine (enzyme non fonctionnelle).
\end{itemize}

\begin{definition}[Génotypes et phénotypes pour l'albinisme]
\begin{itemize}
\item $AA$ : individu sain, à pigmentation normale (homozygote dominant) ;
\item $Aa$ : individu sain, à pigmentation normale, mais \textbf{porteur} de l'allèle albinos (hétérozygote) ;
\item $aa$ : individu \textbf{albinos} : absence de mélanine (homozygote récessif).
\end{itemize}
L'albinisme est donc une anomalie autosomique \textbf{récessive} : il faut recevoir l'allèle $a$ des deux parents pour être atteint.
\end{definition}

\textbf{Question d'investigation.}\par
Deux parents sains, tous deux hétérozygotes $Aa$, peuvent-ils avoir un enfant albinos ? Construisons l'échiquier de croisement : chaque parent produit deux types de gamètes, $A$ et $a$, en proportions égales ($1/2$ chacun).

\begin{popfigure}
\begin{center}
\begin{tikzpicture}[scale=1.2]
% Grille 2x2 pour les résultats (colonnes 1-2, lignes 1-2)
\draw[thick, popdark] (1,1) grid (3,3);
% En-têtes gamètes père (au-dessus des colonnes de résultats)
\node[align=center] at (2,3.4) {\textbf{Gamètes du père}};
\node at (1.5,3.1) {$A$};
\node at (2.5,3.1) {$a$};
% En-têtes gamètes mère (à gauche des lignes de résultats)
\node[align=center, rotate=90] at (0.4,2) {\textbf{Gamètes de la mère}};
\node at (0.7,2.5) {$A$};
\node at (0.7,1.5) {$a$};
% Remplissage des cases résultats
\fill[popgreenL] (1,2) rectangle (2,3); \node at (1.5,2.5) {$AA$};
\fill[popgreenL] (2,2) rectangle (3,3); \node at (2.5,2.5) {$Aa$};
\fill[popgreenL] (1,1) rectangle (2,2); \node at (1.5,1.5) {$Aa$};
\fill[poporangeL] (2,1) rectangle (3,2); \node at (2.5,1.5) {$aa$};
% Légende proportions
\node[align=left, anchor=west] at (3.3,2.5) {\textbf{Résultats :}};
\node[align=left, anchor=west] at (3.3,2.0) {$1/4\ AA$ : sain, non porteur};
\node[align=left, anchor=west] at (3.3,1.5) {$2/4\ Aa$ : sain, porteur};
\node[align=left, anchor=west] at (3.3,1.0) {$1/4\ aa$ : \textbf{albinos}};
\node[align=left, anchor=west] at (3.3,0.4) {Phénotypes : $3/4$ sains,\\ $1/4$ albinos};
\end{tikzpicture}
\end{center}
\legende{Échiquier de croisement $Aa \times Aa$ (albinisme). Les gamètes paternels sont en colonnes, maternels en lignes. Chaque case (issue de fécondation) est équiprobable. Vert : enfants sains ($AA$ ou $Aa$) ; Orange : enfant albinos ($aa$).}
\end{popfigure}

\begin{propriete}[Risque de transmission d'un caractère autosomique récessif]
Deux parents sains mais hétérozygotes ($Aa \times Aa$) ont, \textbf{à chaque naissance}, une probabilité de :
\begin{itemize}
\item $1/4$ d'avoir un enfant homozygote $AA$ (sain, non porteur) ;
\item $2/4$ d'avoir un enfant hétérozygote $Aa$ (sain, porteur) ;
\item $1/4$ d'avoir un enfant $aa$ (atteint).
\end{itemize}
Démonstration exigible : elle se fait par l'échiquier de croisement ci-dessus, en croisant les gamètes $A$ et $a$ de chaque parent. Un couple sain peut donc avoir un enfant atteint d'une anomalie récessive.
\end{propriete}

\begin{popfigure}
\begin{center}
\begin{tikzpicture}[
  male/.style={rectangle, draw=popdark, minimum size=0.6cm, inner sep=0pt, thick},
  female/.style={circle, draw=popdark, minimum size=0.6cm, inner sep=0pt, thick},
  affected/.style={fill=poporange},
  link/.style={thick, popdark}]
% Génération I (Grands-parents)
\node[male] (I1) at (0,3) {};
\node[female] (I2) at (2,3) {};
\draw[link] (I1) -- (I2);
% Ligne de descendance vers Génération II
\draw[link] (1,3) -- (1,2.2);
% Génération II (Parents)
\node[male] (II1) at (0,1.2) {};
\node[female] (II2) at (2,1.2) {};
\draw[link] (II1) -- (II2);
% Ligne de descendance vers Génération III
\draw[link] (1,1.2) -- (1,0.4);
% Génération III (Enfants) - 3 enfants
\node[male] (III1) at (0.2,-0.5) {};
\node[female, affected] (III2) at (1.0,-0.5) {};
\node[male] (III3) at (1.8,-0.5) {};
\draw[link] (0.2,0.4) -- (0.2,-0.1);
\draw[link] (1.0,0.4) -- (1.0,-0.1);
\draw[link] (1.8,0.4) -- (1.8,-0.1);
% Étiquettes
\node[above=2pt of I1] {$I_1$ ($Aa$)};
\node[above=2pt of I2] {$I_2$ ($Aa$)};
\node[left=4pt of II1] {$II_1$ ($Aa$)};
\node[right=4pt of II2] {$II_2$ ($Aa$)};
\node[below=3pt of III1] {$III_1$ : sain};
\node[below=3pt of III2] {$III_2$ : \textbf{albinos} ($aa$)};
\node[below=3pt of III3] {$III_3$ : sain};
\end{tikzpicture}
\end{center}
\legende{Arbre généalogique d'une famille où l'albinisme apparaît chez une fille ($III_2$, symbole colorié) alors que ses parents ($II_1$, $II_2$) et ses grands-parents sont sains. Carrés = garçons, Cercles = filles ; symbole colorié = individu atteint.}
\end{popfigure}

\begin{attention}[Le porteur sain n'est pas malade, mais il transmet]
Une erreur très fréquente consiste à croire qu'un individu $Aa$ est « un peu malade ». C'est faux : le porteur sain possède un allèle $A$ fonctionnel qui suffit à assurer la pigmentation normale. Il est \textbf{sain}, mais chacun de ses gamètes a une chance sur deux de porter l'allèle $a$. De même, un individu $AS$ pour la drépanocytose n'est pas un malade. Autre piège : les probabilités $1/4$, $2/4$, $1/4$ s'appliquent \textbf{à chaque naissance}, indépendamment des enfants déjà nés. Un couple $Aa \times Aa$ qui a déjà un enfant albinos a toujours $1/4$ de risque à la grossesse suivante.
\end{attention}

\courssub{La drépanocytose : une anomalie de l'hémoglobine}

\textbf{Situation concrète.}\par
Au Cameroun, lors des campagnes de santé scolaire ou avant certains mariages, on propose parfois un test sanguin appelé « test d'électrophorèse de l'hémoglobine ». Pourquoi ? Parce que la \cle{drépanocytose} (ou anémie falciforme) est l'anomalie génétique la plus fréquente en Afrique subsaharienne, et que son évitement passe par l'information des couples.

\textbf{Explication scientifique.}\par
Les globules rouges transportent le dioxygène grâce à une protéine, l'\cle{hémoglobine}. Le gène de la chaîne $\beta$-globine de l'hémoglobine, situé sur le chromosome 11 (autosome), existe sous deux allèles principaux :
\begin{itemize}
\item l'allèle $A$ : code une hémoglobine normale (HbA) ;
\item l'allèle $S$ : code une hémoglobine anormale (HbS) qui, en cas de faible oxygénation, polymérise et déforme le globule rouge en \textbf{faucille} (d'où le nom « drépanocytose », du grec \emph{drepanê}, faucille).
\end{itemize}
Les globules rouges en faucille sont rigides, fragiles, se détruisent rapidement (anémie hémolytique) et peuvent boucher les petits vaisseaux sanguins, provoquant des crises vaso-occlusives douloureuses et une sensibilité accrue aux infections.

\begin{popfigure}
\begin{center}
\begin{tikzpicture}[scale=1.1]
% Globule rouge normal
\shade[inner color=poppinkL, outer color=poppink] (0,0) circle (0.9);
\shade[inner color=poppink!40, outer color=poppinkL] (0,0) circle (0.45);
\node at (0,-1.5) {Globule rouge normal};
\node at (0,-2.0) {(disque biconcave souple)};
% Globule rouge drépanocytaire (faucille)
\begin{scope}[xshift=5.2cm]
\fill[poppurple!60] (0,0.7) .. controls (1.1,0.55) and (1.25,-0.15) .. (0.85,-0.75)
 .. controls (0.95,-0.1) and (0.55,0.35) .. (-0.15,0.35)
 .. controls (-0.35,0.45) and (-0.25,0.65) .. (0,0.7) -- cycle;
\node at (0.45,-1.5) {Globule rouge drépanocytaire};
\node at (0.45,-2.0) {(forme de faucille, rigide)};
\end{scope}
% Légendes fléchées
\draw[thin, popmuted] (0.9,0.6) -- (1.7,1.1) node[right, popdark] {membrane souple};
\draw[thin, popmuted] (6.05,0.5) -- (6.9,1.1) node[right, popdark] {cellule rigide et fragile};
\end{tikzpicture}
\end{center}
\legende{Comparaison d'un globule rouge normal (disque souple, riche en hémoglobine A) et d'un globule rouge drépanocytaire (déformé en faucille par l'hémoglobine S polymérisée). La faucille se détruit vite (anémie) et obstrue les capillaires (crises).}
\end{popfigure}

\begin{definition}[Génotypes et phénotypes pour la drépanocytose]
\begin{itemize}
\item $AA$ : individu sain, hémoglobine normale (HbA) ;
\item $AS$ : \textbf{porteur sain} (on dit aussi « trait drépanocytaire ») : il fabrique les deux types d'hémoglobine (HbA et HbS), ne présente généralement aucun symptôme en conditions normales, mais peut transmettre l'allèle $S$ ;
\item $SS$ : individu \textbf{drépanocytaire}, atteint de la forme grave de la maladie (anémie chronique, crises vaso-occlusives douloureuses, infections fréquentes, risque d'accident vasculaire cérébral).
\end{itemize}
La drépanocytose est une anomalie autosomique : l'allèle $S$ est récessif au sens où la maladie grave n'apparaît qu'à l'état homozygote $SS$.
\end{definition}

\textbf{Croisement type.}\par
Si deux porteurs sains se marient ($AS \times AS$), l'échiquier de croisement est identique en structure à celui de l'albinisme :

\begin{center}
\begin{tabular}{|c|c|c|}
\hline
\rowcolor{popblueL} \textbf{Gamètes} & $A$ & $S$ \\
\hline
$A$ & $AA$ (sain) & $AS$ (porteur sain) \\
\hline
$S$ & $AS$ (porteur sain) & $SS$ (drépanocytaire) \\
\hline
\end{tabular}
\end{center}

Soit, à chaque naissance : $1/4\ AA$, $2/4\ AS$, $1/4\ SS$. C'est exactement ici que réside l'intérêt du \cle{dépistage avant mariage} : un couple dont les deux membres sont $AS$ sait qu'il a un risque de $1/4$ d'avoir un enfant atteint de la forme grave, et peut ainsi faire des choix éclairés.

\begin{savaistu}[La drépanocytose au Cameroun]
Au Cameroun, environ \textbf{20 à 25 \%} de la population est porteuse du trait drépanocytaire ($AS$), soit près d'une personne sur quatre ! Le dépistage prénuptial (électrophorèse de l'hémoglobine) est recommandé par les autorités sanitaires. Fait remarquable : l'allèle $S$ confère aux porteurs $AS$ une \textbf{protection partielle contre le paludisme grave} (le parasite \emph{Plasmodium falciparum} se développe mal dans les globules rouges contenant de l'HbS), ce qui explique pourquoi cet allèle s'est maintenu à haute fréquence dans les régions où le paludisme sévit, comme la zone forestière et soudanienne camerounaise. C'est un bel exemple d'équilibre évolutif (avantage hétérozygote) entre avantage et inconvénient d'un allèle.
\end{savaistu}

\courssub{Lecture d'un arbre généalogique autosomique}

\begin{methode}[Reconnaître une hérédité autosomale dans un arbre généalogique]
\begin{enumerate}
\item \textbf{Observer la répartition par sexe} : si des filles ET des garçons sont atteints en proportions comparables, l'hérédité est probablement autosomique (une hérédité liée à X touche presque uniquement les garçons).
\item \textbf{Chercher une fille atteinte} : si une fille est atteinte alors que son père est sain, le caractère ne peut pas être récessif lié à X (une fille $X^aX^a$ exigerait un père $X^aY$, donc atteint). Conclusion : hérédité autosomique.
\item \textbf{Repérer les « sauts de génération »} : si des parents sains ont un enfant atteint, le caractère est \textbf{récessif} et les deux parents sont hétérozygotes.
\item \textbf{Conclure et vérifier} : attribuer les génotypes obligatoires (les atteints sont homozygotes récessifs), puis vérifier la cohérence de tout l'arbre.
\end{enumerate}
\end{methode}

\begin{exoresolu}[Interprétation d'un arbre généalogique]
\textbf{Énoncé.} Dans une famille, les deux parents (II) sont sains. Ils ont quatre enfants : deux garçons sains, une fille saine et une fille drépanocytaire. (1) Le caractère est-il dominant ou récessif ? Est-il autosomique ou lié au sexe ? Justifier. (2) Donner le génotype des parents et de la fille atteinte. (3) Calculer la probabilité que le prochain enfant du couple soit drépanocytaire.
\tcblower
\textbf{Solution détaillée.}

(1) Des parents \textbf{sains} ont un enfant \textbf{atteint} : l'allèle responsable était donc présent mais masqué chez les parents ; le caractère est \textbf{récessif}. De plus, une \textbf{fille} est atteinte alors que son père est sain : si le gène était lié à X, la fille atteinte ($X^S X^S$) aurait dû recevoir un chromosome $X^S$ de son père, qui serait alors atteint — contradiction. Le caractère est donc \textbf{autosomique récessif}.

(2) La fille atteinte est homozygote $SS$. Chaque parent lui a transmis un allèle $S$ tout en étant sain : les deux parents sont donc hétérozygotes $AS$ (porteurs sains).

(3) Croisement $AS \times AS$ : d'après l'échiquier de croisement, les issues sont $1/4\ AA$, $2/4\ AS$, $1/4\ SS$. La probabilité que le prochain enfant soit drépanocytaire ($SS$) est donc de $\dfrac{1}{4}$, soit $25\ \%$, quel que soit son sexe.
\end{exoresolu}

\courssub{Éradiquer les préjugés : comprendre pour respecter}

La science démontre que l'albinisme et la drépanocytose sont des \textbf{caractères héréditaires naturels}, inscrits dans les chromosomes reçus des parents, comme la couleur des yeux ou le groupe sanguin. Il faut donc combattre activement les idées fausses :

\begin{itemize}
\item Ces anomalies ne sont \textbf{pas contagieuses} : on ne « attrape » ni l'albinisme ni la drépanocytose par contact, par le partage de repas ou par l'amitié.
\item Elles n'ont \textbf{rien de surnaturel} : ni malédiction, ni faute des parents, ni punition ; ce sont des combinaisons d'allèles parfaitement expliquées par la génétique.
\item Les personnes albinos ont besoin d'une \textbf{protection solaire} (chapeaux, vêtements couvrants, crèmes écran total) car leur peau, dépourvue de mélanine, est très vulnérable aux rayons ultraviolets (risque de coups de soleil, vieillissement cutané précoce, cancers de la peau) ; elles ont droit à la sécurité, à l'éducation et au respect, comme tout élève.
\item Les personnes drépanocytaires ($SS$) ont besoin d'un \textbf{suivi médical régulier} (hydratation, vaccination renforcée, acide folique, traitement préventif ou curatif des crises) ; les porteurs sains ($AS$) mènent une vie parfaitement normale.
\end{itemize}

\begin{aretenir}[Conclusion de la section]
L'hérédité autosomale concerne les gènes portés par les 22 paires d'autosomes ; elle est indépendante du sexe. L'albinisme ($aa$) et la drépanocytose grave ($SS$) sont des anomalies autosomiques récessives : deux parents sains hétérozygotes ont $1/4$ de risque d'avoir un enfant atteint à chaque naissance. Le dépistage et l'information des couples sont des outils de prévention ; la connaissance scientifique est l'arme la plus efficace contre les préjugés et pour la dignité des personnes concernées.
\end{aretenir}

\begin{exercice}[À toi de jouer]
Un homme de génotype $AS$ épouse une femme de génotype $AA$. (1) Construire l'échiquier de croisement. (2) Ce couple peut-il avoir un enfant drépanocytaire ($SS$) ? (3) Quelle est la probabilité qu'un enfant de ce couple soit porteur sain ?
\end{exercice}

\begin{corrige}[Exercice 1]
(1) Le père $AS$ produit des gamètes $A$ ($1/2$) et $S$ ($1/2$) ; la mère $AA$ ne produit que des gamètes $A$. L'échiquier donne : $1/2\ AA$ et $1/2\ AS$.
(2) Non : l'enfant $SS$ exigerait un allèle $S$ de chaque parent, or la mère $AA$ ne peut pas en fournir. Ce couple ne peut donc pas avoir d'enfant drépanocytaire.
(3) La probabilité qu'un enfant soit porteur sain ($AS$) est de $1/2$, soit $50\ \%$.
\end{corrige}

\coursec{Synthèse : démarche d'interprétation et lutte contre les préjugés}

Après avoir exploré tour à tour l'hérédité liée au sexe (daltonisme, hémophilie) et l'hérédité autosomale (albinisme, drépanocytose), nous disposons maintenant d'une boîte à outils complète pour analyser la transmission des caractères chez l'homme. Cette section finale a un double objectif : \textbf{structurer votre raisonnement} face à tout document (caryotype ou arbre généalogique) pour réussir l'épreuve du Probatoire, et \textbf{déconstruire scientifiquement} les préjugés tenaces qui pèsent sur les familles touchées par ces anomalies, ici au Cameroun comme ailleurs.

\courssub{Bilan comparatif : deux modes de transmission, des logiques distinctes}

Commençons par mettre en regard les deux grands types d'hérédité mendélienne étudiés. La différence fondamentale réside dans le \textbf{support chromosomique} du gène concerné, ce qui dicte toute la suite : sexes touchés, transmission par les pères/mères, et risques de récidive.

\begin{popfigure}
\begin{tikzpicture}[
    node distance=0.5cm and 1cm,
    header/.style={fill=popblue, text=white, font=\bfseries, minimum height=0.9cm, align=center},
    cell/.style={minimum height=0.7cm, align=center, draw=popline},
    cellL/.style={cell, fill=popblueL, font=\bfseries, text=popdark},
    cellC/.style={cell, fill=popcream},
    table/.style={matrix of nodes, nodes in empty cells, row sep=-\pgflinewidth, column sep=-\pgflinewidth, anchor=north west}
]
\matrix (tab) [table, column 1/.style={nodes={cellL, minimum width=3.5cm}}, column 2/.style={nodes={cellC, minimum width=5.5cm}}, column 3/.style={nodes={cellC, minimum width=5.5cm}}]
{
    \node[header, minimum width=3.5cm] {Critère de comparaison}; & \node[header, minimum width=5.5cm] {\textbf{Hérédité autosomale} (ex. albinisme, drépanocytose)}; & \node[header, minimum width=5.5cm] {\textbf{Hérédité liée au sexe} (ex. daltonisme, hémophilie)}; \\
    \textbf{Localisation du gène} & Sur un \textbf{autosome} (paires 1 à 22). & Sur un \textbf{gonosome} (chromosome X le plus souvent). \\
    \textbf{Sexes touchés} & \textbf{Filles et garçons} dans les mêmes proportions. & Principalement les \textbf{garçons} (hémizygotes $X^dY$) ; filles atteintes rares (homozygotes $X^dX^d$). \\
    \textbf{Transmission père $\to$ fils} & \textbf{Possible} (père transmet un autosome). & \textbf{Impossible} (père transmet son Y à son fils, son X à sa fille). \\
    \textbf{Transmission mère $\to$ fils} & Possible (mère transmet un autosome). & \textbf{Caractéristique} : mère porteuse ($X^DX^d$) transmet $X^d$ à 1/2 de ses fils $\to$ atteints. \\
    \textbf{Transmission père $\to$ fille} & Possible. & \textbf{Systématique} si père atteint ($X^dY$) : toutes les filles reçoivent $X^d$ (porteuses ou atteintes). \\
    \textbf{Arbre généalogique typique} & Atteints dans les deux sexes ; transmission horizontale (fratrie) ou verticale ; pas de saut de génération systématique. & \textbf{Saute les générations} (père atteint $\to$ filles porteuses $\to$ petits-fils atteints) ; prédominance masculine. \\
    \textbf{Exemple camerounais} & Couple $Aa \times Aa$ (drépanocytose) : risque 1/4 enfant $aa$ (SS) à chaque grossesse. & Mère porteuse hémophilie ($X^HX^h$) $\times$ père sain ($X^HY$) : 1/2 fils hémophiles, 1/2 filles porteuses. \\
};
\end{tikzpicture}
\legende{Tableau comparatif : hérédité autosomale vs hérédité liée au sexe (récessives). Ce tableau est un outil de révision indispensable pour le Probatoire.}
\end{popfigure}

\begin{aretenir}[Distinction cruciale : anomalie chromosomique vs anomalie génique]
Ne confondez jamais les deux niveaux d'organisation :
\begin{itemize}
    \item \textbf{Anomalie chromosomique} (ex. trisomie 21, syndrome de Turner, Klinefelter) : \textbf{visible sur le caryotype}. Elle concerne un \textbf{excès ou un déficit de chromosomes entiers} (ou de gros fragments). La formule chromosomique s'écrit : $2n+1$, $2n-1$, $47,XX,+21$, $45,X$, etc. Elle n'est \textbf{pas} transmise de manière mendélienne classique (sauf translocations équilibrées parentales).
    \item \textbf{Anomalie génique} (ex. albinisme, drépanocytose, hémophilie, daltonisme) : \textbf{invisible sur le caryotype standard} (les chromosomes ont une morphologie normale). Elle concerne une \textbf{mutation ponctuelle d'un gène}. Le caryotype est normal ($46,XX$ ou $46,XY$). La transmission suit les lois de Mendel (autosomale ou liée au sexe).
\end{itemize}
\end{aretenir}

\courssub{La démarche d'identification en trois étapes : votre « GPS » pour le Probatoire}

Face à un sujet d'examen (arbre généalogique ou caryotype), ne vous précipitez pas. Suivez cette méthode rigoureuse, étape par étape.

\begin{popfigure}
\begin{tikzpicture}[
    node distance=1cm and 1.5cm,
    startstop/.style={rectangle, rounded corners, minimum width=3.5cm, minimum height=1.2cm, text centered, draw=popgreen, fill=popgreenL, font=\bfseries},
    process/.style={rectangle, minimum width=4.5cm, minimum height=1.1cm, text centered, draw=popblue, fill=popblueL, font=\bfseries, text width=4.2cm, align=center},
    decision/.style={diamond, aspect=1.5, minimum width=4cm, minimum height=1.3cm, text centered, draw=poporange, fill=poporangeL, font=\bfseries, text width=3.5cm, align=center},
    arrow/.style={thick, ->, >=Stealth, popdark}
]
\node (start) [startstop, align=center]{DOCUMENT FOURNI\\ (Caryotype OU Arbre généalogique)};
\node (obs) [process, below=1cm of start, align=center]{\textbf{ÉTAPE 1 : OBSERVER \& CARACTÉRISER}\\\textit{Caryotype :} Nombre chr, paires, sexe, anomalies visibles (trisomie, délétion...)\\\textit{Arbre :} Sexes atteints, générations sautées, transmission père/fils ?};
\node (type) [decision, below=1cm of obs, align=center]{\textbf{ÉTAPE 2 : IDENTIFIER LE TYPE}\\\textit{Anomalie chromosomique ?} (Caryotype anormal)\\\textit{Anomalie génique ?} (Caryotype normal + arbre)\\\quad $\to$ Autosomale ou Liée au sexe ?};
\node (auto) [process, below left=1.5cm and 2.5cm of type, align=center]{\textbf{Cas Autosomale}\\\textit{Allèle :} $A$ (sain), $a$ (malade)\\\textit{Génotypes parents ?}\\\textit{Échiquier de croisement}\\\textit{Proportions phénotypes}};
\node (sex) [process, below right=1.5cm and 2.5cm of type, align=center]{\textbf{Cas Liée au sexe (X)}\\\textit{Allèles :} $X^A$, $X^a$\\\textit{Génotypes parents (♀ $X^AX^a$, ♂ $X^AY$)}\\\textit{Échiquier (gonosomes)}\\\textit{Proportions par sexe}};
\node (chromo) [process, below=3.5cm of type, align=center]{\textbf{Cas Chromosomique}\\\textit{Formule chromosomique}\\\textit{Ex: } 47,XY,+21 \textit{ (Trisomie 21)}\\\textit{Mécanisme : non-disjonction méiose}\\\textit{Risque récidive faible (sauf translocation)}};
\node (concl) [startstop, below=1.5cm of chromo, align=center]{\textbf{ÉTAPE 3 : CONCLURE RIGOUREUSEMENT}\\\textit{Écrire le croisement, les proportions, le risque de récidive}\\\textit{Utiliser le vocabulaire précis : homozygote, hétérozygote, hémizygote, porteur sain}};

\draw [arrow] (start) -- (obs);
\draw [arrow] (obs) -- (type);
\draw [arrow] (type) -- node[left, pos=0.5, fill=white, font=\small, align=center] {Anomalie génique\\autosomale} (auto);
\draw [arrow] (type) -- node[right, pos=0.5, fill=white, font=\small, align=center] {Anomalie génique\\liée au sexe} (sex);
\draw [arrow] (type) -- node[left, pos=0.5, fill=white, font=\small, align=center] {Caryotype\\anormal} (chromo);
\draw [arrow] (auto) |- (concl);
\draw [arrow] (sex) |- (concl);
\draw [arrow] (chromo) -- (concl);
\end{tikzpicture}
\legende{Organigramme de la démarche d'identification en trois étapes. Suivez la flèche correspondant à votre document.}
\end{popfigure}

\courssub{Étape 1 : Observer et caractériser le document}

\begin{methode}[Lecture active du document]
\begin{enumerate}
    \item \textbf{Si caryotype :} Comptez le nombre total de chromosomes. Repérez les gonosomes (XX ou XY). Cherchez un chromosome en triple (trisomie), manquant (monosomie), ou un fragment déplacé (translocation). \textit{Exemple :} « Caryotype à 47 chromosomes, présence de 3 chromosomes 21 $\to$ Trisomie 21, sexe masculin (XY) ».
    \item \textbf{Si arbre généalogique :} Notez les symboles : $\square$ = homme, $\bigcirc$ = femme, noircis = atteints, hachurés = porteurs sains (parfois). Repérez : \textbf{qui est atteint ?} (hommes seulement ? les deux sexes ?), \textbf{comment ça se transmet ?} (père $\to$ fils ? mère $\to$ fils ? saut de génération ?).
\end{enumerate}
\end{methode}

\courssub{Étape 2 : Identifier le type d'anomalie et de transmission}

C'est l'étape diagnostique. Utilisez le tableau comparatif ci-dessus.

\begin{itemize}
    \item \textbf{Caryotype anormal} $\to$ \textbf{Anomalie chromosomique}. Donnez la formule (ex. $47,XX,+21$). N'essayez pas de faire un échiquier mendélien.
    \item \textbf{Caryotype normal ($46,XX$ ou $46,XY$) + Arbre} $\to$ \textbf{Anomalie génique}.
    \begin{itemize}
        \item Atteints des deux sexes, transmission père $\to$ fils possible $\to$ \textbf{Autosomale} (souvent récessive si parents sains ont enfant atteint).
        \item Atteints quasi-exclusivement mâles, pas de transmission père $\to$ fils, saut de génération (grand-père $\to$ petit-fils via fille porteuse) $\to$ \textbf{Liée au sexe (récessive X)}.
    \end{itemize}
\end{itemize}

\begin{attention}[Piège mortel : conclure sur un seul individu]
\textbf{Ne déduisez JAMAIS le mode de transmission à partir d'un seul individu atteint} (ex. « le fils est atteint, donc c'est récessif »). Un arbre généalogique se lit \textbf{globalement} : il faut au minimum observer une fratrie, ou deux générations successives, pour distinguer une autosomale récessive (parents sains, 1/4 atteints) d'une liée au sexe (mère porteuse, 1/2 fils atteints). Si l'arbre ne montre qu'un couple et un seul enfant, \textbf{on ne peut pas conclure} : il faut écrire « Information insuffisante pour déterminer le mode de transmission ».
\end{attention}

\courssub{Étape 3 : Conclure par un croisement ou une formule chromosomique}

C'est là que vous marquez les points. La rédaction doit être standardisée.

\begin{methode}[Fiche de résolution type Probatoire — 4 étapes gagnantes]
\begin{enumerate}
    \item \textbf{Lire le document} et énoncer le type d'anomalie (chromosomique / génique autosomale / génique liée au sexe) et le caractère (récessif/dominant).
    \item \textbf{Poser les allèles et les génotypes des parents} (déduits de leurs phénotypes et de l'existence d'enfants atteints).
    \item \textbf{Construire l'échiquier de croisement} (carré de Punnett) : gamètes en lignes/colonnes, zygotes dans les cases.
    \item \textbf{Conclure avec les proportions phénotypiques} (et génotypiques si demandé), \textbf{séparées par sexe} si hérédité liée au sexe. Écrire une phrase de conclusion : « Le risque de récidive à chaque grossesse est de... ».
\end{enumerate}
\end{methode}

\begin{exemplebox}[Synthèse des risques : les valeurs repères à connaître par cœur]
Voici les trois configurations classiques du programme de Première A. Maîtrisez-les pour ne pas recalculer à l'examen.

\begin{center}
\begin{tabularx}{\linewidth}{|l|X|X|X|}\hline
\rowcolor{popblue} \textbf{Configuration} & \textbf{Croisement type} & \textbf{Résultat (Proportions phénotypiques)} & \textbf{Risque de récidive (enfant atteint)} \\ \hline
\textbf{Autosomale récessive} (Albinisme, Drépanocytose) & $Aa \times Aa$ (deux porteurs sains) & $\frac{1}{4}\ AA$ (sain) ; $\frac{1}{2}\ Aa$ (porteur sain) ; $\frac{1}{4}\ aa$ (\textbf{atteint}) & \textbf{1/4} (25\%) pour \textbf{chaque enfant}, \textbf{indépendamment du sexe} \\ \hline
\textbf{Liée au sexe récessive} (Hémophilie, Daltonisme) & $\text{♀ } X^AX^a \times \text{♂ } X^AY$ (mère porteuse, père sain) & \textbf{Filles :} 1/2 $X^AX^A$ (saines), 1/2 $X^AX^a$ (porteuses)\\
\textbf{Garçons :} 1/2 $X^AY$ (sains), 1/2 $X^aY$ (\textbf{atteints}) & \textbf{1/2 des garçons} (50\%) \\
\textbf{0\% des filles} (atteintes) \\ \hline
\textbf{Liée au sexe récessive} (Père atteint $\times$ Mère saine) & $\text{♀ } X^AX^A \times \text{♂ } X^aY$ & \textbf{Filles :} 100\% $X^AX^a$ (\textbf{porteuses saines})\\
\textbf{Garçons :} 100\% $X^AY$ (\textbf{sains}) & \textbf{0\%} (aucun enfant atteint) \\
& & & \textit{Toutes les filles sont porteuses} \\ \hline
\end{tabularx}
\end{center}
\legende{Tableau récapitulatif des risques de récidive pour les trois situations types du programme. Notez l'indépendance du sexe pour l'autosomale, et la dépendance stricte au sexe pour l'hérédité X.}
\end{exemplebox}

\begin{exoresolu}[Analyse complète d'un arbre généalogique : Drépanocytose]
\textbf{Énoncé :} On étudie une famille camerounaise. Le père et la mère sont phénotypiquement sains. Ils ont 4 enfants : 2 filles saines, 1 garçon atteint de drépanocytose (SS), 1 garçon sain. On demande : 1) Le mode de transmission. 2) Les génotypes des parents. 3) La proportion théorique d'enfants atteints. 4) Le risque pour le garçon sain d'être porteur.

\textbf{Solution détaillée :}
\begin{enumerate}
    \item \textbf{Observation :} Les deux parents sont sains. Ils ont un enfant atteint (garçon). Les deux sexes sont représentés dans la fratrie (filles saines, garçons dont un atteint). \textbf{Transmission père $\to$ fils possible} (père sain $\to$ fils atteint). \textbf{Conclusion : Hérédité autosomale récessive.} (Si c'était lié au sexe, un père sain $X^AY$ ne pourrait pas avoir de fils atteint $X^aY$).
    \item \textbf{Allèles :} $A$ = allèle normal (hémoglobine A), $a$ = allèle muté (hémoglobine S). Phénotypes : $AA$ = sain (AA), $Aa$ = porteur sain (AS), $aa$ = atteint (SS).
    \item \textbf{Génotypes parents :} Parents sains ayant un enfant $aa$ $\to$ obligatoirement hétérozygotes $Aa$ (AS) tous les deux.
    \item \textbf{Échiquier de croisement $Aa \times Aa$ :}
    \begin{center}
    \begin{tabular}{|c|c|c|}\hline
     & $A$ & $a$ \\ \hline
    $A$ & $AA$ (AA sain) & $Aa$ (AS porteur) \\ \hline
    $a$ & $Aa$ (AS porteur) & $aa$ (\textbf{SS atteint}) \\ \hline
    \end{tabular}
    \end{center}
    \item \textbf{Proportions :} 1/4 $AA$ (sain non porteur), 1/2 $Aa$ (porteur sain AS), 1/4 $aa$ (atteint SS). \textbf{Risque de récidive = 1/4 (25\%)} pour chaque grossesse future, fille ou garçon.
    \item \textbf{Risque pour le garçon sain :} Il est phénotypiquement sain. Parmi les phénotypes « sains » ($AA$ ou $Aa$), la proportion de porteurs ($Aa$) est $\frac{1/2}{1/4+1/2} = \frac{2}{3}$. \textbf{Il a 2 chances sur 3 d'être porteur sain (AS).}
\end{enumerate}
\end{exoresolu}

\courssub{Argumentaire scientifique : déconstruire les préjugés, un devoir de citoyenneté}

Au Cameroun, comme dans de nombreuses sociétés, les anomalies héréditaires (albinisme, drépanocytose, hémophilie, trisomie 21) sont encore trop souvent entourées de croyances non scientifiques : « c'est la faute de la mère », « c'est un sort jeté par la belle-famille », « l'enfant est un esprit/« enfant-sorcier » », « c'est contagieux », « il faut cacher l'enfant ». En tant que futurs citoyens éclairés et élèves de SVT, vous devez opposer la \textbf{rigueur biologique} à ces préjugés nuisibles.

\begin{definition}[Préjugé vs Fait scientifique]
Un \textbf{préjugé} est un jugement porté à l'avance, sans examen critique, souvent basé sur la tradition, la peur de l'inconnu ou la recherche d'un responsable. Un \textbf{fait scientifique} est une observation reproductible, validée par la communauté scientifique, explicable par des mécanismes connus (mécanismes moléculaires, cellulaires, lois statistiques).
\end{definition}

\begin{propriete}[Trois préjugés majeurs et leur réfutation scientifique]
\begin{enumerate}
    \item \textbf{Préjugé : « C'est la femme qui détermine le sexe de l'enfant / C'est sa faute s'il est fille/garçon ou malade. »} \\
    \textbf{Réfutation :} La détermination du sexe dépend \textbf{uniquement du spermatozoïde} (X ou Y) qui féconde l'ovule (toujours X). Le père fournit le chromosome sexuel décisif. Pour les anomalies autosomales (drépanocytose, albinisme), \textbf{les deux parents transmettent un allèle} : le père et la mère sont \textbf{également responsables} de la transmission de l'allèle muté. Accuser la mère est biologiquement faux et socialement injuste.
    \item \textbf{Préjugé : « L'albinisme est une malédiction, une punition divine, ou l'enfant est un esprit/« blanc » mystique. »} \\
    \textbf{Réfutation :} L'albinisme est une \textbf{maladie génétique autosomale récessive} due à une mutation du gène \emph{TYR} (tyrosinase) ou d'autres gènes de la mélanogenèse. L'absence de mélanine est un \textbf{défaut biochimique} (\ce{Tyrosine ->[Tyrosinase] DOPA -> Melanine}). C'est une pathologie médicale (photosensibilité extrême, risque élevé de cancers cutanés, troubles visuels), \textbf{pas une entité surnaturelle}. Les personnes albinos ont les mêmes droits, la même intelligence, la même humanité.
    \item \textbf{Préjugé : « La drépanocytose / l'hémophilie / l'albinisme sont contagieux. »} \\
    \textbf{Réfutation :} Ces maladies sont \textbf{héréditaires} (transmises par les gènes lors de la fécondation), \textbf{pas infectieuses}. Il n'y a \textbf{aucun agent pathogène} (virus, bactérie) transmissible par contact, salive, sang (sauf transfusion pour hémophilie/drépanocytose, mais ce n'est pas la maladie elle-même qui se transmet, c'est le sang porteur du gène/déficit). On ne « attrape » pas la drépanocytose en serrant la main, en mangeant ensemble, ou en allant à l'école avec un enfant atteint. La stigmatisation par peur du contact est scientifiquement infondée.
\end{enumerate}
\end{propriete}

\begin{savaistu}[Le conseil génétique existe au Cameroun !]
Contrairement à une idée reçue, le conseil génétique n'est pas réservé aux pays riches. Au Cameroun, des \textbf{services de génétique médicale et de conseil génétique} existent dans les hôpitaux de référence :
\begin{itemize}
    \item \textbf{Hôpital Central de Yaoundé} (Service de Génétique / Pédiatrie / Médecine Interne).
    \item \textbf{Hôpital Laquintinie de Douala} (Service d'Hématologie pour la drépanocytose).
    \item \textbf{Centre Pasteur du Cameroun} (Diagnostic moléculaire).
    \item Certaines facultés de médecine (Yaoundé I, Douala) et laboratoires privés spécialisés.
\end{itemize}
\textbf{Le conseil génétique, c'est quoi ?} C'est une consultation où un médecin généticien ou un conseiller formé : 1) \textbf{Confirme le diagnostic} (caryotype, biologie moléculaire), 2) \textbf{Explique le mode de transmission} (schémas, risques de récidive), 3) \textbf{Informe sur les options} (diagnostic prénatal, DPI - diagnostic préimplantatoire, prise en charge médicale), 4) \textbf{Soutient la famille} sans la juger, dans le respect de ses choix. \textbf{C'est un droit pour les familles}, pas une obligation. Encouragez les familles concernées à s'y rendre.
\end{savaistu}

\courssub{Le rôle du dépistage, du conseil et de l'éducation : agir pour la santé publique}

La lutte contre les préjugés ne se fait pas seulement en cours de SVT. Elle s'appuie sur trois piliers opérationnels au Cameroun :

\begin{enumerate}
    \item \textbf{Le dépistage néonatal et prénuptial :}
    \begin{itemize}
        \item \textbf{Dépistage néonatal de la drépanocytose :} Effectué au talon (test de Guthrie ou électrophorèse hémoglobine) dès la naissance dans les grandes maternités (Yaoundé, Douala, Garoua, etc.). \textbf{Objectif :} Prise en charge \textbf{précoce} (pénicilline prophylactique, vaccins, acide folique, éducation des parents) avant les premières crises vaso-occlusives. Cela réduit drastiquement la mortalité infantile.
        \item \textbf{Dépistage prénuptial / préconceptionnel :} Connaître son statut hémoglobinique (AA, AS, SS, AC, SC...) \textbf{AVANT} de s'engager dans un projet de grossesse. Un couple AS $\times$ AS a 1/4 de risque d'enfant SS. S'ils le savent, ils peuvent choisir : diagnostic prénatal (amniocentèse), adoption, ou accepter le risque en étant préparés. \textbf{Le savoir libère le choix.}
    \end{itemize}
    \item \textbf{Le conseil génétique (voir encadré ci-dessus) :} Il transforme l'angoisse en information. Il explique que « 1/4 de risque » ne veut pas dire « 1 enfant sur 4 sera malade » (c'est un risque \textbf{par grossesse}, indépendant), et que 3/4 des enfants seront sains (ou porteurs sains).
    \item \textbf{L'éducation des communautés :} C'est le rôle de l'école, des médias, des leaders d'opinion (chefs traditionnels, religieux), des associations de patients (ex. \emph{ASDREP} - Association pour la Lutte contre la Drépanocytose au Cameroun, associations d'albinos). Il faut :
    \begin{itemize}
        \item Remplacer « enfant sorcier » par « enfant porteur d'une mutation génétique ».
        \item Expliquer que le porteur sain (AS) est \textbf{avantagé} face au paludisme (hétérozygote avantage), ce qui explique la fréquence du gène S en zone d'endémie palustre (Cameroun entier).
        \item Promouvoir l'inclusion scolaire et professionnelle : un enfant drépanocytaire bien suivi va à l'école ; un albinos protégé du soleil (crèmes, chapeaux, lunettes, vêtements couvrants) réussit ses études.
    \end{itemize}
\end{enumerate}

\begin{exemplebox}[Situation concrète : Briser le silence dans une famille]
\textit{Contexte :} Dans un village du Nord, un couple a déjà deux enfants SS (drépanocytose). La belle-famille accuse la femme d'avoir « mangé » les enfants (sorcellerie). Le mari veut la répudier.
\textit{Intervention scientifique :}
\begin{enumerate}
    \item \textbf{Explication biologique :} « Papa et Maman sont tous les deux AS (porteurs sains). À chaque grossesse, la nature fait un tirage au sort : 1 chance sur 4 d'avoir un enfant SS, 1 chance sur 2 AS, 1 chance sur 4 AA. Ce n'est la faute de personne, c'est la loi du hasard (Mendel). »
    \item \textbf{Dépistage des frères/sœurs :} Proposer le test aux autres membres de la famille élargie pour identifier les porteurs AS.
    \item \textbf{Conseil génétique :} Orienter le couple vers l'hôpital de district ou régional pour une consultation. Discuter du diagnostic prénatal pour une future grossesse.
    \item \textbf{Médiation sociale :} Impliquer le chef de village et le pasteur/imam formés aux bases de la génétique pour apaiser les tensions.
\end{enumerate}
\textbf{Résultat attendu :} La femme n'est plus coupable. Le couple est informé. Les prochains enfants (si souhaités) bénéficieront d'un suivi. La communauté apprend que la génétique n'est pas de la sorcellerie.
\end{exemplebox}

\courssub{Exercices de synthèse pour le Probatoire}

\begin{exercice}[Entraînement Probatoire — Cas mixte]
Un couple consulte pour un projet de grossesse. L'homme est daltonien. La femme est phénotypiquement normale, mais son père est daltonien.
\begin{enumerate}
    \item Préciser le mode de transmission du daltonisme. Justifier.
    \item Déterminer les génotypes du couple (allèles : $X^D$ = vision normale, $X^d$ = daltonisme).
    \item Réaliser le croisement et donner les proportions phénotypiques attendues chez la descendance (séparer filles/garçons).
    \item Calculer le risque (en \%) d'avoir un enfant daltonien à chaque grossesse.
    \item La femme dit : « Si on a une fille, elle sera daltonienne comme son grand-père. » Vrai ou Faux ? Justifier scientifiquement.
\end{enumerate}
\end{exercice}

\begin{corrige}[Exercice 1]
\begin{enumerate}
    \item \textbf{Hérédité liée au sexe, récessive.} Justification : Atteint quasi-exclusivement les hommes (père daltonien, grand-père maternel daltonien) ; pas de transmission père $\to$ fils (père daltonien $\to$ fils ? pas d'info, mais grand-père $\to$ fille $\to$ petit-fils possible).
    \item \textbf{Homme (daltonien) :} $X^dY$. \textbf{Femme (normale, père $X^dY$) :} Elle a reçu $X^d$ de son père. Elle est phénotypiquement normale $\to$ génotype $X^DX^d$ (porteuse).
    \item \textbf{Croisement $X^DX^d \times X^dY$ :}
    \begin{center}
    \begin{tabular}{|c|c|c|}\hline
     & $X^d$ & $Y$ \\ \hline
    $X^D$ & $X^DX^d$ (♀ porteuse) & $X^DY$ (♂ normal) \\ \hline
    $X^d$ & $X^dX^d$ (\textbf{♀ daltonienne}) & $X^dY$ (\textbf{♂ daltonien}) \\ \hline
    \end{tabular}
    \end{center}
    Proportions : Filles : 1/2 normales (porteuses), 1/2 daltoniennes. Garçons : 1/2 normaux, 1/2 daltoniens.
    \item Risque enfant daltonien = 1/2 filles daltoniennes $\times$ 1/2 (proba fille) + 1/2 garçons daltoniens $\times$ 1/2 (proba garçon) = 1/4 + 1/4 = \textbf{1/2 (50\%)}. \textit{Attention : 1/2 des enfants au total, mais 100\% des filles atteintes si on conditionne sur le sexe fille !}
    \item \textbf{Faux.} La fille reçoit $X^d$ de son père (obligatoirement) et $X^D$ ou $X^d$ de sa mère. Elle a 1/2 chance d'être $X^DX^d$ (porteuse normale) et 1/2 chance d'être $X^dX^d$ (daltonienne). Elle n'est pas \textbf{sûrement} daltonienne.
\end{enumerate}
\end{corrige}

\begin{exercice}[Analyse de caryotype et préjugés]
On présente le caryotype suivant : 47 chromosomes, dont deux gonosomes X et un gonosome Y, et trois chromosomes 21.
\begin{enumerate}
    \item Écrire la formule chromosomique normalisée (notation ISCN).
    \item Déterminer le sexe et l'anomalie chromosomique.
    \item Un voisin affirme : « C'est la mère qui a fait une mauvaise grossesse, elle a trop porté de charges. » Répondre scientifiquement à cette affirmation en deux arguments.
    \item Cette anomalie est-elle transmissible aux descendants de cet individu ? Justifier.
\end{enumerate}
\end{exercice}

\begin{corrige}[Exercice 2]
\begin{enumerate}
    \item \textbf{47,XY,+21}
    \item Sexe : \textbf{Masculin} (XY). Anomalie : \textbf{Trisomie 21 libre} (non-disjonction lors de la méiose I ou II, le plus souvent maternelle).
    \item \textbf{Arguments :} 1) L'origine est un \textbf{accident méiotique aléatoire} (non-disjonction) survenu lors de la formation des gamètes (ovocyte ou spermatozoïde), \textbf{bien avant la grossesse}. Les efforts physiques pendant la grossesse n'altèrent pas le nombre de chromosomes déjà fixé à la fécondation. 2) L'âge maternel avancé est un \textbf{facteur de risque statistique} (cohésion des chromatides affaiblie), mais ce n'est pas une « faute » ni une conséquence du mode de vie. C'est un aléa biologique.
    \item \textbf{Généralement non.} La trisomie 21 libre n'est pas héréditaire au sens mendélien. Le risque de récidive pour les parents est faible ($\approx$ 1\% + risque lié à l'âge maternel). \textit{Exception rare :} si l'un des parents est porteur d'une \textbf{translocation robertsonienne équilibrée} (ex. 14;21), le risque est alors élevé (10-15\%). Un caryotype des parents est recommandé pour le confirmer.
\end{enumerate}
\end{corrige}

\begin{aretenir}[Check-list finale pour le Probatoire - Leçon NA02]
\begin{itemize}
    \item \textbf{Caryotype :} Savoir compter, sexer (XX/XY), repérer trisomie/monosomie, écrire formule ($47,XX,+21$).
    \item \textbf{Arbre généalogique :} Différencier Autosomale (2 sexes, père$\to$fils OK) vs Liée au sexe (♂ majoritaires, pas père$\to$fils, saut génération).
    \item \textbf{Croisements :} Maîtriser les 3 cas types (Aa$\times$Aa ; $X^AX^a \times X^AY$ ; $X^AX^A \times X^aY$). Écrire gamètes, échiquier, proportions \textbf{par sexe si lié au sexe}.
    \item \textbf{Vocabulaire précis :} Allèle, gène, locus, homozygote, hétérozygote, hémizygote, porteur sain, phénotype, génotype, non-disjonction.
    \item \textbf{Argumentaire anti-préjugés :} Sexe déterminé par le père ; Albinisme = mutation gène \emph{TYR} (pas sorcellerie) ; Héréditaire $\neq$ Contagieux ; Dépistage/Conseil = prévention, pas fatalité.
\end{itemize}
\end{aretenir}

\courssub{Conclusion : La science au service de l'humain}

Cette leçon sur la génétique humaine n'est pas qu'un chapitre à apprendre pour un examen. C'est une \textbf{formation à la pensée critique} et à la responsabilité citoyenne. Comprendre que le caryotype révèle l'architecture de nos chromosomes, que la méiose explique le brassage et les accidents, que les lois de Mendel (autosomales ou liées au sexe) régissent la transmission des allèles, vous donne le pouvoir de \textbf{dire non à l'ignorance}.

Quand vous entendrez : « Cet enfant albinos est un esprit », vous répondrez : « Non, c'est un enfant homozygous pour une mutation du gène \emph{TYR}, il a besoin de crème solaire et de lunettes, pas de rejet. »
Quand vous entendrez : « Cette femme a fait des enfants drépanocytaires, elle est maudite », vous répondrez : « Non, le père et la mère sont tous deux AS. C'est un risque de 1/4 par grossesse, connu, dépistable, gérable. »
Quand vous verrez un caryotype 47,XY,+21, vous ne verrez pas une « erreur de la nature », mais un \textbf{accident méiotique aléatoire} dont la prise en charge médicale et éducative permet à l'enfant de s'épanouir.

\emph{La génétique n'est pas une fatalité qui enferme, c'est une connaissance qui libère.} C'est tout le sens de votre formation en SVT. Bon courage pour la suite du programme, et surtout pour le Probatoire !

\newpage

\coursec{Activité d'intégration}

\coursec{Génétique et hérédité humaines : éradication des préjugés autour des anomalies}

\courssub{Activité d'intégration : Enquête au cœur des préjugés — Conseillers génétiques du quartier}

\courssub{Objectifs de l'activité}

À l'issue de cette activité d'intégration, tu seras capable de :
\begin{itemize}
\item \cle{identifier} un caryotype humain, déterminer le sexe de l'individu et repérer une anomalie du nombre de chromosomes ;
\item \cle{écrire} correctement une formule chromosomique (normale ou anormale) ;
\item \cle{interpréter} un arbre généalogique simple et déterminer les génotypes des individus par un échiquier de croisement ;
\item \cle{calculer} les risques de transmission d'une anomalie génique (drépanocytose) à la descendance ;
\item \cle{argumenter} scientifiquement contre les préjugés sociaux qui accablent les personnes atteintes d'anomalies héréditaires (sorcellerie, « faute » de la mère, malédiction familiale, infidélité supposée).
\end{itemize}

\courssub{Situation déclenchante}

Tu es élève de Première A au Lycée de Nkongsamba. Dans le cadre de la semaine de la santé organisée par le club de sciences, le médecin du Centre Médical d'Arrondissement (CMA) a sollicité ta classe : il reçoit régulièrement des familles bouleversées par la naissance d'un enfant « différent » ou par la répétition d'une maladie dans la famille, et les rumeurs du quartier (sorcellerie, infidélité, malédiction) détruisent des couples et isolent des enfants. Il propose à ta classe de jouer le rôle de \textbf{conseillers génétiques du quartier} et de traiter trois dossiers anonymisés.

\textbf{Dossier 1 — Le nouveau-né de Bonabéri.} Une jeune mère accouche d'un garçon présentant un visage arrondi, les yeux bridés et une faiblesse musculaire (hypotonie). L'analyse des chromosomes de ses cellules (caryotype) révèle \textbf{47 chromosomes}, dont \textbf{trois exemplaires du chromosome 21} et une paire de gonosomes \textbf{XY}. La grand-mère paternelle accuse la mère d'avoir « mal porté la grossesse ».

\textbf{Dossier 2 — La famille Essomba.} Dans cette famille, le père (I-1) et la mère (I-2) voient tous deux normalement les couleurs. Ils ont quatre enfants : deux filles (II-1, II-3) à la vision normale et deux garçons : II-2 daltonien et II-4 daltonien. Le grand-père maternel (I-3, père de I-2) était daltonien. Le voisinage prétend que « le père n'est pas le vrai père ».

\textbf{Dossier 3 — Le couple de fiancés.} Aïcha et Moussa, tous deux en bonne santé, souhaitent se marier. Le test d'électrophorèse de l'hémoglobine réalisé au CMA montre qu'ils sont tous les deux de génotype \textbf{AS} (hétérozygotes pour le gène de l'hémoglobine ; S désigne l'allèle drépanocytaire). Leurs familles hésitent : « est-ce dangereux de les laisser se marier ? »

\courssub{Consignes}

Travail en groupes de quatre élèves. Chaque groupe traite un dossier, puis présente ses résultats au tableau devant la classe.

\textbf{Pour le dossier 1 :}
\begin{enumerate}
\item Décris le caryotype attendu d'un homme normal, puis compare-le à celui du nouveau-né.
\item Identifie l'anomalie chromosomique et nomme la maladie correspondante.
\item Écris la formule chromosomique du nouveau-né.
\item Rédige en cinq lignes maximum une explication rassurante et scientifique destinée aux parents, en précisant l'origine de l'anomalie.
\end{enumerate}

\textbf{Pour le dossier 2 :}
\begin{enumerate}
\item Rappelle sur quel chromosome est porté l'allèle responsable du daltonisme et précise s'il est dominant ou récessif.
\item Écris les génotypes possibles du père I-1 et de la mère I-2, sachant qu'ils ont des fils daltoniens.
\item Construis l'échiquier de croisement du couple I-1 $\times$ I-2 et vérifie qu'il rend compte des quatre enfants.
\item Explique pourquoi le grand-père maternel I-3 confirme l'origine de l'allèle, et réfute la rumeur du voisinage.
\end{enumerate}

\textbf{Pour le dossier 3 :}
\begin{enumerate}
\item Rappelle le mode de transmission de la drépanocytose (autosomique récessif).
\item Construis l'échiquier de croisement AS $\times$ AS et calcule les proportions génotypiques et phénotypiques de la descendance.
\item Rédige un court avis de conseil génétique (dix lignes maximum) : le mariage est-il interdit ? Quels sont les risques réels ? Quelles précautions proposer ?
\end{enumerate}

\textbf{Travail collectif final :} la classe rédige ensemble un \textbf{tract anti-préjugés} comportant cinq affirmations courtes, chacune appuyée par un argument scientifique tiré des trois dossiers.

\courssub{Ressources à mobiliser}

\begin{aretenir}[Boîte à outils du conseiller génétique]
\begin{itemize}
\item Le caryotype humain normal comporte \cle{46 chromosomes}, soit \cle{23 paires} : 22 paires d'\cle{autosomes} (numérotées de 1 à 22 par taille décroissante) et 1 paire de \cle{gonosomes} (XX chez la femme, XY chez l'homme).
\item Formules chromosomiques normales : femme $46,XX$ ; homme $46,XY$. Une anomalie du nombre s'écrit en indiquant le nombre total puis les gonosomes et l'anomalie : trisomie 21 chez un garçon $= 47,XY,+21$ ; monosomie X (syndrome de Turner) $= 45,X$.
\item Le daltonisme et l'hémophilie sont dus à des allèles \cle{récessifs portés par le chromosome X} (hérédité liée au sexe / hétérochromosomique) : un garçon ($XY$) est atteint dès qu'il reçoit l'allèle muté de sa mère ; une fille ne l'est que si elle reçoit l'allèle muté des deux parents (rare). Une femme hétérozygote est une \cle{conductrice} (porteuse saine).
\item L'albinisme et la drépanocytose sont des anomalies \cle{autosomiques récessives} : le gène est sur un autosome ; un individu sain peut être \cle{porteur sain} (hétérozygote, ex. $AS$ pour la drépanocytose). Seuls les homozygotes récessifs ($aa$ ou $SS$) sont atteints.
\item Échiquier de croisement : gamètes d'un parent en ligne, de l'autre en colonne ; chaque case = un génotype possible à probabilité égale ($1/4$ pour un croisement dihybride ou $Aa \times Aa$, $1/2$ pour les croisements test).
\end{itemize}
\end{aretenir}

\courssub{Démarche et corrigé commenté}

\begin{exoresolu}[Corrigé commenté des trois dossiers]

\textbf{Dossier 1 — Le nouveau-né de Bonabéri : lecture de caryotype et formule chromosomique}

\begin{popfigure}
\begin{tikzpicture}[scale=0.7, every node/.style={font=\small}]
% Normal Male Karyotype
\node[align=center] at (0,5.5) {\textbf{Homme normal (46,XY)}};
\foreach \i in {1,...,22} {
    \pgfmathsetmacro{\row}{int((\i-1)/7)}
    \pgfmathsetmacro{\col}{int(mod(\i-1,7))}
    \draw[popblue, thick] (1.5*\col, 4.5-\row*1.2) rectangle ++(1.2, 1.0);
    \node at (1.5*\col+0.6, 5.0-\row*1.2) {\i};
}
% Sex chromosomes
\draw[poporange, thick] (10.5, 4.5) rectangle ++(1.2, 1.0); \node at (11.1, 5.0) {X};
\draw[poporange, thick] (11.8, 4.5) rectangle ++(1.2, 1.0); \node at (12.4, 5.0) {Y};

% Trisomy 21 Karyotype
\node[align=center] at (0,0.5) {\textbf{Trisomie 21 (47,XY,+21)}};
\foreach \i in {1,...,22} {
    \pgfmathsetmacro{\row}{int((\i-1)/7)}
    \pgfmathsetmacro{\col}{int(mod(\i-1,7))}
    \draw[popblue, thick] (1.5*\col, -0.5-\row*1.2) rectangle ++(1.2, 1.0);
    \node at (1.5*\col+0.6, 0.0-\row*1.2) {\i};
}
% Extra chromosome 21
\draw[popred, thick, dashed] (10.5, -0.5) rectangle ++(1.2, 1.0); \node[popred] at (11.1, 0.0) {21*};
% Sex chromosomes
\draw[poporange, thick] (11.8, -0.5) rectangle ++(1.2, 1.0); \node at (12.4, 0.0) {X};
\draw[poporange, thick] (13.1, -0.5) rectangle ++(1.2, 1.0); \node at (13.7, 0.0) {Y};
\end{tikzpicture}
\legende{Comparaison schématique des caryotypes : homme normal (46 chromosomes) vs trisomie 21 (47 chromosomes avec un chromosome 21 surnuméraire en rouge pointillé). Les autosomes sont en bleu, les gonosomes en orange.}
\end{popfigure}

1. Un homme normal possède \textbf{46 chromosomes} organisés en 23 paires : \textbf{22 paires d'autosomes} (homologues deux à deux) et \textbf{1 paire de gonosomes XY} (hétérologues). Le nouveau-né possède \textbf{47 chromosomes} : il y a \textbf{un chromosome en plus} par rapport à la normale.

2. L'analyse montre \textbf{trois exemplaires du chromosome 21} (au lieu de deux). Il s'agit d'une \textbf{trisomie 21} (syndrome de Down). C'est une anomalie du \emph{nombre} de chromosomes (anomalie numérique), due le plus souvent à une \textbf{non-disjonction} de la paire de chromosomes 21 lors de la méiose (généralement la méiose I ou II de l'oogenèse, plus rarement la spermatogenèse) : un gamète se retrouve avec deux chromosomes 21 (disomie 21) ; la fécondation par un gamète normal (haploïde, 1 chromosome 21) donne un zygote à 47 chromosomes (trisomie 21). \textbf{Cet accident méiotique survient au hasard, avant la conception.}

3. Formule chromosomique du nouveau-né (notation ISCN) :
\[ 47,XY,+21 \]
(47 chromosomes au total, gonosomes XY $\rightarrow$ sexe masculin, $+21$ $\rightarrow$ un chromosome 21 en excès).

4. \textbf{Message aux parents (exemple) :}
« Votre enfant présente une trisomie 21 : il possède un chromosome 21 en plus dans toutes ses cellules. Cette anomalie chromosomique survient par hasard, au moment de la formation des gamètes (ovule ou spermatozoïde), \textbf{bien avant la conception et l'implantation}. \textbf{Personne n'en est responsable} : ni la mère, ni le père, ni aucune pratique pendant la grossesse. Ce n'est ni une sorcellerie, ni une malédiction, ni une punition. Avec un suivi médical régulier (cardiologie, ORL, ophtalmologie), une stimulation précoce, une scolarisation adaptée et beaucoup d'affection, votre enfant peut se développer, apprendre, acquérir une autonomie et vivre heureux. »

\tcblower

\textbf{Dossier 2 — La famille Essomba : hérédité liée au sexe (daltonisme)}

\begin{popfigure}
\begin{tikzpicture}[scale=0.9, every node/.style={font=\small}, node distance=1.2cm and 1.5cm]
% Generation I
\node[draw, rectangle, minimum size=0.8cm] (I1) {I-1};
\node[draw, circle, minimum size=0.8cm, right=of I1] (I2) {I-2};
\node[draw, rectangle, minimum size=0.8cm, fill=popred!30, right=of I2] (I3) {I-3};
\node[above=0.2cm of I1] {\textbf{Père}};
\node[above=0.2cm of I2] {\textbf{Mère}};
\node[above=0.2cm of I3] {\textbf{Grand-père maternel}};
\node[below=0.2cm of I1] {Sain ($X^D Y$)};
\node[below=0.2cm of I2] {Saine conductrice ($X^D X^d$)};
\node[below=0.2cm of I3] {Daltonien ($X^d Y$)};

% Marriage line
\draw (I1) -- (I2);
\draw (I2) -- (I3); % Not a marriage line, just lineage. Better: vertical line from I3 down to a connector.
% Actually standard pedigree: I-3 above I-2.
\end{tikzpicture}
\legende{Arbre généalogique simplifié de la famille Essomba. Carrés = hommes, Cercles = femmes. Hachuré/coloré = phénotype daltonien. I-3 (grand-père maternel) est daltonien ($X^d Y$). I-2 (mère) est conductrice ($X^D X^d$). I-1 (père) est sain ($X^D Y$).}
\end{popfigure}

1. L'allèle responsable du daltonisme (vision anormale des couleurs rouge/vert) est porté par le \textbf{chromosome X} : c'est une hérédité \textbf{liée au sexe (hétérochromosomique)}. Cet allèle est \textbf{récessif} par rapport à l'allèle de la vision normale. Notons : $X^D$ = chromosome X porteur de l'allèle normal (dominant) ; $X^d$ = chromosome X porteur de l'allèle daltonien (récessif).

2. \begin{itemize}
    \item Père I-1 : phénotype « vision normale ». Comme il n'a qu'un seul chromosome X, son génotype est nécessairement \textbf{$X^D Y$}.
    \item Mère I-2 : phénotype « vision normale » mais elle a des fils daltoniens. Un fils reçoit son unique chromosome X \textbf{obligatoirement de sa mère}. Pour avoir un fils $X^d Y$, la mère doit posséder l'allèle $d$. Comme elle est phénotypiquement normale, elle est hétérozygote : \textbf{$X^D X^d$} (conductrice / porteuse saine).
\end{itemize}

3. \textbf{Échiquier de croisement $X^D Y \times X^D X^d$ :}

\begin{center}
\begin{tabular}{|c|c|c|}
\hline
\rowcolor{popblueL} \textbf{Gamètes maternels $\downarrow$ / Gamètes paternels $\rightarrow$} & \textbf{$X^D$} & \textbf{$Y$} \\
\hline
\textbf{$X^D$} & $X^D X^D$ \newline Fille saine non conductrice & $X^D Y$ \newline Garçon sain \\
\hline
\textbf{$X^d$} & $X^D X^d$ \newline Fille saine \textbf{conductrice} & $X^d Y$ \newline \textbf{Garçon daltonien} \\
\hline
\end{tabular}
\end{center}

\textbf{Interprétation :} Probabilités à chaque conception : $1/4$ fille $X^D X^D$ (saine), $1/4$ fille $X^D X^d$ (saine conductrice), $1/4$ garçon $X^D Y$ (sain), $1/4$ garçon $X^d Y$ (daltonien).
La famille observée (2 filles saines, 2 garçons daltoniens) est compatible avec ces probabilités (échantillon de 4 enfants, fluctuations possibles autour des 1/4 théoriques).

4. Le grand-père maternel I-3 est daltonien : son génotype est \textbf{$X^d Y$}. Il transmet son unique chromosome X \textbf{à toutes ses filles}. Sa fille I-2 a donc \textbf{obligatoirement} reçu le chromosome $X^d$. Elle est ainsi devenue conductrice ($X^D X^d$) et a pu transmettre $X^d$ à ses propres fils (II-2 et II-4).
\textbf{Conclusion scientifique :} L'allèle daltonien vient de la lignée maternelle (grand-père $\rightarrow$ mère $\rightarrow$ fils). Le père I-1 transmet son chromosome Y à ses fils ; il ne peut \textbf{pas} leur transmettre l'allèle daltonien situé sur le X. La rumeur d'infidélité est \textbf{scientifiquement infondée}.

\tcblower

\textbf{Dossier 3 — Le couple Aïcha et Moussa : hérédité autosomique récessive (drépanocytose)}

1. La drépanocytose (maladie des globules rouges falciformes) est une anomalie \textbf{autosomique récessive}. Le gène concerné ($\beta$-globine) est sur le \textbf{chromosome 11} (autosome).
\begin{itemize}
    \item Allèle $A$ : hémoglobine normale (HbA), dominant.
    \item Allèle $S$ : hémoglobine drépanocytaire (HbS), récessif.
    \item Génotypes : $AA$ = sain non porteur ; $AS$ = \textbf{porteur sain} (hétérozygote, asymptomatique en conditions normales, résistant au paludisme grave) ; $SS$ = \textbf{drépanocytaire} (malade, anémie hémolytique chronique, crises vaso-occlusives).
\end{itemize}
Aïcha et Moussa sont tous deux $AS$ : ils sont \textbf{porteurs sains}, cliniquement indemnes.

2. \textbf{Échiquier de croisement $AS \times AS$ :}

\begin{center}
\begin{tabular}{|c|c|c|}
\hline
\rowcolor{popblueL} \textbf{Gamètes $\downarrow$ / $\rightarrow$} & \textbf{$A$} & \textbf{$S$} \\
\hline
\textbf{$A$} & $AA$ \newline Sain non porteur & $AS$ \newline Porteur sain \\
\hline
\textbf{$S$} & $AS$ \newline Porteur sain & $SS$ \newline \textbf{Drépanocytaire} \\
\hline
\end{tabular}
\end{center}

\textbf{Proportions théoriques à chaque grossesse (loi de Mendel) :}
\begin{itemize}
    \item $1/4$ ($25\%$) $AA$ : enfant sain, non porteur.
    \item $1/2$ ($50\%$) $AS$ : enfant porteur sain (comme les parents).
    \item $1/4$ ($25\%$) $SS$ : enfant atteint de drépanocytose majeure.
\end{itemize}
Le risque d'avoir un enfant malade est de \textbf{25 \% à chaque grossesse}, indépendamment des issues des grossesses précédentes (indépendance des événements).

3. \textbf{Avis de conseil génétique (exemple) :}
« Aïcha et Moussa sont tous deux en parfaite santé. Leur mariage n'a \textbf{aucune raison d'être interdit} sur le plan médical ou génétique. En tant que couple $AS \times AS$, ils ont, \textbf{à chaque grossesse}, une probabilité de 25 \% d'avoir un enfant drépanocytaire ($SS$), 50 \% d'avoir un enfant porteur sain ($AS$) et 25 \% d'avoir un enfant non porteur ($AA$). Nous leur recommandons : (1) une information claire sur ce risque \textbf{avant} la première grossesse ; (2) un suivi prénatal précoce avec possibilité de diagnostic prénatal (biopsie des villosités choriales ou amniocentèse) pour connaître le génotype fœtal ; (3) en cas de naissance d'un enfant $SS$, une prise en charge immédiate et continue au CMA (prophylaxie anti-paludique, vaccination pneumocoque/Haemophilus, acide folique, hydratation, éducation thérapeutique) qui transforme le pronostic vital. Rappel : être $AS$ confère un avantage sélectif contre le paludisme grave ; ce n'est pas une maladie. »

\end{exoresolu}

\begin{attention}[Pièges à éviter dans les présentations et au Probatoire]
\begin{itemize}
\item \textbf{Formule chromosomique :} Ne jamais écrire « $47,XXY$ » pour une trisomie 21 (cela désigne le syndrome de Klinefelter). Le chromosome surnuméraire est un \emph{autosome} (le 21) : on écrit \textbf{$47,XY,+21$} (garçon) ou \textbf{$47,XX,+21$} (fille).
\item \textbf{Hérédité liée au sexe :} Un père \textbf{ne transmet jamais son chromosome X à ses fils} (il leur donne son Y). Dire « le fils est daltonien à cause de son père » ou « le père a transmis l'allèle » est une \textbf{erreur éliminatoire}.
\item \textbf{Probabilités :} « 25 \% de risque » ne signifie \textbf{pas} « un enfant malade sur quatre garanti » dans la fratrie. Chaque grossesse est un tirage indépendant : une famille de 4 enfants peut avoir 0, 1, 2, 3 ou 4 enfants atteints.
\item \textbf{Porteur vs Malade :} Un porteur sain ($AS$, ou femme conductrice $X^D X^d$) n'est \textbf{pas malade}. Ne jamais confondre « porteur » (hétérozygote asymptomatique) et « atteint » (homozygote récessif ou hémizygote récessif).
\item \textbf{Non-disjonction :} Préciser qu'elle a lieu lors de la \textbf{méiose} (formation des gamètes), pas lors de la mitose (sauf mosaïcisme, hors programme). Ne pas dire « pendant la grossesse ».
\end{itemize}
\end{attention}

\courssub{Bilan de l'activité}

Cette activité t'a permis de mobiliser ensemble tous les savoir-faire de la leçon dans des situations authentiques et socialement pertinentes :
\begin{itemize}
\item Lecture d'un caryotype, détermination du sexe, identification d'une anomalie numérique et écriture de la formule chromosomique normalisée (Dossier 1).
\item Interprétation d'un arbre généalogique, détermination des génotypes, construction et exploitation d'un échiquier de croisement pour une hérédité liée au sexe (Dossier 2).
\item Calcul de risques mendéliens pour une hérédité autosomique récessive et rédaction d'un avis de conseil génétique (Dossier 3).
\end{itemize}
Elle a surtout montré que la génétique est un outil de \cle{justice sociale} et de \cle{santé publique} : elle innocente les mères accusées à tort, réfute les rumeurs d'infidélité, replace la responsabilité partagée au niveau du couple (pour l'autosomique) et permet des choix reproductifs éclairés grâce au dépistage prénuptial et prénatal. La science protège, le préjugé exclut.

\begin{exemplebox}[Exemple de tract anti-préjugés rédigé par la classe (à afficher au CMA et au lycée)]
\begin{enumerate}
\item « La trisomie 21 est due à un chromosome 21 en plus survenu par hasard lors de la formation des gamètes : \textbf{aucun parent n'en est responsable}, ce n'est pas une "faute" de la grossesse. »
\item « Un garçon daltonien ou hémophile a reçu l'allèle muté de sa \textbf{mère}, qui le tenait souvent de son propre père : accuser le père d'infidélité est une erreur biologique. »
\item « Être porteur sain (AS pour la drépanocytose, conductrice pour l'hémophilie/daltonisme) n'est \textbf{pas être malade} : on vit, travaille et se marie normalement. »
\item « Un couple AS $\times$ AS a 25 \% de risque d'enfant drépanocytaire à chaque grossesse : le dépistage avant mariage permet de \textbf{choisir en connaissance de cause} (suivi, diagnostic prénatal), pas d'interdire l'union. »
\item « L'albinisme, la drépanocytose, le daltonisme, la trisomie s'expliquent par les chromosomes et les gènes : \textbf{la science éclaire, le préjugé détruit}. »
\end{enumerate}
\end{exemplebox}

\begin{savaistu}[Le dépistage et la prise en charge de la drépanocytose au Cameroun]
La drépanocytose est la maladie génétique la plus fréquente au Cameroun. On estime qu'environ \textbf{20 à 25 \% de la population} est porteuse saine (génotype $AS$) et qu'environ \textbf{2 à 3 \% des nouveau-nés} sont atteints (génotype $SS$), soit plusieurs milliers de naissances par an. L'allèle $S$ s'est maintenu à haute fréquence dans les zones d'endémie palustre parce que les hétérozygotes $AS$ résistent mieux au paludisme grave (avantage sélectif de l'hétérozygote / \emph{balanced polymorphism}).

Depuis 2019, le Cameroun a intégré le \textbf{dépistage néonatal systématique} de la drépanocytose dans le Programme Élargi de Vaccination (PEV) dans plusieurs régions (Centre, Littoral, Nord-Ouest, Sud-Ouest) : une goutte de sang sur papier buvard (test de Guthrie) permet de détecter les $SS$ dès la naissance pour une prise en charge précoce (pénicilline prophylactique, vaccins, acide folique, éducation des parents). Des \textbf{consultations de conseil génétique} existent notamment à l'Hôpital Central de Yaoundé, à l'Hôpital Général de Douala, à l'Hôpital Laquintinie de Douala et dans les centres de référence régionaux. N'hésite pas à t'y renseigner ou à orienter ton entourage : \textbf{connaître son statut (AA, AS, SS) avant le mariage est un acte de responsabilité et d'amour.}
\end{savaistu}

\newpage

\coursec{Méthodes et exercices résolus}

\coursec{Génétique et hérédité humaines : éradication des préjugés autour des anomalies}

\courssub{Le caryotype humain : notion, caractéristiques et identification}

\begin{definition}[Caryotype humain]
Le \cle{caryotype} est la photographie, au microscope optique, des chromosomes d'une cellule en \cle{métaphase} de mitose, après coloration (bande G) et classement informatisé. Il permet de visualiser le \cle{complement chromosomique} d'un individu.
\end{definition}

\begin{propriete}[Caractéristiques du caryotype humain normal]
\begin{itemize}
\item \textbf{Nombre} : \cle{46 chromosomes} soit \cle{23 paires} ($2n = 46$). L'espèce humaine est \cle{diploïde} ($2n$).
\item \textbf{Autosomes} : 22 paires (paires 1 à 22), identiques chez l'homme et la femme. Elles sont classées par taille décroissante (groupe A à G selon la classification de Denver).
\item \textbf{Gonosomes (chromosomes sexuels)} : 23\ieme{} paire (notation $23^{\text{ème}}$).
  \begin{itemize}
  \item Femelle : deux chromosomes X homologues $\rightarrow$ \cle{XX} (homogamétique).
  \item Mâle : un chromosome X et un chromosome Y non homologues (hétéromorphes) $\rightarrow$ \cle{XY} (hétérogamétique).
  \end{itemize}
\end{itemize}
\end{propriete}

\begin{popfigure}
\begin{tikzpicture}[scale=0.85, every node/.style={font=\small}]
% Dessiner les paires 1-22 schématiquement (groupes)
% Groupe A (1-3) : grands métacentriques
\foreach \i/\y in {1/11.5, 2/10.5, 3/9.5} {
  \draw[popdark, thick] (0.5,\y) -- (0.5,\y-0.7) (1.5,\y) -- (1.5,\y-0.7);
  \draw[popdark, thick] (0.5,\y-0.35) -- (1.5,\y-0.35); % centromère
  \node at (2.2, \y-0.35) {Paire \i};
}
% Groupe B (4-5) : grands sous-métacentriques
\foreach \i/\y in {4/8.5, 5/7.5} {
  \draw[popdark, thick] (0.5,\y) -- (0.5,\y-0.65) (1.5,\y) -- (1.5,\y-0.65);
  \draw[popdark, thick] (0.5,\y-0.2) -- (1.5,\y-0.2);
  \node at (2.2, \y-0.325) {Paire \i};
}
% ... Groupes C à G (représentés par une zone pointillée pour compacité)
\draw[popmuted, dashed] (0,7) rectangle (2, -0.5);
\node[popmuted, align=center] at (1, 3.2) {Paires 6 à 12\\(Groupe C)};
\node[popmuted, align=center] at (1, 1.5) {Paires 13 à 22\\(Groupes D, E, F, G)};
% Gonosomes (Paire 23)
\draw[popblue, thick] (0.5,-0.5) -- (0.5,-1.3) (1.5,-0.5) -- (1.5,-1.3); % X
\draw[popblue, thick] (0.5,-0.9) -- (1.5,-0.9); % centromère X
\draw[poporange, thick] (2.5,-0.5) -- (2.5,-1.1) (3.0,-0.5) -- (3.0,-1.1); % Y (acrocentrique)
\draw[poporange, thick] (2.5,-0.7) -- (3.0,-0.7); % centromère Y
\node at (1.0, -1.5) {X};
\node at (2.75, -1.5) {Y};
\node[popblue, align=center] at (1.75, -0.9) {Paire 23\\(Gonosomes XY)};
% Légende
\node[draw, popblueL, fill=popcream, rounded corners, anchor=south west, align=left] at (4, 10) {
\textbf{Légende :}\\
Chromosomes appariés par taille et position du centromère.\\
\color{popblue} Bleu : Autosomes (paires 1-22).\\
\color{poporange} Orange : Gonosomes (XY = Mâle).\\
Centromère = trait transversal.
};
\end{tikzpicture}
\legende{Caryotype humain masculin normal ($46, XY$) : représentation schématique des 23 paires de chromosomes classées par taille décroissante. Les autosomes (paires 1 à 22) sont homologues deux à deux. La 23\ieme{} paire (gonosomes) est hétéromorphe : un grand chromosome X et un petit chromosome Y.}
\end{popfigure}

\begin{methode}[Décrire et identifier un caryotype humain]
Pour analyser un caryotype (photographie des chromosomes d'une cellule en métaphase, rangés par paires) :
\begin{enumerate}
\item \textbf{Compter} le nombre total de chromosomes : chez l'Homme sain, on doit trouver \cle{46 chromosomes}, soit \cle{23 paires} ($2n = 46$).
\item \textbf{Classer} les chromosomes par paires d'\cle{homologues} (même taille, même forme, même position du centromère, même bandeamento) et par ordre de taille décroissante : paires 1 à 22 = \cle{autosomes} ; 23\ieme{} paire = \cle{gonosomes} (chromosomes sexuels X et Y).
\item \textbf{Examiner la 23\ieme{} paire} : deux chromosomes identiques de grande taille = $XX$ (femelle) ; un grand chromosome et un très petit = $XY$ (mâle).
\item \textbf{Vérifier chaque paire} : toute paire à 3 chromosomes (trisomie) ou à 1 seul chromosome (monosomie) signale une anomalie du nombre. Une structure anormale (délétion, translocation, inversion) signale une anomalie de structure.
\item \textbf{Conclure} en rédigeant la \cle{formule chromosomique} : nombre total, virgule, gonosomes, puis mention de l'anomalie le cas échéant. Exemples : $46, XX$ ; $46, XY$ ; $47, XY, +21$ ; $45, X$.
\end{enumerate}
\textbf{Réflexe Probatoire} : toujours justifier l'identification par deux arguments chiffrés (nombre total + nature de la paire sexuelle) et décrire la morphologie des chromosomes (métacentrique, sous-métacentrique, acrocentrique).
\end{methode}

\begin{exoresolu}[Analyse d'un caryotype à l'Hôpital Central de Yaoundé]
Au laboratoire de cytogénétique, on réalise le caryotype des cellules d'un patient. On observe 46 chromosomes répartis en 23 paires ; les paires 1 à 22 sont constituées de chromosomes deux à deux identiques (même taille, même centromère) ; la 23\ieme{} paire comporte un grand chromosome (métacentrique) et un chromosome très petit (acrocentrique).
\begin{enumerate}
\item Établis la formule chromosomique de cet individu.
\item Précise son sexe en justifiant.
\item Ce caryotype est-il normal ?
\end{enumerate}
\tcblower
\textbf{1. Formule chromosomique.}
\begin{itemize}
\item Dénombrement : 46 chromosomes au total, soit 23 paires $\Rightarrow$ le nombre est conforme à la formule humaine $2n = 46$.
\item Paires 1 à 22 : chromosomes homologues deux à deux semblables $\Rightarrow$ ce sont les 22 paires d'autosomes, sans anomalie numérique ni structurelle visible.
\item 23\ieme{} paire : un grand chromosome (X, métacentrique) + un très petit chromosome (Y, acrocentrique) $\Rightarrow$ gonosomes $XY$.
\end{itemize}
La formule chromosomique s'écrit donc : \[ 46, XY \]
\textbf{2. Sexe de l'individu.} La 23\ieme{} paire est hétéromorphe (un X et un Y). Or le chromosome Y est spécifique du sexe masculin chez l'Homme (porteur du gène \emph{SRY} déterminant la testicule). L'individu est donc de \textbf{sexe masculin}.
\textbf{3. Normalité.} Le nombre total est de 46 chromosomes, toutes les paires d'autosomes sont complètes (2 exemplaires chacune) et la paire sexuelle est $XY$ sans surnuméraire ni délétion visible. Le caryotype est donc \textbf{normal} : il s'agit d'un homme sans anomalie chromosomique décelable à cette résolution.
\textbf{Conclusion} : formule $46, XY$ ; individu de sexe masculin ; caryotype normal.
\end{exoresolu}

\courssub{Anomalies du nombre de chromosomes et formule chromosomique}

\begin{definition}[Anomalies numériques : Trisomie et Monosomie]
Une \cle{anomalie chromosomique numérique} résulte d'une \cle{non-disjonction} chromosomique au cours de la méiose I ou II (ou plus rarement en mitose précoce), conduisant à des gamètes anormaux ($n+1$ ou $n-1$ chromosomes).
\begin{itemize}
\item \textbf{Trisomie} : présence de 3 exemplaires d'un chromosome au lieu de 2. Caryotype à $2n+1 = 47$ chromosomes. Exemple majeur : \cle{trisomie 21} (syndrome de Down) : $47, XX, +21$ ou $47, XY, +21$. Autres trisomies viables : 13 (Patau), 18 (Edwards), XXX, XXY (Klinefelter), XYY.
\item \textbf{Monosomie} : absence d'un chromosome d'une paire. Caryotype à $2n-1 = 45$ chromosomes. La monosomie autosomale est létale in utero. Seule monosomie viable : \cle{syndrome de Turner} ($45, X$ ou mosaique $45, X / 46, XX$).
\end{itemize}
\end{definition}

\begin{propriete}[Écriture rigoureuse de la formule chromosomique]
La notation internationale (ISCN) s'écrit dans l'ordre :
\[ \text{Nombre total}, \text{ Gonosomes}, \text{ Anomalies (par ordre croissant de n° de paire)} \]
Exemples :
\begin{itemize}
\item Homme normal : $46, XY$
\item Femme trisomique 21 : $47, XX, +21$
\item Homme syndrome de Klinefelter : $47, XXY$
\item Femme syndrome de Turner : $45, X$
\item Translocation robertsonienne 14;21 équilibrée : $45, XX, \text{rob}(14;21)$
\end{itemize}
\end{propriete}

\begin{methode}[Identifier une anomalie chromosomique numérique sur un caryotype]
\begin{enumerate}
\item \textbf{Compter} le nombre total de chromosomes.
\item Si $2n = 47$ : chercher la paire qui en compte 3 $\rightarrow$ \cle{trisomie}. Noter $+n$ (n = n° de la paire).
\item Si $2n = 45$ : chercher la paire qui en compte 1 (souvent les gonosomes) $\rightarrow$ \cle{monosomie}. Noter $-n$ ou simplement l'absence (ex $45, X$).
\item \textbf{Rédiger la formule complète} : total, virgule, gonosomes, virgule, anomalie.
\item \textbf{Relier au phénotype} : trisomie 21 $\rightarrow$ retard mental, traits faciaux caractéristiques, cardiopathies ; Turner $\rightarrow$ nanisme, stérilité, plis du cou ; Klinefelter $\rightarrow$ stérilité, gynécomastie, grande taille.
\end{enumerate}
\end{methode}

\begin{exoresolu}[Trisomie 21 ou mutation génique ?]
Deux enfants sont examinés dans un centre de santé de Douala. L'enfant M présente un retard mental, des yeux bridés, un pli cutané transversal de la paume et un caryotype de 47 chromosomes dont trois chromosomes de la paire 21. L'enfant N, de caryotype $46, XY$ sans particularité morphologique chromosomique, présente des crises douloureuses vaso-occlusives et une anémie hémolytique ; l'examen de son sang au microscope montre des globules rouges en forme de faucille (drepanocytes).
\begin{enumerate}
\item Établis la formule chromosomique de l'enfant M et nomme son affection.
\item L'anomalie de l'enfant N est-elle visible sur son caryotype ? Justifie et nomme la maladie.
\end{enumerate}
\tcblower
\textbf{1. Enfant M.} Le caryotype comporte 47 chromosomes au lieu de 46, avec \textbf{trois chromosomes 21} (paire 21 à 3 exemplaires) : il s'agit d'une trisomie 21 due à une non-disjonction de la paire 21 lors de la méiose (souvent maternelle, liée à l'âge). La formule chromosomique s'écrit :
\[ 47, XY, +21 \quad \text{(ou } 47, XX, +21 \text{ si c'est une fille)} \]
L'affection est la \textbf{trisomie 21} (syndrome de Down), anomalie du \emph{nombre} de chromosomes.
\textbf{2. Enfant N.} Son caryotype est normal ($46, XY$) : l'anomalie n'est donc \textbf{pas visible} sur le caryotype (les gènes sont microscopiques). Les globules rouges en faucille (drepanocytes) et l'anémie hémolytique révèlent la \textbf{drépanocytose} (anémie falciforme), maladie due à la \emph{mutation ponctuelle d'un gène} autosomal (chromosome 11) codant la chaîne $\beta$-globine de l'hémoglobine (hémoglobine normale HbA $\rightarrow$ hémoglobine anormale HbS). C'est une anomalie \emph{génique} (ou allélique).
\textbf{Conclusion} : M présente une anomalie chromosomique (trisomie 21, $47, XY, +21$) décelable sur le caryotype ; N présente une anomalie génique (drépanocytose) que seuls le phénotype clinique et l'analyse biochimique/génétique révèlent.
\end{exoresolu}

\courssub{Chromosomes et détermination du sexe}

\begin{popfigure}
\begin{tikzpicture}[scale=0.9, every node/.style={font=\small}]
% Méiose mâle
\node[draw, popblueL, fill=popblue!10, rounded corners, align=center] at (0, 5) {Spermatogonie\\$46, XY$ ($2n$)};
\draw[->, thick, popdark] (0, 4.3) -- (0, 3.5) node[midway, right] {Méiose I \& II};
\node[draw, poporangeL, fill=poporange!10, rounded corners, align=center] at (-2.5, 2.5) {Spermatozoïde\\$23, X$ ($n$) \\ 50 \%};
\node[draw, poporangeL, fill=poporange!10, rounded corners, align=center] at (2.5, 2.5) {Spermatozoïde\\$23, Y$ ($n$) \\ 50 \%};
% Méiose femelle
\node[draw, poppinkL, fill=poppink!10, rounded corners, align=center] at (7, 5) {Oogonie\\$46, XX$ ($2n$)};
\draw[->, thick, popdark] (7, 4.3) -- (7, 3.5) node[midway, left] {Méiose I \& II};
\node[draw, poppurpleL, fill=poppurple!10, rounded corners, align=center] at (7, 2.5) {Ovule\\$23, X$ ($n$) \\ 100 \%};
% Fécondation
\draw[->, thick, popgreen] (-2.5, 1.8) -- (-1, 1) node[midway, above, sloped] {Fécondation};
\draw[->, thick, popgreen] (2.5, 1.8) -- (1, 1) node[midway, above, sloped] {Fécondation};
\draw[->, thick, popgreen] (7, 1.8) -- (0, 1) node[midway, above] {Fécondation};
\node[draw, popgreenL, fill=popgreen!10, rounded corners, align=center] at (-1, 0) {Zygote\\$46, XX$\\$\rightarrow$ Fille};
\node[draw, popgreenL, fill=popgreen!10, rounded corners, align=center] at (1, 0) {Zygote\\$46, XY$\\$\rightarrow$ Garçon};
% Flèches croisées
\draw[dashed, popmuted] (-2.5, 2.5) -- (-1, 1);
\draw[dashed, popmuted] (2.5, 2.5) -- (1, 1);
\draw[dashed, popmuted] (7, 2.5) -- (-1, 1);
\draw[dashed, popmuted] (7, 2.5) -- (1, 1);
\end{tikzpicture}
\legende{Détermination chromosomique du sexe chez l'Homme. La méiose chez le mâle ($XY$) produit deux types de spermatozoïdes (hétérogamétie) : 50 \% portent un chromosome X, 50 \% un chromosome Y. La méiose chez la femelle ($XX$) ne produit qu'un seul type d'ovule (homogamétie) : 100 \% portent un chromosome X. Le sexe de l'enfant est déterminé par le spermatozoïde fécondant : $X_{sperm} + X_{ovule} \rightarrow XX$ (fille) ; $Y_{sperm} + X_{ovule} \rightarrow XY$ (garçon). Probabilité = 1/2 pour chaque sexe.}
\end{popfigure}

\begin{methode}[Exploiter les gamètes pour déterminer le sexe d'un individu ou d'un futur enfant]
\begin{enumerate}
\item \textbf{Rappel} : les \cle{gamètes} sont haploïdes ($n = 23$ chromosomes : 22 autosomes + 1 gonosome).
\item \textbf{Chez la femme} ($46, XX$) : la méiose ne produit qu'\textbf{un seul type} d'ovule, tous porteurs d'un chromosome $X$.
\item \textbf{Chez l'homme} ($46, XY$) : la méiose produit \textbf{deux types} de spermatozoïdes en proportions égales (loi de ségrégation) : 50 \% portent $X$, 50 \% portent $Y$.
\item \textbf{Déduire le sexe du parent} à partir de ses gamètes observés : un seul type de gonosome ($X$ seulement) $\Rightarrow$ femme ; deux types ($X$ et $Y$) $\Rightarrow$ homme.
\item \textbf{Déterminer le sexe de l'enfant} : c'est \textbf{le spermatozoïde} qui décide du sexe ($X + X \Rightarrow$ fille ; $Y + X \Rightarrow$ garçon), avec une probabilité de $\dfrac{1}{2}$ pour chaque sexe à chaque fécondation (indépendance des événements).
\end{enumerate}
\end{methode}

\begin{exoresolu}[Gamètes et sexe d'un nouveau-né]
On observe au microscope les cellules reproductrices d'un individu A : elles contiennent chacune 23 chromosomes (noyau haploïde), et l'on distingue deux catégories de cellules en quantités égales, l'une possédant un grand chromosome sexuel (X), l'autre un très petit chromosome sexuel (Y).
\begin{enumerate}
\item Quel est le sexe de l'individu A ? Justifie.
\item L'individu A et son épouse attendent un enfant. À l'aide d'un échiquier de croisement, détermine la probabilité d'avoir une fille.
\end{enumerate}
\tcblower
\textbf{1. Sexe de l'individu A.}
Les cellules observées ont 23 chromosomes : ce sont des gamètes ($n = 23$). Il existe \textbf{deux types} de gamètes en proportions égales : les uns portent un grand gonosome ($X$), les autres un très petit ($Y$). Or seul un individu de formule $46, XY$ produit par méiose deux catégories de spermatozoïdes : 50 \% à $X$ et 50 \% à $Y$. L'individu A est donc un \textbf{homme}.
\textbf{2. Probabilité d'avoir une fille.}
L'épouse ($46, XX$) ne produit que des ovules à $X$. Échiquier de croisement (gamètes en lignes/colonnes) :
\begin{center}
\begin{tabular}{|c|c|c|}
\hline
\rowcolor{popblueL} \textbf{Gamètes} & \textbf{Ovule $X$} & \textbf{Ovule $X$} \\
\hline
\textbf{Spermatozoïde $X$} & $XX$ (fille) & $XX$ (fille) \\
\hline
\textbf{Spermatozoïde $Y$} & $XY$ (garçon) & $XY$ (garçon) \\
\hline
\end{tabular}
\end{center}
Sur 4 issues équiprobables, 2 donnent $XX$ (fille) et 2 donnent $XY$ (garçon). La probabilité d'avoir une fille est donc :
\[ P(\text{fille}) = \frac{2}{4} = \frac{1}{2} = 50\ \% \]
\textbf{Conclusion} : l'individu A est un homme car il produit deux types de spermatozoïdes ($X$ et $Y$) ; la probabilité d'avoir une fille est de $\dfrac{1}{2}$, car c'est le gonosome du spermatozoïde qui détermine le sexe de l'enfant.
\end{exoresolu}

\courssub{Hérédité liée au sexe (hétérochromosomique) : daltonisme et hémophilie}

\begin{definition}[Hérédité liée au chromosome X (hétérochromosomique)]
C'est la transmission d'un caractère dont le gène responsable est situé sur le \cle{chromosome X} (et non sur le Y, qui porte très peu de gènes, dont \emph{SRY}). Comme l'homme n'a qu'un seul X ($XY$), il exprime \textbf{tous} les allèles portés par son unique chromosome X (hémizygose). La femme ($XX$) a deux allèles : dominance/récessivité s'applique.
\end{definition}

\begin{propriete}[Règles de transmission X-liée récessive (Daltonisme, Hémophilie A et B)]
\begin{itemize}
\item \textbf{Allèles} : $X^N$ (normal, dominant) ; $X^d$ ou $X^h$ (muté, récessif).
\item \textbf{Génotypes / Phénotypes} :
  \begin{tabular}{lllc}
  Femelle : & $X^N X^N$ = Sain ; & $X^N X^d$ = \cle{Conductrice} (saine) ; & $X^d X^d$ = Malade (rare) \\
  Mâle : & $X^N Y$ = Sain ; & $X^d Y$ = \textbf{Malade} (un seul allèle suffit)
  \end{tabular}
\item \textbf{Règles clés} :
  \begin{enumerate}
  \item Un père malade transmet son $X^d$ à \textbf{toutes ses filles} (conductrices obligées) et son $Y$ à tous ses fils (sains).
  \item Une mère conductrice transmet $X^d$ à 50 \% de ses enfants (filles $\rightarrow$ conductrices ; fils $\rightarrow$ malades).
  \item Un fils malade a \textbf{toujours} reçu l'allèle muté de sa \textbf{mère}.
  \item Pas de transmission père $\rightarrow$ fils.
  \end{enumerate}
\end{itemize}
\end{propriete}

\begin{methode}[Interpréter un arbre généalogique : hérédité liée au sexe (X-récessif)]
\begin{enumerate}
\item \textbf{Lire l'arbre} : $\square$ = homme, $\bigcirc$ = femme ; symbole noir = atteint, blanc = sain, hachuré/point = conducteur connu.
\item \textbf{Repérer les indices X-récessif} :
  \begin{itemize}
  \item Majorité de malades de sexe \textbf{masculin}.
  \item Transmission \textbf{grand-père maternel $\rightarrow$ petit-fils} (via la fille conductrice).
  \item Pas de transmission père $\rightarrow$ fils.
  \item Femelles malades rares (nécessitent père malade + mère conductrice).
  \end{itemize}
\item \textbf{Poser les allèles} : $X^N$ (dominant), $X^a$ (récessif, $a$ = d pour daltonisme, h pour hémophilie).
\item \textbf{Déduire les génotypes} en remontant l'arbre (déduction logique).
\item \textbf{Construire l'échiquier} du couple concerné pour calculer les risques.
\end{enumerate}
\textbf{Phrase-clé Probatoire} : « Le père transmet son chromosome $X$ à toutes ses filles et son $Y$ à tous ses fils ; un garçon reçoit donc toujours son $X$ de sa mère. »
\end{methode}

\begin{exoresolu}[L'hémophilie dans une famille de Bafoussam]
Dans une famille, un couple sain a quatre enfants : deux filles saines, un garçon sain et un garçon hémophile. Le frère de la mère était également hémophile. On note $X^H$ l'allèle normal (dominant) et $X^h$ l'allèle responsable de l'hémophilie (récessif), porté par le chromosome X.
\begin{enumerate}
\item Le gène de l'hémophilie est-il dominant ou récessif ? Justifie.
\item Donne le génotype de la mère, du père et du garçon malade.
\item Quelle est la probabilité que le prochain enfant du couple soit un garçon hémophile ?
\end{enumerate}
\tcblower
\textbf{1. Dominance.} Les deux parents sont phénotypiquement sains et ont pourtant un enfant malade : l'allèle maladif était donc présent mais « caché » (non exprimé) chez un parent. L'allèle $X^h$ est \textbf{récessif} ; l'allèle normal $X^H$ est dominant. (Si $X^h$ était dominant, un parent sain ne pourrait pas le porter).
\textbf{2. Génotypes.}
\begin{itemize}
\item \textbf{Père sain} : il possède un seul X, qui est normal : $X^H Y$.
\item \textbf{Mère saine} : elle a un frère hémophile ($X^h Y$). Ce frère a reçu son $X^h$ de \textbf{leur mère} (grand-mère de l'enfant). La mère du couple a donc 1 chance sur 2 d'avoir reçu ce $X^h$. De plus, elle a un fils hémophile ($X^h Y$) : ce fils a obligatoirement reçu son $Y$ de son père et son $X^h$ de \textbf{sa mère}. La mère est donc certainement \textbf{conductrice} : $X^H X^h$.
\item \textbf{Garçon malade} : il a reçu $Y$ de son père et $X^h$ de sa mère : $X^h Y$.
\end{itemize}
\textbf{3. Probabilité d'un garçon hémophile.} Échiquier de croisement $X^H X^h \times X^H Y$ :
\begin{center}
\begin{tabular}{|c|c|c|}
\hline
\rowcolor{popblueL} & \textbf{Ovule $X^H$} & \textbf{Ovule $X^h$} \\
\hline
\textbf{Spermatozoïde $X^H$} & $X^H X^H$ (fille saine) & $X^H X^h$ (fille conductrice saine) \\
\hline
\textbf{Spermatozoïde $Y$} & $X^H Y$ (garçon sain) & $X^h Y$ (\textbf{garçon hémophile}) \\
\hline
\end{tabular}
\end{center}
Sur 4 issues équiprobables, une seule correspond à un garçon hémophile :
\[ P(\text{garçon hémophile}) = \frac{1}{4} = 25\ \% \]
(Calcul alternatif : $P(\text{garçon}) \times P(\text{transmission de } X^h \text{ par la mère}) = \dfrac{1}{2} \times \dfrac{1}{2} = \dfrac{1}{4}$.)
\textbf{Conclusion} : l'allèle hémophile est récessif et lié au chromosome X ; la mère est conductrice ($X^H X^h$), le père est $X^H Y$, le fils malade $X^h Y$ ; la probabilité d'un garçon hémophile à la prochaine naissance est de $\dfrac{1}{4}$.
\end{exoresolu}

\courssub{Hérédité autosomale : albinisme et drépanocytose}

\begin{definition}[Hérédité autosomale]
Le gène responsable est porté par un \cle{autosome} (paires 1 à 22). La transmission est \textbf{indépendante du sexe} : hommes et femmes sont atteints en proportions égales et transmettent le caractère de la même manière.
\end{definition}

\begin{propriete}[Spécificités de l'albinisme et de la drépanocytose (Autosomales Récessives)]
\begin{itemize}
\item \textbf{Albinisme} : gène \emph{TYR} (chromosome 11) ou \emph{OCA2} (chromosome 15). Allèles : $N$ (normal, dominant) / $a$ (albinos, récessif). $a/\!/a$ = albinos (absence de mélanine) ; $N/\!/a$ = porteur sain (phénotype normal).
\item \textbf{Drépanocytose} : gène \emph{HBB} (chromosome 11). Allèles : $A$ (HbA normale) / $S$ (HbS mutée). \textbf{Codominance au niveau moléculaire} (HbA et HbS synthétisées chez $A/\!/S$), \textbf{récessivité au niveau clinique}.
  \begin{itemize}
  \item $A/\!/A$ : Sain non porteur.
  \item $A/\!/S$ : \cle{Porteur sain} (trait drépanocytaire). Asymptomatique en conditions normales ; résistance au paludisme (avantage sélectif \cle{hétérozygote} en zone endémique comme le Cameroun).
  \item $S/\!/S$ : \cle{Drépanocytaire majeur} (anémie hémolytique chronique, crises vaso-occlusives, infections, séquestration splénique).
  \end{itemize}
\end{itemize}
\end{propriete}

\begin{methode}[Interpréter un arbre généalogique : hérédité autosomale récessive]
\begin{enumerate}
\item \textbf{Vérifier l'indépendance du sexe} : hommes et femmes atteints en proportions comparables $\Rightarrow$ autosomale.
\item \textbf{Déterminer la dominance} : deux parents sains ayant un enfant malade $\Rightarrow$ allèle maladif \textbf{récessif}.
\item \textbf{Noter les allèles} : Majuscule = dominant (normal), Minuscule = récessif (maladif). Ex: $N/a$ ou $A/S$.
\item \textbf{Attribuer les génotypes} :
  \begin{itemize}
  \item Malade = homozygote récessif ($a/\!/a$ ou $S/\!/S$).
  \item Parent sain d'un enfant malade = \textbf{obligatoirement hétérozygote} ($N/\!/a$ ou $A/\!/S$) = \cle{porteur sain} (ou \cle{conducteur}).
  \end{itemize}
\item \textbf{Construire l'échiquier de croisement} et calculer les proportions (génotypiques et phénotypiques).
\end{enumerate}
\end{methode}

\begin{exoresolu}[Albinisme : parents sains, enfant atteint à Bamenda]
Un couple à la peau normalement pigmentée, vivant à Bamenda, donne naissance à un enfant albinos (photophobie, peau blanche, cheveux blonds). On note $N$ l'allèle normal (pigmentation présente, dominant) et $a$ l'allèle responsable de l'albinisme (récessif).
\begin{enumerate}
\item L'allèle de l'albinisme est-il dominant ou récessif ? Justifie.
\item Donne les génotypes des deux parents et de l'enfant.
\item Ce couple peut-il avoir un autre enfant albinos ? Avec quelle probabilité ?
\end{enumerate}
\tcblower
\textbf{1. Dominance.} Deux parents au phénotype normal ont un enfant albinos : chacun porte donc l'allèle $a$ sans l'exprimer. L'allèle $a$ est \textbf{récessif} ; $N$ est dominant. Le gène est autosomal : l'albinisme touche indifféremment filles et garçons.
\textbf{2. Génotypes.}
\begin{itemize}
\item L'enfant albinos exprime le caractère récessif : il est forcément homozygote $a/\!/a$.
\item Chaque parent a fourni un allèle $a$ à l'enfant tout en étant sain : les deux parents sont \textbf{hétérozygotes porteurs sains} : $N/\!/a$.
\end{itemize}
\textbf{3. Probabilité d'un autre enfant albinos.} Échiquier de croisement $N/\!/a \times N/\!/a$ :
\begin{center}
\begin{tabular}{|c|c|c|}
\hline
\rowcolor{popblueL} & $N$ & $a$ \\
\hline
$N$ & $N/\!/N$ (sain non porteur) & $N/\!/a$ (porteur sain) \\
\hline
$a$ & $N/\!/a$ (porteur sain) & $a/\!/a$ (\textbf{albinos}) \\
\hline
\end{tabular}
\end{center}
Résultats : $\dfrac{1}{4}\ N/\!/N$ (sain non porteur), $\dfrac{2}{4}\ N/\!/a$ (porteurs sains), $\dfrac{1}{4}\ a/\!/a$ (albinos). La probabilité d'avoir un enfant albinos est donc :
\[ P(a/\!/a) = \frac{1}{4} = 25\ \% \]
et cela \textbf{à chaque grossesse}, quel que soit le sexe de l'enfant (gène autosomal) et indépendamment des naissances précédentes.
\textbf{Conclusion} : l'albinisme est une anomalie autosomale récessive ; les parents sont tous deux porteurs sains $N/\!/a$ ; chaque enfant à naître a 25 \% de risque d'être albinos, 50 \% d'être porteur sain et 25 \% d'être non porteur.
\end{exoresolu}

\begin{exoresolu}[Drépanocytose : le dépistage prénuptial à Yaoundé]
Avant leur mariage à Yaoundé, deux jeunes gens effectuent un test de dépistage (émigration de l'hémoglobine sur gel) : tous deux sont hétérozygotes $A/\!/S$ (porteurs sains du trait drépanocytaire). On note $A$ l'allèle codant l'hémoglobine normale (HbA) et $S$ l'allèle codant l'hémoglobine anormale (HbS).
\begin{enumerate}
\item Pourquoi dit-on que ces deux personnes sont « porteuses saines » ?
\item Détermine, à l'aide d'un échiquier de croisement, les proportions de génotypes attendus chez leurs enfants.
\item Quelle est la probabilité qu'un de leurs enfants soit atteint de drépanocytose majeure ?
\end{enumerate}
\tcblower
\textbf{1. Porteurs sains.} Chacun possède un allèle normal $A$ et un allèle muté $S$. L'allèle $A$ assure la production d'hémoglobine normale (HbA) en quantité suffisante pour assurer le transport de l'oxygène : les hétérozygotes $A/\!/S$ ne présentent pas les crises de la maladie (ils sont \textbf{sains}), mais ils peuvent transmettre l'allèle $S$ à leurs enfants (ils sont \textbf{porteurs}). C'est l'avantage sélectif de l'hétérozygote (résistance au paludisme) qui maintient la fréquence de l'allèle $S$ élevée au Cameroun.
\textbf{2. Échiquier de croisement} $A/\!/S \times A/\!/S$ :
\begin{center}
\begin{tabular}{|c|c|c|}
\hline
\rowcolor{popblueL} & $A$ & $S$ \\
\hline
$A$ & $A/\!/A$ (sain non porteur) & $A/\!/S$ (porteur sain) \\
\hline
$S$ & $A/\!/S$ (porteur sain) & $S/\!/S$ (\textbf{drépanocytaire majeur}) \\
\hline
\end{tabular}
\end{center}
Proportions attendues (loi de Mendel) : $\dfrac{1}{4}\ A/\!/A$ ; $\dfrac{2}{4} = \dfrac{1}{2}\ A/\!/S$ ; $\dfrac{1}{4}\ S/\!/S$.
\textbf{3. Probabilité de drépanocytose majeure.} Seuls les homozygotes $S/\!/S$ sont gravement malades (globules rouges falciformes, anémie chronique, crises douloureuses, risque infectieux majeur) :
\[ P(S/\!/S) = \frac{1}{4} = 25\ \% \]
\textbf{Conclusion} : deux porteurs sains $A/\!/S$ ont, à chaque naissance, 25 \% de risque d'avoir un enfant atteint de drépanocytose majeure ($S/\!/S$), 50 \% d'avoir un porteur sain et 25 \% d'avoir un enfant non porteur. C'est pourquoi le dépistage prénuptial (et prénatal) est fortement encouragé au Cameroun, où la prévalence du trait drépanocytaire ($A/\!/S$) atteint 20-30 \% de la population.
\end{exoresolu}

\courssub{Synthèse : démarche d'interprétation et lutte contre les préjugés}

\begin{methode}[Démarche universelle d'interprétation d'un arbre généalogique (Probatoire)]
\begin{enumerate}
\item \textbf{Observer} : Symboles (carré/cerclé), noircis/blancs, générations, consanguinité (double trait).
\item \textbf{Tester le sexe} : Atteints $\approx$ 50/50 H/F $\rightarrow$ \textbf{Autosomale} ; Atteints quasi-exclusivement H $\rightarrow$ \textbf{X-liée}.
\item \textbf{Tester la dominance} (si autosomale) : Parents sains + enfant malade $\rightarrow$ \textbf{Récessive} ; Malade à chaque génération $\rightarrow$ \textbf{Dominante}.
\item \textbf{Tester la transmission} (si X-liée) : Père malade $\rightarrow$ toutes filles atteintes/conductrices, 0 fils atteints $\rightarrow$ \textbf{X-dominante} (rare) ; Père malade $\rightarrow$ toutes filles conductrices, 0 fils atteints + Mère conductrice $\rightarrow$ 50\% fils malades $\rightarrow$ \textbf{X-récessive}.
\item \textbf{Noter les allèles} (convention : majuscule dominant, minuscule récessif ; $X^A$, $X^a$ pour X-liée).
\item \textbf{Déduire les génotypes} de tous les individus de l'arbre.
\item \textbf{Calculer les risques} pour une union future via l'échiquier de croisement.
\end{enumerate}
\end{methode}

\begin{aretenir}[Arguments scientifiques contre les préjugés sociaux (Éducation à la citoyenneté)]
\begin{itemize}
\item \textbf{« C'est la femme qui détermine le sexe de l'enfant »} $\rightarrow$ \textbf{FAUX}. Scientifiquement, la femme ne fournit que des ovules $X$. C'est le spermatozoïde ($X$ ou $Y$) qui détermine le sexe. Accuser la femme est une injustice sans fondement biologique.
\item \textbf{« L'albinisme / la drépanocytose sont des malédictions ou des sorts »} $\rightarrow$ \textbf{FAUX}. Ce sont des maladies \textbf{génétiques autosomales récessives}, dues à des mutations ponctuelles de gènes précis (\emph{TYR/OCA2}, \emph{HBB}). Elles suivent des lois mathématiques prévisibles (Mendel). Les parents porteurs sains ($N/\!/a$ ou $A/\!/S$) sont biologiquement normaux.
\item \textbf{« L'hémophilie / le daltonisme viennent de la mère »} $\rightarrow$ \textbf{PARTIELLEMENT VRAI mais trompeur}. L'allèle muté est sur le X. Un fils malade l'a reçu de sa mère (conductrice), mais cette mère l'a souvent reçu de \textbf{son père} (grand-père maternel). La responsabilité est partagée génétiquement ; aucune « faute » morale n'existe.
\item \textbf{« Un couple de porteurs sains ne doit pas avoir d'enfants »} $\rightarrow$ \textbf{CHOIX PERSONNEL ÉCLAIRÉ}. Le risque est de 25 \% par grossesse. Le conseil génétique (existant au Cameroun aux centres de dépistage de la drépanocytose) permet d'informer sur les options : diagnostic prénatal (DPN), fécondation in vitro avec DPI, adoption, ou acceptation du risque. La science éclaire le choix, ne l'impose pas.
\item \textbf{La trisomie 21 n'est pas une « faute » maternelle} $\rightarrow$ C'est un accident méiotique (non-disjonction) dont la fréquence augmente avec l'âge maternel, mais qui peut survenir à tout âge et aussi d'origine paternelle. Le dépistage prénatal (marqueurs sériques, ADN fœtal circulant, caryotype sur liquide amniotique) permet une prise en charge adaptée.
\end{itemize}
\end{aretenir}

\begin{savaistu}[Contexte camerounais : Drépanocytose et Santé Publique]
Le Cameroun est en zone de forte prévalence de l'hémoglobine S. Environ \textbf{25 à 30 \% de la population} est porteuse saine ($A/\!/S$). Chaque année, environ \textbf{6 000 à 8 000 enfants} naissent avec une drépanocytose majeure ($S/\!/S$). Le Ministère de la Santé Publique (MINSANTE) a mis en place un \textbf{Programme National de Lutte contre la Drépanocytose (PNLCD)} incluant : le dépistage néonatal systématique (test de Guthrie modifié ou électrophorèse Hb), la prise en charge gratuite des enfants de moins de 5 ans (pénicilline, acide folique, vaccins), et la promotion du \textbf{dépistage prénuptial} (carte de génotype Hb). Connaître son génotype ($AA$, $AS$, $SS$) avant le mariage est un acte de responsabilité citoyenne et de santé publique.
\end{savaistu}

\begin{attention}[Les 5 erreurs qui coûtent des points au Probatoire]
\begin{enumerate}
\item \textbf{Confondre anomalie chromosomique et anomalie génique.} La trisomie 21 se \emph{voit} sur le caryotype (47 chromosomes) ; l'albinisme, la drépanocytose, le daltonisme et l'hémophilie sont des mutations de \emph{gènes}, invisibles sur le caryotype. Ne jamais écrire « le caryotype montre le gène de l'hémophilie ».
\item \textbf{Mal rédiger la formule chromosomique.} Elle s'écrit : nombre total, virgule, gonosomes : $46, XX$ ; $46, XY$ ; $47, XY, +21$. Écrire « 46 chromosomes XY » sans virgule, ou oublier le nombre total, fait perdre des points.
\item \textbf{Dire que « la mère donne le sexe de l'enfant ».} C'est faux : la femme ne produit que des ovules à $X$ ; c'est le \textbf{spermatozoïde} ($X$ ou $Y$) qui détermine le sexe. Ce préjugé, source d'injustices dans les familles, doit être combattu par l'argument scientifique.
\item \textbf{Oublier que les gènes liés au sexe sont portés par X, jamais par Y.} Un homme ne peut pas transmettre le daltonisme ou l'hémophilie à ses fils (il leur donne son $Y$) ; il transmet son $X^d$ à \emph{toutes} ses filles, qui deviennent conductrices. Un garçon malade a toujours reçu l'allèle de sa \textbf{mère}.
\item \textbf{Se tromper dans les proportions de l'échiquier de croisement.} Pour deux hétérozygotes $N/\!/a \times N/\!/a$, les proportions sont $\dfrac{1}{4} - \dfrac{2}{4} - \dfrac{1}{4}$ et non « 1/3 de malades » ; et la probabilité « garçon hémophile » est $\dfrac{1}{2} \times \dfrac{1}{2} = \dfrac{1}{4}$, pas $\dfrac{1}{2}$ : il faut multiplier la probabilité du sexe par celle de la transmission de l'allèle. Toujours construire l'échiquier complet avant de conclure.
\end{enumerate}
\end{attention}

\begin{exercice}[Entraînement Probatoire : Arbre généalogique mixte]
On étudie une famille sur trois générations.
\begin{itemize}
\item G-I : Grand-père (I-1) daltonien, Grand-mère (I-2) saine.
\item G-II : Leurs 3 enfants : Fille (II-1) saine, Fils (II-2) daltonien, Fille (II-3) saine.
\item G-III : La fille II-1 épouse un homme sain. Ils ont 2 fils : III-1 sain, III-2 daltonien. La fille II-3 épouse un homme daltonien. Ils ont une fille III-3 daltonienne.
\end{itemize}
Le gène du daltonisme est récessif et porté par le chromosome X ($X^D$ normal, $X^d$ daltonien).
\begin{enumerate}
\item Représentez l'arbre généalogique complet avec les symboles standards.
\item Déterminez les génotypes de tous les individus (I-1 à III-3).
\item Calculez la probabilité qu'un futur enfant du couple II-3 $\times$ homme daltonien soit une fille daltonienne.
\end{enumerate}
\end{exercice}

\begin{corrige}[Exercice n°1]
\textbf{1. Arbre généalogique} : (À dessiner : I-1 $\square$ noir, I-2 $\bigcirc$ blanc ; II-1 $\bigcirc$ blanc, II-2 $\square$ noir, II-3 $\bigcirc$ blanc ; III-1 $\square$ blanc, III-2 $\square$ noir, III-3 $\bigcirc$ noir. Unions : II-1 x $\square$ blanc ; II-3 x $\square$ noir).
\textbf{2. Génotypes} :
\begin{itemize}
\item I-1 (GP daltonien) : $X^d Y$
\item I-2 (GM saine) : le fils II-2 est daltonien ($X^d Y$) : il a reçu $Y$ du père I-1 et $X^d$ de la mère I-2. La mère I-2, phénotypiquement saine mais porteuse de l'allèle $X^d$, est donc conductrice : $X^D X^d$.
\item II-1 (Fille saine) : A reçu $X^d$ du père I-1, $X^D$ de la mère I-2 $\rightarrow$ $X^D X^d$ (conductrice).
\item II-2 (Fils daltonien) : A reçu $Y$ du père, $X^d$ de la mère $\rightarrow$ $X^d Y$.
\item II-3 (Fille saine) : A reçu $X^d$ du père, $X^D$ de la mère $\rightarrow$ $X^D X^d$ (conductrice).
\item Mari de II-1 (Homme sain) : $X^D Y$.
\item III-1 (Fils sain de II-1) : $X^D Y$ (a reçu Y du père, $X^D$ de la mère).
\item III-2 (Fils daltonien de II-1) : $X^d Y$ (a reçu Y du père, $X^d$ de la mère II-1).
\item Mari de II-3 (Homme daltonien) : $X^d Y$.
\item III-3 (Fille daltonienne de II-3) : A reçu $X^d$ du père, $X^d$ de la mère II-3 $\rightarrow$ $X^d X^d$.
\end{itemize}
\textbf{3. Probabilité fille daltonienne (Couple II-3 $X^D X^d$ x $X^d Y$)} :
Échiquier :
\begin{center}
\begin{tabular}{|c|c|c|}
\hline
\rowcolor{popblueL} & Ovule $X^D$ & Ovule $X^d$ \\
\hline
Spermatozoïde $X^d$ & $X^D X^d$ (fille conductrice) & $X^d X^d$ (\textbf{fille daltonienne}) \\
\hline
Spermatozoïde $Y$ & $X^D Y$ (garçon sain) & $X^d Y$ (garçon daltonien) \\
\hline
\end{tabular}
\end{center}
$P(\text{fille daltonienne}) = \dfrac{1}{4} = 25\ \%$.
\end{corrige}

\newpage

\coursec{Exercices}

\courssub{Je vérifie mes connaissances}

\begin{exercice}[QCM : le caryotype humain]
Pour chaque question, choisis la (ou les) bonne(s) réponse(s) en justifiant brièvement.

\textbf{1.} Le caryotype normal d'un homme s'écrit :
\begin{itemize}
\item[a)] 44 autosomes + XX
\item[b)] 46, XY
\item[c)] 46, XX
\item[d)] 22 paires d'autosomes + XY
\end{itemize}

\textbf{2.} La trisomie 21 est caractérisée par :
\begin{itemize}
\item[a)] l'absence d'un chromosome 21
\item[b)] la présence de trois chromosomes 21
\item[c)] une formule chromosomique 47, XY, +21 chez un garçon
\item[d)] une anomalie d'un gène du chromosome 21
\end{itemize}

\textbf{3.} Le daltonisme est une maladie :
\begin{itemize}
\item[a)] autosomale dominante
\item[b)] autosomale récessive
\item[c)] liée au chromosome X, récessive
\item[d)] liée au chromosome Y
\end{itemize}

\textbf{4.} Un individu de génotype AS pour le gène de l'hémoglobine est :
\begin{itemize}
\item[a)] atteint de drépanocytose grave
\item[b)] sain, porteur du caractère drépanocytaire
\item[c)] dépourvu d'hémoglobine anormale
\item[d)] incapable de transmettre l'allèle S
\end{itemize}
\end{exercice}

\begin{exercice}[Vrai ou faux ? Justifie]
Indique si chaque affirmation est vraie ou fausse, puis justifie ta réponse en une ou deux phrases.
\begin{enumerate}
\item Le caryotype humain normal comporte 23 paires de chromosomes, dont 22 paires d'autosomes.
\item Une femme peut transmettre l'allèle de l'hémophilie à son fils alors qu'elle n'est pas malade.
\item Un père daltonien transmet toujours son allèle muté à ses fils.
\item Deux parents à la peau normalement pigmentée peuvent avoir un enfant albinos.
\item La trisomie 21 résulte d'une mutation d'un gène du chromosome 21.
\end{enumerate}
\end{exercice}

\begin{exercice}[Définitions et vocabulaire]
Définis précisément les termes suivants en une ou deux phrases chacun :
\begin{enumerate}
\item Caryotype.
\item Formule chromosomique.
\item Femme conductrice (dans le cas de l'hémophilie).
\item Hérédité autosomale récessive.
\item Non-disjonction chromosomique.
\end{enumerate}
\end{exercice}

\begin{exercice}[Formules chromosomiques]
Écris la formule chromosomique complète (nombre total de chromosomes, chromosomes sexuels, anomalie éventuelle) des individus suivants :
\begin{enumerate}
\item Une fille saine.
\item Un garçon atteint de trisomie 21.
\item Une fille atteinte de trisomie 21.
\item Un homme dont les cellules ne possèdent qu'un seul chromosome sexuel X (syndrome de Turner, 45 chromosomes).
\item Un garçon dont les cellules possèdent 47 chromosomes avec deux chromosomes X et un chromosome Y (syndrome de Klinefelter).
\end{enumerate}
\end{exercice}

\courssub{Je m'entraîne}

\begin{exercice}[Réalisation d'un caryotype]
Au laboratoire de cytogénétique de l'hôpital, on photographie les chromosomes d'une cellule en métaphase. On observe : 22 paires de chromosomes deux à deux semblables, plus une paire de chromosomes sexuels de tailles inégales (un grand chromosome et un petit).
\begin{enumerate}
\item Combien de chromosomes cette cellule contient-elle au total ?
\item Quel est le sexe de l'individu ? Justifie.
\item Écris la formule chromosomique de cet individu.
\item Explique en deux phrases comment on réalise pratiquement un caryotype à partir de la photographie des chromosomes.
\end{enumerate}
\end{exercice}

\begin{exercice}[Un caryotype anormal]
L'analyse des cellules d'un nouveau-né montre 47 chromosomes, dont trois chromosomes 21 ; la paire sexuelle est XX.
\begin{enumerate}
\item Identifie l'anomalie chromosomique et nomme la maladie correspondante.
\item Écris la formule chromosomique de cet enfant.
\item Cette anomalie résulte d'une non-disjonction des chromosomes 21 au cours de la méiose chez l'un des parents. Schématise simplement (deux gamètes anormaux possibles) comment un gamète peut porter deux chromosomes 21 au lieu d'un.
\item Cite deux caractères modifiés observables chez un enfant atteint de cette anomalie.
\end{enumerate}
\end{exercice}

\begin{exercice}[Daltonisme : échiquier de croisement]
On note $D$ l'allèle normal (vision des couleurs normale) et $d$ l'allèle responsable du daltonisme, portés par le chromosome X. Une femme conductrice épouse un homme à la vision normale.
\begin{enumerate}
\item Écris les génotypes des deux parents.
\item Construis l'échiquier de croisement et détermine les génotypes possibles des enfants.
\item Donne les phénotypes correspondants (fille saine, fille conductrice, garçon sain, garçon daltonien) et leurs proportions.
\item Quel est le risque, pour cette famille, d'avoir un garçon daltonien ? Une fille daltonienne ? Justifie.
\end{enumerate}
\end{exercice}

\begin{exercice}[Drépanocytose et dépistage prénuptial]
Un homme de génotype AS épouse une femme de génotype AS (allèle A : hémoglobine normale ; allèle S : hémoglobine drépanocytaire).
\begin{enumerate}
\item Construis l'échiquier de croisement.
\item Calcule les proportions attendues de descendants AA, AS et SS.
\item Précise le phénotype de chaque génotype.
\item Explique en quelques lignes pourquoi le dépistage prénuptial de la drépanocytose est recommandé avant le mariage.
\end{enumerate}
\end{exercice}

\courssub{Je me prépare au Probatoire}

\begin{exercice}[L'hémophilie dans une famille]
Dans une famille, on étudie la transmission de l'hémophilie (allèle $h$ récessif porté par le chromosome X ; allèle normal $H$). L'arbre généalogique est le suivant : à la génération I, un homme hémophile (I-1) épouse une femme saine non conductrice (I-2). Ils ont deux enfants (génération II) : une fille (II-1) et un garçon (II-2), tous deux phénotypiquement sains. La fille II-1 épouse un homme sain (II-3) ; ils ont trois enfants (génération III) : une fille saine (III-1), un garçon sain (III-2) et un garçon hémophile (III-3).
\begin{enumerate}
\item Représente cet arbre généalogique en utilisant les conventions : carrés pour les garçons, cercles pour les filles, symboles noircis pour les individus malades.
\item Écris les génotypes de I-1, I-2 et justifie que la fille II-1 est obligatoirement conductrice.
\item Détermine le génotype de II-2 et explique pourquoi il ne peut pas être conducteur.
\item Par un échiquier de croisement du couple II-1 $\times$ II-3, montre que la naissance du garçon III-3 était prévisible et calcule le risque qu'avait ce couple d'avoir un garçon hémophile.
\item Justifie l'affirmation suivante : « un père hémophile ne transmet jamais l'allèle de la maladie à ses fils, mais le transmet à toutes ses filles ».
\end{enumerate}
\end{exercice}

\begin{exercice}[Albinisme : interprétation d'une observation familiale]
Dans un village, un couple à la peau normalement pigmentée donne naissance à un enfant albinos. Certains voisins accusent la mère d'infidélité ou parlent de « mauvais sort ». Tu es sollicité(e) pour expliquer scientifiquement la situation. On note $N$ l'allèle dominant commandant une pigmentation normale et $a$ l'allèle récessif responsable de l'albinisme.
\begin{enumerate}
\item Rappelle ce qu'est l'albinisme et précise s'il s'agit d'une anomalie chromosomique ou génique.
\item Explique, à l'aide des notions d'allèle dominant et récessif, comment deux parents à la peau normale peuvent avoir un enfant albinos. Donne les génotypes des parents et de l'enfant.
\item Construis l'échiquier de croisement correspondant et calcule le risque que le prochain enfant de ce couple soit albinos.
\item Ce risque dépend-il du sexe de l'enfant ? Justifie.
\item Rédige un court texte (5 à 8 lignes) destiné aux villageois pour combattre les préjugés et montrer que l'albinisme est une maladie héréditaire comme les autres, qui ne justifie ni accusation ni exclusion.
\end{enumerate}
\end{exercice}

\begin{exercice}[Arbre généalogique muet : démarche de synthèse]
On te présente un arbre généalogique sans indication sur la maladie transmise. Conventions : carrés = garçons, cercles = filles, symboles noircis = individus atteints. Voici les données : à la génération I, un couple sain (I-1 homme, I-2 femme) a quatre enfants (génération II) : deux filles saines (II-1, II-2), un garçon sain (II-3) et une fille atteinte (II-4). La fille atteinte II-4 épouse un homme sain (II-5) ; leurs trois enfants (génération III) sont tous sains : un garçon (III-1) et deux filles (III-2, III-3).
\begin{enumerate}
\item Dessine l'arbre généalogique correspondant à ces données.
\item La transmission est-elle dominante ou récessive ? Justifie en t'appuyant sur le couple I-1 $\times$ I-2 et l'enfant II-4.
\item La transmission est-elle autosomale ou liée au chromosome X ? Raisonne par l'absurde : montre que si l'allèle était récessif lié à X, la fille II-4 ne pourrait pas être atteinte avec un père sain. Conclus.
\item Écris les génotypes de I-1, I-2, II-4 (en notant $A$ l'allèle dominant et $a$ l'allèle récessif).
\item Détermine le ou les génotypes possibles de l'époux II-5 sachant que tous ses enfants sont sains, et calcule, pour chaque hypothèse, le risque d'avoir un enfant atteint.
\item Propose deux exemples de maladies humaines correspondant au mode de transmission ainsi déterminé.
\end{enumerate}
\end{exercice}

\newpage

\coursec{Corrigés des exercices}

\begin{corrige}[Exercice 1 — QCM : le caryotype humain]
\textbf{1. Réponses b) et d).} Le caryotype normal d'un homme comporte 46 chromosomes : 22 paires d'autosomes (soit 44 autosomes) et une paire de chromosomes sexuels XY. Les notations « 46, XY » et « 22 paires d'autosomes + XY » sont donc équivalentes. La réponse a) est fausse (XX est la paire sexuelle féminine) ; c) correspond à une femme.

\textbf{2. Réponses b) et c).} La trisomie 21 est la présence de \cle{trois chromosomes 21} au lieu de deux ; chez un garçon, la formule chromosomique s'écrit 47, XY, +21. La réponse a) décrirait une monosomie (non viable pour le chromosome 21) ; d) est fausse : la trisomie 21 est une anomalie du \emph{nombre} de chromosomes, pas d'un gène.

\textbf{3. Réponse c).} Le daltonisme est dû à un allèle récessif porté par le chromosome X (hérédité hétérochromosomique récessive). Le chromosome Y ne porte pas ce gène, et le gène n'est pas sur un autosome : a), b) et d) sont fausses.

\textbf{4. Réponse b).} L'allèle S est récessif : un individu AS possède une hémoglobine normale (A) et une hémoglobine anormale (S), il est sain mais \cle{porteur du caractère drépanocytaire} (on dit « drépanocytaire sain »). Il peut transmettre l'allèle S à ses enfants (d est donc faux) ; seul le génotype SS provoque la drépanocytose grave (a faux) ; c est faux car il produit de l'hémoglobine S.
\end{corrige}

\begin{corrige}[Exercice 2 — Vrai ou faux ? Justifie]
\begin{enumerate}
\item \textbf{Vrai.} Le caryotype humain normal comporte 23 paires, soit 46 chromosomes : 22 paires d'autosomes et 1 paire de chromosomes sexuels (XX ou XY).
\item \textbf{Vrai.} Une femme conductrice $X^{H}X^{h}$ possède l'allèle $h$ sur l'un de ses chromosomes X sans être malade (l'allèle $h$ est récessif) ; elle peut transmettre ce X porteur de $h$ à son fils, qui sera alors hémophile.
\item \textbf{Faux.} Un père transmet son chromosome Y à ses fils, jamais son X. L'allèle du daltonisme étant porté par X, un père daltonien transmet son allèle muté à toutes ses \emph{filles}, et à \emph{aucun} de ses fils.
\item \textbf{Vrai.} Si les deux parents sont hétérozygotes $Na$ (phénotype normal car $N$ domine $a$), chacun peut transmettre l'allèle $a$ ; l'enfant $aa$ sera albinos (probabilité 1/4 à chaque naissance).
\item \textbf{Faux.} La trisomie 21 est une anomalie du \emph{nombre} de chromosomes (trois chromosomes 21), due à une non-disjonction au cours de la méiose, et non à la mutation d'un gène.
\end{enumerate}
\end{corrige}

\begin{corrige}[Exercice 3 — Définitions et vocabulaire]
\begin{enumerate}
\item \textbf{Caryotype} : ensemble des chromosomes d'une cellule, photographiés en métaphase, puis classés par paires selon leur taille et la position de leur centromère. Il permet d'identifier le sexe et les anomalies chromosomiques.
\item \textbf{Formule chromosomique} : notation conventionnelle qui indique le nombre total de chromosomes, la paire sexuelle, puis l'anomalie éventuelle. Exemple : 47, XX, +21 (fille trisomique 21).
\item \textbf{Femme conductrice} : femme hétérozygote $X^{H}X^{h}$, phénotypiquement saine car l'allèle $h$ de l'hémophilie est récessif, mais capable de transmettre l'allèle malade à ses descendants.
\item \textbf{Hérédité autosomale récessive} : mode de transmission d'un caractère dont le gène est porté par un autosome (ni X ni Y) et dont l'allèle responsable ne s'exprime qu'à l'état homozygote ($aa$). Les hétérozygotes $Aa$ sont sains mais porteurs.
\item \textbf{Non-disjonction chromosomique} : défaut de séparation de deux chromosomes homologues (ou de deux chromatides) au cours de la méiose, produisant des gamètes à deux exemplaires d'un chromosome et des gamètes dépourvus de ce chromosome.
\end{enumerate}
\end{corrige}

\begin{corrige}[Exercice 4 — Formules chromosomiques]
\begin{enumerate}
\item Fille saine : \textbf{46, XX}.
\item Garçon trisomique 21 : \textbf{47, XY, +21}.
\item Fille trisomique 21 : \textbf{47, XX, +21}.
\item Syndrome de Turner : \textbf{45, X} (un seul chromosome sexuel X ; on note parfois 45, X0).
\item Syndrome de Klinefelter : \textbf{47, XXY}.
\end{enumerate}
\textbf{Vérification} : la formule chromosomique indique toujours, dans l'ordre, le nombre total de chromosomes, la paire sexuelle, puis l'anomalie numérique éventuelle précédée du signe + (chromosome surnuméraire) ou $-$ (chromosome manquant).
\end{corrige}

\begin{corrige}[Exercice 5 — Réalisation d'un caryotype]
\begin{enumerate}
\item La cellule contient 22 paires d'autosomes + 1 paire sexuelle = 23 paires, soit \textbf{46 chromosomes}.
\item La paire sexuelle est formée de deux chromosomes de tailles inégales : un grand chromosome X et un petit chromosome Y. L'individu est donc de \textbf{sexe masculin} (XY).
\item Formule chromosomique : \textbf{46, XY}.
\item Pratiquement : on cultive des cellules (lymphocytes sanguins) et on bloque la division en métaphase, stade où les chromosomes sont les plus condensés. On les colore, on les photographie au microscope, puis on découpe et on classe les chromosomes par paires d'homologues, par ordre décroissant de taille (paires 1 à 22), les chromosomes sexuels étant placés en dernier.
\end{enumerate}
\end{corrige}

\begin{corrige}[Exercice 6 — Un caryotype anormal]
\begin{enumerate}
\item L'enfant possède trois chromosomes 21 : il s'agit d'une \textbf{trisomie 21} (anomalie du nombre de chromosomes).
\item Formule chromosomique : \textbf{47, XX, +21} (47 chromosomes, paire sexuelle XX : c'est une fille).
\item Schéma de la non-disjonction : au lieu que les deux chromosomes 21 homologues se séparent en méiose I (ou les deux chromatides en méiose II), ils migrent ensemble dans le même gamète.
\begin{popfigure}
\begin{center}
\begin{tikzpicture}[scale=0.9]
\node at (0,2.6) {\textbf{Méiose normale}};
\draw (0,1.8) ellipse (0.9 and 0.5);
\node at (0,1.8) {\small 21\;21};
\draw[->, thick] (-0.4,1.2) -- (-1.2,0.4);
\draw[->, thick] (0.4,1.2) -- (1.2,0.4);
\draw (-1.6,0) ellipse (0.55 and 0.35); \node at (-1.6,0) {\small 21};
\draw (1.6,0) ellipse (0.55 and 0.35); \node at (1.6,0) {\small 21};
\node at (0,-0.8) {\small 2 gamètes à 1 chromosome 21};
\node at (7,2.6) {\textbf{Non-disjonction}};
\draw (7,1.8) ellipse (0.9 and 0.5);
\node at (7,1.8) {\small 21\;21};
\draw[->, thick] (6.6,1.2) -- (5.8,0.4);
\draw[->, thick] (7.4,1.2) -- (8.2,0.4);
\draw (5.4,0) ellipse (0.55 and 0.35); \node at (5.4,0) {\small 21\;21};
\draw (8.6,0) ellipse (0.55 and 0.35); \node at (8.6,0) {\small 0};
\node at (7,-0.8) {\small 1 gamète à 2 chr. 21 + 1 gamète sans chr. 21};
\end{tikzpicture}
\end{center}
\legende{Schéma simplifié de la non-disjonction du chromosome 21 (méiose I ou II) conduisant à un gamète disomique et un gamète nullisomique.}
\end{popfigure}
La fécondation d'un gamète à deux chromosomes 21 par un gamète normal donne un zygote à trois chromosomes 21.
\item Caractères modifiés observables (exemples) : un léger handicap mental, un visage caractéristique (yeux bridés, face aplatie), une petite taille, une faiblesse musculaire (hypotonie), parfois des malformations cardiaques.
\end{enumerate}
\end{corrige}

\begin{corrige}[Exercice 7 — Daltonisme : échiquier de croisement]
\begin{enumerate}
\item Génotypes des parents : femme conductrice $X^{D}X^{d}$ ; homme à vision normale $X^{D}Y$.
\item Gamètes : la mère produit $X^{D}$ et $X^{d}$ ; le père produit $X^{D}$ et $Y$.
\begin{center}
\begin{tabular}{|c|c|c|}
\hline
\rowcolor{popblueL} Mère $\backslash$ Père & $X^{D}$ & $Y$ \\
\hline
$X^{D}$ & $X^{D}X^{D}$ & $X^{D}Y$ \\
\hline
$X^{d}$ & $X^{D}X^{d}$ & $X^{d}Y$ \\
\hline
\end{tabular}
\end{center}
\item Phénotypes et proportions :
\begin{itemize}
\item $X^{D}X^{D}$ : fille saine non conductrice (1/4) ;
\item $X^{D}X^{d}$ : fille saine conductrice (1/4) ;
\item $X^{D}Y$ : garçon sain (1/4) ;
\item $X^{d}Y$ : garçon daltonien (1/4).
\end{itemize}
\item Risque d'avoir un garçon daltonien : $1/4$ de tous les enfants, soit \textbf{1 garçon sur 2} parmi les garçons. Risque d'avoir une fille daltonienne : \textbf{nul}, car une fille reçoit toujours le $X^{D}$ normal de son père, qui masque l'allèle $d$ éventuellement reçu de la mère.
\end{enumerate}
\end{corrige}

\begin{corrige}[Exercice 8 — Drépanocytose et dépistage prénuptial]
\begin{enumerate}
\item Chaque parent AS produit deux types de gamètes : A et S.
\begin{center}
\begin{tabular}{|c|c|c|}
\hline
\rowcolor{popblueL} Père $\backslash$ Mère & A & S \\
\hline
A & AA & AS \\
\hline
S & AS & SS \\
\hline
\end{tabular}
\end{center}
\item Proportions attendues : \textbf{1/4 AA ; 2/4 (soit 1/2) AS ; 1/4 SS}.
\item Phénotypes : AA = individu sain, hémoglobine normale ; AS = individu sain mais porteur du caractère drépanocytaire (drépanocytaire sain) ; SS = individu atteint de \textbf{drépanocytose grave} (anémie sévère, crises douloureuses).
\item Lorsque deux porteurs AS se marient, ils ont à chaque naissance un risque de 1/4 d'avoir un enfant SS gravement malade. Le dépistage prénuptial (électrophorèse de l'hémoglobine, pratiquée dans de nombreux centres de santé au Cameroun) informe les futurs époux de leur génotype : il ne s'agit pas d'interdire un mariage, mais de permettre au couple de décider en connaissance de cause et de bénéficier d'un suivi médical adapté.
\end{enumerate}
\end{corrige}

\begin{corrige}[Exercice 9 — L'hémophilie dans une famille (type Probatoire)]
\begin{enumerate}
\item Arbre généalogique :
\begin{popfigure}
\begin{center}
\begin{tikzpicture}[
male/.style={draw, rectangle, minimum size=0.6cm, inner sep=1pt},
female/.style={draw, circle, minimum size=0.6cm, inner sep=1pt},
sick/.style={fill=popdark},
scale=1]
% Génération I
\node[male, sick, label=above:I-1] (i1) at (0,0) {};
\node[female, label=above:I-2] (i2) at (2,0) {};
\draw (i1) -- (i2); % mariage
\draw (1,0) -- (1,-0.5); % descente
\draw (0,-0.5) -- (2,-0.5); % ligne de fratrie
% Génération II
\node[female, label=above:II-1] (ii1) at (0,-1.5) {};
\node[male, label=above:II-2] (ii2) at (1.5,-1.5) {};
\node[male, label=above:II-3] (ii3) at (3,-1.5) {};
\draw (ii1) -- (ii3); % mariage II-1 x II-3
\draw (1.5,-1.5) -- (1.5,-2.5); % descente couple II-1 x II-3
\draw (0.5,-2.5) -- (2.5,-2.5); % ligne fratrie III
% Génération III
\node[male, label=below:III-1] (iii1) at (0.5,-3.5) {};
\node[female, label=below:III-2] (iii2) at (1.5,-3.5) {};
\node[male, sick, label=below:III-3] (iii3) at (2.5,-3.5) {};
% Légende
\node[male, sick] at (4.2,-0.2) {}; \node[right] at (4.55,-0.2) {\small homme malade};
\node[male] at (4.2,-0.8) {}; \node[right] at (4.55,-0.8) {\small homme sain};
\node[female] at (4.2,-1.4) {}; \node[right] at (4.55,-1.4) {\small femme saine};
\end{tikzpicture}
\end{center}
\legende{Arbre généalogique de la transmission de l'hémophilie. Carrés : garçons ; cercles : filles ; symboles noircis : individus hémophiles.}
\end{popfigure}
\item Génotypes : I-1 est un homme hémophile : $X^{h}Y$ ; I-2 est une femme saine non conductrice : $X^{H}X^{H}$. La fille II-1 reçoit obligatoirement le chromosome $X^{h}$ de son père (un père transmet son unique X à toutes ses filles) et un $X^{H}$ de sa mère : elle est donc $X^{H}X^{h}$, c'est-à-dire \textbf{obligatoirement conductrice}, tout en étant saine (l'allèle $h$ est récessif).
\item II-2 est un garçon sain : il a reçu le Y de son père et un $X^{H}$ de sa mère ; son génotype est $X^{H}Y$. Un homme ne possède qu'un seul X : il est soit $X^{H}Y$ (sain), soit $X^{h}Y$ (malade). La notion de « conducteur » n'existe donc pas chez l'homme pour un caractère lié à X.
\item Échiquier du couple II-1 ($X^{H}X^{h}$) $\times$ II-3 ($X^{H}Y$) :
\begin{center}
\begin{tabular}{|c|c|c|}
\hline
\rowcolor{popblueL} Mère $\backslash$ Père & $X^{H}$ & $Y$ \\
\hline
$X^{H}$ & $X^{H}X^{H}$ & $X^{H}Y$ \\
\hline
$X^{h}$ & $X^{H}X^{h}$ & $X^{h}Y$ \\
\hline
\end{tabular}
\end{center}
Le génotype $X^{h}Y$ (garçon hémophile) apparaît avec une probabilité de 1/4 : la naissance de III-3 était donc prévisible. Risque d'avoir un garçon hémophile : $1/4$ des enfants, soit \textbf{1 garçon sur 2}.
\item Un père transmet son chromosome \textbf{Y} à ses fils : il ne peut donc jamais leur transmettre l'allèle $h$ porté par son X. En revanche, il transmet son unique chromosome X, porteur de $h$, à \textbf{toutes} ses filles, qui seront au minimum conductrices.
\end{enumerate}
\textbf{Barème indicatif (sur 10)} : arbre correct et légendé (2 pts) ; génotypes de I-1, I-2 et justification de II-1 conductrice (2,5 pts) ; génotype de II-2 et explication (1,5 pt) ; échiquier et calcul du risque (2,5 pts) ; justification finale X/Y (1,5 pt).

\textbf{Points de vigilance} : toujours écrire les allèles en exposant du X ($X^{H}$, $X^{h}$) et jamais sur le Y ; justifier chaque génotype par l'origine parentale des chromosomes sexuels ; distinguer clairement « risque parmi tous les enfants » (1/4) et « risque parmi les garçons » (1/2).
\end{corrige}

\begin{corrige}[Exercice 10 — Albinisme : interprétation d'une observation familiale (type Probatoire)]
\begin{enumerate}
\item L'albinisme est une maladie héréditaire caractérisée par l'absence de mélanine (pigment) dans la peau, les poils, les cheveux et les yeux. C'est une anomalie \textbf{génique} (mutation d'un gène de la mélanogenèse), et non chromosomique : le caryotype d'une personne albinos est normal (46, XX ou 46, XY).
\item L'allèle $a$ est récessif : il ne s'exprime qu'à l'état homozygote $aa$. Deux parents à la peau normale mais hétérozygotes $Na$ (porteurs sains) peuvent chacun transmettre l'allèle $a$. Génotypes : parents $Na \times Na$ ; enfant albinos $aa$.
\item Échiquier de croisement :
\begin{center}
\begin{tabular}{|c|c|c|}
\hline
\rowcolor{popblueL} Père $\backslash$ Mère & $N$ & $a$ \\
\hline
$N$ & $NN$ & $Na$ \\
\hline
$a$ & $Na$ & $aa$ \\
\hline
\end{tabular}
\end{center}
Descendance : 1/4 $NN$, 2/4 $Na$ (peau normale), 1/4 $aa$ (albinos). Le risque que le prochain enfant soit albinos est de \textbf{1/4 (25 \%)} à chaque naissance, indépendamment des naissances précédentes.
\item \textbf{Non}, ce risque ne dépend pas du sexe : le gène est porté par un \textbf{autosome}, présent en deux exemplaires chez les deux sexes. Garçons et filles ont le même risque (1/4).
\item Texte destiné aux villageois : « L'enfant albinos de ce couple n'est ni un prodige ni une punition : c'est un enfant comme les autres, porteur d'une maladie héréditaire bien connue des médecins. Ses parents, tous deux à la peau normale, portaient chacun, sans le savoir, un allèle caché de l'albinisme ; en le transmettant ensemble, ils ont donné naissance à un enfant albinos, ce qui arrive une fois sur quatre dans cette situation. Nulle infidélité, nul mauvais sort : la science l'explique et de nombreuses familles dans le monde connaissent la même chose. Accueillons cet enfant avec affection, protégeons sa peau et ses yeux du soleil, et refusons toute exclusion : les personnes albinos ont les mêmes droits et les mêmes talents que chacun. »
\end{enumerate}
\textbf{Barème indicatif (sur 10)} : définition et nature génique (2 pts) ; explication avec génotypes (2 pts) ; échiquier et calcul du risque (2,5 pts) ; indépendance vis-à-vis du sexe justifiée (1,5 pt) ; texte de lutte contre les préjugés (2 pts).

\textbf{Points de vigilance} : préciser que les parents sont \emph{obligatoirement} hétérozygotes puisqu'ils ont un enfant $aa$ ; ne pas confondre « porteur sain » (hétérozygote) et « malade » ; citer le caractère autosomique pour justifier l'égalité des risques entre garçons et filles.
\end{corrige}

\begin{corrige}[Exercice 11 — Arbre généalogique muet : démarche de synthèse (type Probatoire)]
\begin{enumerate}
\item Arbre généalogique :
\begin{popfigure}
\begin{center}
\begin{tikzpicture}[
male/.style={draw, rectangle, minimum size=0.6cm, inner sep=1pt},
female/.style={draw, circle, minimum size=0.6cm, inner sep=1pt},
sick/.style={fill=popdark},
scale=1]
% Génération I
\node[male, label=above:I-1] (i1) at (0,0) {};
\node[female, label=above:I-2] (i2) at (2,0) {};
\draw (i1) -- (i2);
\draw (1,0) -- (1,-0.5);
\draw (-2,-0.5) -- (4,-0.5); % ligne fratrie
% Génération II
\node[female, label=above:II-1] (ii1) at (-2,-1.5) {};
\node[female, label=above:II-2] (ii2) at (-0.5,-1.5) {};
\node[male, label=above:II-3] (ii3) at (1,-1.5) {};
\node[female, sick, label=above:II-4] (ii4) at (2.5,-1.5) {};
\node[male, label=above:II-5] (ii5) at (4.5,-1.5) {};
\draw (ii4) -- (ii5); % mariage II-4 x II-5
\draw (3.5,-1.5) -- (3.5,-2.5); % descente
\draw (2,-2.5) -- (5,-2.5); % ligne fratrie III
% Génération III
\node[male, label=below:III-1] (iii1) at (2,-3.5) {};
\node[female, label=below:III-2] (iii2) at (3.5,-3.5) {};
\node[female, label=below:III-3] (iii3) at (5,-3.5) {};
\end{tikzpicture}
\end{center}
\legende{Arbre généalogique muet. Carrés : garçons ; cercles : filles ; symbole noirci : individu atteint (II-4).}
\end{popfigure}
\item Le couple I-1 $\times$ I-2 est sain et donne naissance à une enfant atteinte (II-4). Or deux parents sains ne peuvent avoir un enfant malade que s'ils portent chacun l'allèle de la maladie à l'état caché : la transmission est donc \textbf{récessive}.
\item Raisonnement par l'absurde : supposons l'allèle récessif lié au chromosome X. La fille II-4, atteinte, serait alors $X^{a}X^{a}$ : elle aurait dû recevoir un $X^{a}$ de son père I-1, qui serait donc $X^{a}Y$, c'est-à-dire \emph{malade}. Or I-1 est sain : l'hypothèse est impossible. La transmission est donc \textbf{autosomale récessive}.
\item Génotypes : I-1 : $Aa$ ; I-2 : $Aa$ (obligatoirement hétérozygotes, car sains avec un enfant $aa$) ; II-4 : $aa$.
\item II-4 ($aa$) transmet un allèle $a$ à tous ses enfants. Ses enfants étant tous sains, ils ont reçu un allèle $A$ de leur père II-5 ; ils sont tous $Aa$. Deux hypothèses pour II-5 :
\begin{itemize}
\item II-5 est $AA$ : tous les enfants sont $Aa$ (sains) ; risque d'enfant atteint : \textbf{0}.
\item II-5 est $Aa$ : croisement $aa \times Aa$ donne 1/2 $Aa$ (sains) et 1/2 $aa$ (atteints) ; risque d'enfant atteint : \textbf{1/2} à chaque naissance (le fait que les trois premiers enfants soient sains n'exclut pas cette hypothèse : probabilité $(1/2)^3 = 1/8$, tout à fait possible).
\end{itemize}
\item Exemples de maladies à transmission autosomale récessive : l'\textbf{albinisme} et la \textbf{drépanocytose}.
\end{enumerate}
\textbf{Barème indicatif (sur 10)} : arbre correct (2 pts) ; caractère récessif justifié (1,5 pt) ; raisonnement par l'absurde pour l'autosomie (2,5 pts) ; génotypes (1 pt) ; hypothèses sur II-5 et calculs de risque (2 pts) ; exemples de maladies (1 pt).

\textbf{Points de vigilance} : la démonstration par l'absurde doit être rédigée en trois temps (hypothèse, conséquence contradictoire avec les données, conclusion) ; ne pas conclure hâtivement que II-5 est $AA$ : les deux génotypes $AA$ et $Aa$ sont compatibles avec des enfants sains ; citer précisément les données de l'arbre à chaque justification.
\end{corrige}

\newpage

\coursec{Fiche bilan}

\courssub{Carte mentale de la leçon}

\begin{popfigure}
\begin{tikzpicture}[
  mindmap,
  grow cyclic,
  every node/.style={concept, execute at begin node=\hskip0pt},
  root concept/.append style={concept color=popblue, minimum size=3.5cm, text width=3cm, font=\bfseries\large, align=center},
  level 1 concept/.append style={sibling angle=72, minimum size=2.5cm, text width=2.2cm, font=\bfseries\normalsize, level distance=5.5cm, align=center},
  level 2 concept/.append style={sibling angle=40, minimum size=1.8cm, text width=1.8cm, font=\small, level distance=4cm, align=center},
  concept color=popblue!70,
  edge from parent/.style={draw=popline, thick}
]
\node[root concept, align=center]{Génétique et\\ Hérédité humaines\\(NA02)}
  child[concept color=popgreen] { node[level 1 concept, align=center]{1. Caryotype\\ humain}
    child { node[level 2 concept] {46 chr (23 paires)} }
    child { node[level 2 concept] {22 paires autos. + 1 paire gon.} }
    child { node[level 2 concept] {Formule : 46,XX / 46,XY} }
  }
  child[concept color=poporange] { node[level 1 concept, align=center]{2. Anomalies\\ du nombre}
    child { node[level 2 concept] {Trisomies (21, 13, 18)} }
    child { node[level 2 concept] {Anomalies gonosomes (XO, XXY, XYY)} }
    child { node[level 2 concept] {Non-disjonction méiotique} }
    child { node[level 2 concept] {Formule chromosomique} }
  }
  child[concept color=poppurple] { node[level 1 concept, align=center]{3. Détermination\\ du sexe}
    child { node[level 2 concept] {Système XX/XY} }
    child { node[level 2 concept] {Gamètes : X ou Y (♂), X (♀)} }
    child { node[level 2 concept] {Rôle gène SRY (chr Y)} }
  }
  child[concept color=poppink] { node[level 1 concept, align=center]{4. Hérédité liée\\ au sexe (X)}
    child { node[level 2 concept] {Daltonisme, Hémophilie (récessif)} }
    child { node[level 2 concept] {♂ hémizygote : phénotype = allèle X} }
    child { node[level 2 concept] {♀ : porteuse ou atteinte} }
    child { node[level 2 concept] {Pas de transmission ♂$\to$♂} }
  }
  child[concept color=popgold] { node[level 1 concept, align=center]{5. Hérédité\\ autosomale}
    child { node[level 2 concept] {Albinisme (récessif)} }
    child { node[level 2 concept] {Drépanocytose (codominance A/S)} }
    child { node[level 2 concept] {Arbre généalogique : ♂/♀ atteints = proportion équivalente} }
  };
\end{tikzpicture}
\legende{Carte mentale récapitulative de la leçon NA02 : 5 piliers structurant la génétique humaine au programme de 1ère A.}
\end{popfigure}

\courssub{Bilan synthétique}

\begin{bilan}

\textbf{1. Définitions clés à connaître par cœur}
\begin{itemize}
  \item \cle{Caryotype} : Photographie des chromosomes d'une cellule au métaphase de la mitose, classés par paires homologues (taille, position du centromère, bandage G/R/Q). Espèce humaine : \textbf{2n = 46 chromosomes} (23 paires).
  \item \cle{Formule chromosomique} : Notation standardisée ISCN : \textbf{46,XX} (femme) ou \textbf{46,XY} (homme). Anomalie numérique : ex. \textbf{47,XX,+21} (trisomie 21), \textbf{45,X} (Turner), \textbf{47,XXY} (Klinefelter), \textbf{47,XYY}, \textbf{47,XXX}.
  \item \cle{Autosomes} : 22 paires de chromosomes non sexuels (identiques chez ♂ et ♀), numérotés 1 à 22 par taille décroissante.
  \item \cle{Gonosomes (chromosomes sexuels)} : Paire 23. ♀ : XX (homogamétique, gamètes tous X) ; ♂ : XY (hétérogamétique, gamètes X ou Y).
  \item \cle{Hérédité liée au sexe (X-récessive)} : Gène porté par le chromosome X (région non homologue au Y). ♂ \textbf{hémizygote} : un seul allèle $\rightarrow$ expression phénotypique directe. Pas de transmission père $\to$ fils.
  \item \cle{Hérédité autosomale} : Gène porté par un autosome (paires 1 à 22). Transmission indépendante du sexe des parents et de la descendance.
  \item \cle{Codominance} : Les deux allèles s'expriment phénotypiquement chez l'hétérozygote (ex. Drépanocytose : $AS$ = drépanocytose mineure, $SS$ = majeure, $AA$ = normal).
  \item \cle{Pénétrance / Expressivité} : Pénétrance = \% d'individus de génotype concerné exprimant le phénotype. Expressivité = intensité variable du phénotype chez les exprimants.
\end{itemize}

\textbf{2. Règles de transmission et formules (Probatoire : savoir les appliquer et les justifier)}
\begin{center}
\begin{tabularx}{\linewidth}{|l|X|X|}\hline
\rowcolor{popblueL}
\textbf{Type} & \textbf{Règles clés / Formules} & \textbf{Phénotypes / Génotypes types} \\ \hline
\textbf{Trisomie 21} & Non-disjonction méiotique (ovocyte I le plus souvent). Formule : \textbf{47,XX,+21} ou \textbf{47,XY,+21}. Risque maternel $\uparrow$ avec l'âge ($>$35 ans). & Syndrome de Down : retard mental variable, traits faciaux caractéristiques, cardiopathies congénitales fréquentes, hypotonie. \\ \hline
\textbf{Autres trisomies} & Trisomie 13 (Patau) : 47,XX,+13 / 47,XY,+13. Trisomie 18 (Edwards) : 47,XX,+18 / 47,XY,+18. Létales précoces le plus souvent. & Malformations multiples graves, survie courte. \\ \hline
\textbf{Anomalies gonosomes} & Turner (45,X) : Monosomie X. Klinefelter (47,XXY) : Disomie X. Triple X (47,XXX). XYY (47,XYY). & Turner : ♀ phénotypique, dysgénie gonadique, stérilité, nanisme. Klinefelter : ♂ phénotypique, hypogonadisme, stérilité, grande taille. \\ \hline
\textbf{Sexuation} & Gène \textbf{SRY} (Yp11.3) $\rightarrow$ facteur testicule-determinant (TDF) $\rightarrow$ différenciation testiculaire $\rightarrow$ sécrétion testostérone + AMH (hormone anti-müllérienne). Sans SRY $\rightarrow$ voie ovarienne (par défaut). & ♂ : 46,XY ; ♀ : 46,XX. Anomalies : 45,X (Turner ♀), 47,XXY (Klinefelter ♂), 46,XX SRY+ (♂ XX), 46,XY SRY- (♀ XY). \\ \hline
\textbf{Daltonisme / Hémophilie (X-réc)} & Allèle muté $X^a$, allèle sain $X^A$. ♂ : $X^a$Y (atteint), $X^A$Y (sain). ♀ : $X^aX^a$ (atteinte), $X^AX^a$ (porteuse saine), $X^AX^A$ (saine). \textbf{Pas de transmission père $\to$ fils}. & ♂ atteints $\gg$ ♀ atteintes. Croisement porteuse ($X^AX^a$) $\times$ sain ($X^A$Y) : 50\% fils atteints ($X^a$Y), 50\% filles porteuses ($X^AX^a$). \\ \hline
\textbf{Albinisme (Auto-réc)} & Gène \emph{TYR} (tyrosinase) chr 11. Allèle $a$ (récessif, non fonctionnel), $A$ (dominant). $aa$ = albinos (absence mélanine). $Aa$ = porteur sain. $AA$ = sain. Consanguinité $\uparrow$ fréquence. & Phénotype : peau/cheveux/yeux dépigmentés, photophobie, nystagmus, acuité visuelle réduite, risque majeur cancers cutanés (UV). \\ \hline
\textbf{Drépanocytose (Auto-codom)} & Gène \emph{HBB} (globine $\beta$) chr 11. Allèles $A$ (HbA normal) et $S$ (HbS muté Glu6Val). $AA$ : sain. $AS$ : \textbf{drépanocytose mineure} (porteur asymptomatique au repos, résistant paludisme). $SS$ : \textbf{drépanocytose majeure} (anémie hémolytique chronique, crises vaso-occlusives, infections, séquestration splénique). & Hétérozygote $AS$ : avantage sélectif zone paludéenne (Cameroun : $\sim$25\% porteurs, 2\% SS). Diagnostic : électrophorèse hémoglobines (HbA $>$ HbS chez AS ; HbS seule chez SS). \\ \hline
\end{tabularx}
\end{center}

\textbf{3. Méthodes express (Pour l'examen)}

\begin{methode}[Identifier un caryotype et son sexe]
\begin{enumerate}
  \item Compter le nombre total de chromosomes (vérifier 2n = 46, 45 ou 47).
  \item Observer la paire 23 (gonosomes) : XX $\rightarrow$ \textbf{Femelle} ; XY $\rightarrow$ \textbf{Mâle}.
  \item Vérifier le nombre d'autosomes (22 paires = 44). Si total $\neq$ 46 $\rightarrow$ anomalie numérique (trisomie ou monosomie).
  \item Identifier le chromosome en excès ou en défaut (bande, taille).
  \item Écrire la formule ISCN : \textit{Nombre total, Gonosomes, [+chr n° ou -chr n°]} (ex: 47,XY,+21).
\end{enumerate}
\end{methode}

\begin{methode}[Interpréter un arbre généalogique (Daltonisme, Hémophilie, Albinisme, Drépanocytose)]
\begin{enumerate}
  \item \textbf{Repérer le mode de transmission} :
    \begin{itemize}
      \item Saute une génération + atteints nés de parents sains (phénotypiquement) $\rightarrow$ \textbf{Récessif}.
      \item Ne saute pas + au moins un parent atteint $\rightarrow$ \textbf{Dominant} (rare au programme).
      \item Atteints quasi-exclusivement ♂ + pas de transmission ♂$\to$♂ + transmission ♂$\to$toutes les filles $\rightarrow$ \textbf{Lié au sexe (X-réc)}.
      \item Atteints ♂/♀ proportion équivalente + transmission ♂$\to$♂ possible $\rightarrow$ \textbf{Autosomal}.
    \end{itemize}
  \item \textbf{Déduire les génotypes} : Partir des phénotypes connus (atteints = génotype certain pour récessif : $aa$, $X^a$Y, $SS$). Remonter vers les parents (obligatoirement hétérozygotes/porteurs), descendre vers la descendance (tableau de Punnett par couple).
  \item \textbf{Calculer probabilités} : Croisement des gamètes (échiquier de Punnett) pour les unions à risque. Exprimer en fractions (1/4, 1/2) ou pourcentages (25\%, 50\%).
\end{enumerate}
\end{methode}

\begin{methode}[Rédiger une conclusion génétique rigoureuse (épreuve écrite)]
\begin{enumerate}
  \item Préciser : Nature de l'allèle (dominant/récessif/codominant), Chromosome porteur (autosome n°... / chromosome X), Statut des parents (homozygotes, hétérozygotes, hémizygotes, porteurs sains).
  \item Utiliser le vocabulaire normé : \emph{« L'allèle responsable est récessif et porté par un autosome (hérédité autosomale récessive) »} ou \emph{« L'allèle est récessif et porté par le chromosome X (hérédité liée au sexe récessive) »}.
  \item Conclure sur le risque de récidive pour la fratrie (frères/sœurs) ou la descendance (enfants) de l'individu consulté.
\end{enumerate}
\end{methode}

\end{bilan}

\courssub{Checklist avant le Probatoire}

\begin{aretenir}[Checklist avant le Probatoire -- Leçon NA02]
  $\square$ Je sais définir \textbf{caryotype}, \textbf{formule chromosomique}, \textbf{autosome}, \textbf{gonosome}, \textbf{hémizygote}, \textbf{codominance}, \textbf{non-disjonction}.
  $\square$ Je sais écrire la formule chromosomique normale (\textbf{46,XX} / \textbf{46,XY}) et celle des anomalies : trisomie 21 (\textbf{47,XX,+21} / \textbf{47,XY,+21}), trisomie 13, trisomie 18, Turner (\textbf{45,X}), Klinefelter (\textbf{47,XXY}), Triple X (\textbf{47,XXX}), XYY (\textbf{47,XYY}).
  $\square$ J'identifie le sexe d'un individu sur un caryotype (paire 23) et je déduis les types de gamètes produits (X ou Y pour ♂ ; X pour ♀).
  $\square$ J'explique le rôle du gène \textbf{SRY} dans la détermination testiculaire (voie mâle) et l'absence de SRY (voie femelle par défaut) ; je cite les hormones effectrices (testostérone, AMH).
  $\square$ Je distingue \textbf{hérédité liée au sexe (X-récessive)} et \textbf{hérédité autosomale (récessive ou codominante)} sur un arbre généalogique (critères : sexe des atteints, transmission père$\to$fils, saut de génération).
  $\square$ Je connais les génotypes/phénotypes types : Daltonisme/Hémophilie ($X^a$Y, $X^AX^a$) ; Albinisme ($aa$, $Aa$, $AA$) ; Drépanocytose ($AA$, $AS$, $SS$).
  $\square$ Je construis un croisement (Punnett) et je calcule les probabilités phénotypiques/génotypiques pour une descendance (ex: porteuse $\times$ sain ; $AS \times AS$ ; $Aa \times Aa$).
  $\square$ J'interprète un arbre généalogique complet : je propose le mode de transmission justifié, je déduis les génotypes des individus clés (numérotés), je calcule le risque pour un fœtus/individu donné.
  $\square$ Je sais argumenter contre les préjugés : \emph{« La drépanocytose n'est pas une "maladie de sorcellerie" mais une maladie génétique à transmission autosomale codominante »} ; \emph{« Le sexe de l'enfant est déterminé par le spermatozoïde (X ou Y), pas par la mère »} ; \emph{« L'albinisme n'est pas une malédiction, c'est une anomalie enzymatique (tyrosinase) à transmission autosomale récessive »}.
  $\square$ Je connais l'avantage sélectif de l'hétérozygote $AS$ (résistance au paludisme \emph{Plasmodium falciparum}) expliquant la fréquence élevée de l'allèle $S$ au Cameroun (équilibre polymorphismique).
  $\square$ Je maîtrise le vocabulaire précis pour la rédaction : \textbf{allèle}, \textbf{locus}, \textbf{homozygote}, \textbf{hétérozygote}, \textbf{hémizygote}, \textbf{phénotype}, \textbf{génotype}, \textbf{non-disjonction}, \textbf{méiose}, \textbf{pénétrance}, \textbf{expressivité}.
\end{aretenir}

\findecours

\end{document}
